\chapter{Refinements of the eight rules}
\label{prov:ch:kern}
\label{kern:ch}

The native scalar core $\Nat$ has eight rules (\cref{base:sec:kernel}).  This
chapter collects the changes of that rule basis which keep the language of
equations, $\Q{}$ and $\D{}{}$ as it is: there are no clause contexts, no
restricted class of public formulas and, with the one exception of the
calculus $\mathsf{U}_{\mathrm{c}}$ in \cref{kern:sec:urforms}, no normal
forms of terms.  The changes fall into five groups.  \Cref{kern:sec:intro} treats the two
introduction rules: progress may introduce a closure directly or through a
witnessed $\D{}{}$, and the two introductions may be merged into one rule
whose mode is read off the matrix and the witness.  \Cref{kern:sec:transport}
treats the transport rule: the number of templates, the crossing of the
conclusion, the exclusion of the literal $0$ as a source, the reversal of the
root equation, replacement with no template in the forward and in the
backward direction, and replacement of a single occurrence.
\Cref{kern:sec:aqe} replaces the elimination \rulename{QE} of a closure by a
rule that computes its matrix from the goal.  \Cref{kern:sec:uniq} replaces induction for equations by
uniqueness of iteration.  \Cref{kern:sec:prog} changes the reading of $\Q{}$
at a literal equation to progressivity and presents the seven-rule calculi
that result.  \Cref{kern:sec:recursors} records two results on recursors other
than the iterator, which two of the packages prove on the side.
\Cref{kern:sec:depth} records two limits that concern the native language and
its presentation rather than a rule: the classes of relations defined by
native formulas of bounded binder depth, and an obstruction to finite
ordinary axiomatizations of $\iPRA$.

Every calculus of
\cref{kern:sec:intro,kern:sec:transport,kern:sec:aqe,kern:sec:uniq}
derives exactly the sequents of $\Nat$, on the same formulas.  The calculi of
\cref{kern:sec:prog} do not; they are equivalent to $\Nat$ through a pair of
translations with both formula round trips.  \Cref{kern:tab:calculi} at the
end of \cref{kern:sec:prog} lists all calculi of the chapter.

% ---------------------------------------------------------------------
\section{Conventions and the comparison method}
\label{kern:sec:conv}

This section fixes what all packages of the chapter restate before
their own results; it draws on no package in particular.

The rules of $\Nat$ are named as in \cref{base:sec:kernel}: \rulename{Hyp},
\rulename{I0}, \rulename{IS}, \rulename{Tr}, \rulename{QE}, \rulename{DE},
\rulename{DI} (introduction of $\D{}{}$ by cases) and \rulename{QI}
(introduction of $\Q{}$ by progress at a literal equation).  A \emph{term
template} $h(v)$ is a term with a designated variable $v$, which may occur
several times, once, or not at all; in the last case the template is
\emph{constant}.  An \emph{equation template} is $a(v)=b(v)$ for two term
templates, and a \emph{formula template} $F(v)$ is a native formula with a
designated variable.  The sign $\equiv$ is identity of syntax (up to the
names of bound variables), so that $s\not\equiv 0$ says that $s$ is not the
term $0$.  This is a test on syntax and not on values:
$\It{0}{0}{u}{0}\not\equiv 0$, although the term equals $0$ by \rulename{I0}.
A formula is \emph{compound} if it is not a literal equation, that is, if its
root is $\Q{}$ or $\D{}{}$.  Substitution preserves the root of a formula and
may change whether a term is the literal $0$ or a literal successor.

All calculi of the chapter have the terms of \cref{base:sec:terms}, the
formulas of \cref{base:sec:reading}, finite lists of formulas as contexts, and
the rules \rulename{Hyp}, \rulename{I0}, \rulename{IS} and \rulename{DE} in
their forms of \cref{base:sec:kernel}.  A rule inventory below therefore
names these four rules and displays the others.  The one exception is the
calculus $\mathsf{U}_{\mathrm{c}}$ of \cref{kern:def:Uc}, which is stated on
normal terms and has reflexivity in place of the two computation equations.
Derivability in a calculus
$\mathsf{X}$ is written $\proves_{\mathsf{X}}$, and $\proves$ without a
subscript is derivability in $\Nat$.

\begin{definition}[Equivalence on the same formulas]
\label{kern:def:equiv}
Two calculi $\mathsf{X}$ and $\mathsf{Y}$ on the native language are
\emph{equivalent} if $\Gamma\proves_{\mathsf{X}}A$ holds exactly when
$\Gamma\proves_{\mathsf{Y}}A$ holds, for every context $\Gamma$ and every
native formula $A$.
\end{definition}

The equivalence is about sequents with arbitrary contexts and arbitrary
formulas, not about closed equations.  It is always proved by two translations
of derivations that copy the common rules and replace each instance of a
changed rule by a fixed finite derivation; the translations are effective and
leave terms, formulas, contexts and conclusions unchanged.

\begin{proposition}[Transfer of the bridge]
\label{kern:prop:transfer}
Let $\mathsf{X}$ be equivalent to $\Nat$.  Then every statement of
\cref{base:thm:bridge} holds for $\mathsf{X}$ with the same translations:
$\mathsf{X}$, $\Teq$ and $\iPRA$ are bi-interpretable with whole domains,
actual equality and identity numerical composites, both formula round trips
hold, and $\proves_{\mathsf{X}}0=\Suc 0$ holds exactly when $\iPRA$ is
inconsistent.  Weakening, substitution of terms, cut and the Leibniz rule are
admissible in $\mathsf{X}$.
\end{proposition}

\begin{proof}
The statements of \cref{base:thm:bridge,base:sec:structural} about $\Nat$ are
statements about its derivable sequents, and $\mathsf{X}$ has the same
derivable sequents.
\end{proof}

Each package of the chapter restates this transfer for its own calculus; it
is not repeated below.  The last sentence of the proposition concerns the
consequence relation of $\mathsf{X}$.  Whether an operation maps each rule
instance to an instance of the same rule is a different question.

\begin{lemma}[Structural operations by direct induction]
\label{kern:lem:struct}
In every calculus of this chapter, weakening of the context, injective
renaming of free variables, and cut are admissible by a direct induction on
derivations.  Substitution of terms for variables maps rule instances to rule
instances of the same rule in a calculus none of whose side conditions tests
a term for being the literal $0$, a literal successor or a syntactic match of
another term.
\end{lemma}

\begin{proof}
The side conditions of the rules below are of four kinds: freshness of an
eigenvariable; a test on the root of the matrix; a test whether a term of the
instance is the literal $0$ or a literal successor; and, for the two rules of
replacement with no template and for the rule \rulename{AQE} of
\cref{kern:sec:aqe}, a syntactic match of subterms.  An injective renaming preserves all four.  For
weakening, rename the eigenvariables of the derivation away from the added
formulas and keep the added formulas in every context; each inference is
still an instance of its rule with the same terms.  For cut, let
$\Gamma\proves_{\mathsf{X}}A$ and $\Gamma,A\proves_{\mathsf{X}}C$.  Rename
the eigenvariables of the second derivation away from the first, and replace
every use of \rulename{Hyp} for the assumption $A$ by the first derivation,
weakened by the local assumptions.  No term of an inference changes, so every
side condition is preserved.  For substitution the first two kinds of side
condition are preserved after a renaming of eigenvariables, as in the proof
for $\Nat$ in \cref{base:sec:structural}; the last two kinds need not be.
\end{proof}

Several derived rules of \cref{base:sec:structural} use only part of $\Nat$,
and they are used in that form.

\begin{lemma}[Derived rules and the rules they use]
\label{kern:lem:local}
Let $\mathsf{X}$ be a calculus of this chapter.
\begin{enumerate}[label=\textup{(\roman*)}]
\item If $\mathsf{X}$ contains \rulename{QE} and \rulename{DI}, then
 $\mathsf{X}$ derives: $\D{x}{y}B\dashv\vdash\Q{y}B$ for
 $x\notin\FV(B)$; the introduction of $\Q{y}B$ by cases, from
 $\Gamma,B(0)\proves B(\Suc z)$ with $z$ fresh, for an arbitrary matrix $B$;
 numerical cases, from $\Gamma\proves B(0)$ and $\Gamma\proves B(\Suc z)$
 with $z$ fresh to $\Gamma\proves B(t)$; and $\D{x}{y}B(x,y)$ from
 $\Q{y}B(t,y)$.
\item If $\mathsf{X}$ contains \rulename{QE} and \rulename{QI}, then
 $\mathsf{X}$ derives equation induction: from $\Gamma\proves E(0)$ and
 $\Gamma,E(z)\proves E(\Suc z)$, with $E$ a literal equation and $z$ fresh,
 to $\Gamma\proves E(t)$.
\item If $\mathsf{X}$ contains \rulename{Tr}, then $\mathsf{X}$ derives
 reflexivity, symmetry, transitivity, congruence in every term template, and
 transport in an equation template, from $s=t$ and $a(s)=b(s)$ to
 $a(t)=b(t)$.
\end{enumerate}
\end{lemma}

\begin{proof}
The derivations of these rules in \cref{base:sec:structural} use
\rulename{Hyp}, \rulename{DE}, weakening and the rules named in each item;
for (iii) they use \rulename{I0} and \rulename{Tr} alone.
\end{proof}

In (iii) each of the four equality rules is one application of
\rulename{Tr}, and transport in an equation template is two.  Symmetry needs
a constant template: \rulename{Tr} with the template $h(v)\defeq s$, any
derivable equation as its first premise and $s=t$ as its second premise
concludes $t=s$.

% ---------------------------------------------------------------------
\section{The introduction rules}
\label{kern:sec:intro}

This section draws on \pkg{direct\_closure}, \pkg{direct\_closure\_qprogress}
and \pkg{witness\_rule}, and on the obstruction section of
\pkg{reversal\_core}.

\subsection{Progress through a witness and progress at the closure}
\label{kern:sec:direct}

The reports state induction for equations as an introduction of $\D{}{}$
with a witness:
\[
 \infrule{\Gamma,\ E(t,z)\proves E(t,\Suc z)}
         {\Gamma\proves\D{x}{y}E(x,y)}\ \rulename{DP}
 \qquad(E\text{ a literal equation}),
\]
where $z$ is not free in $\Gamma$, in $t$ or in $\D{x}{y}E$.

\begin{definition}
\label{kern:def:ND}
$\Nat_{\mathrm{D}}$ is the calculus with the rules \rulename{Hyp},
\rulename{I0}, \rulename{IS}, \rulename{Tr}, \rulename{QE}, \rulename{DE},
\rulename{DI} and \rulename{DP}.  It is $\Nat$ with \rulename{QI} replaced by
\rulename{DP}.
\end{definition}

The rule \rulename{DP} has a matrix in two variables and a witness term; the
rule \rulename{QI} has a matrix in one variable and no witness.

\begin{theorem}[Direct closure]
\label{kern:thm:direct}
$\Nat_{\mathrm{D}}$ and $\Nat$ are equivalent.  The proof uses only the
rules \rulename{Hyp}, \rulename{QE}, \rulename{DE}, \rulename{DI} and
weakening; it holds for any common set of further rules for which weakening
and renaming are admissible.
\end{theorem}

\begin{proof}
All rules but one are common.  \emph{\rulename{QI} in $\Nat_{\mathrm{D}}$.}
Let $\Gamma,E(z)\proves E(\Suc z)$ be the translated premise.  Choose a
variable $x$ not free in $E$ and $\Gamma$.  The rule \rulename{DP} with the
matrix $E(y)$, in which the coordinate $x$ is not used, and the witness $0$
has exactly this premise and concludes $\Gamma\proves\D{x}{y}E(y)$.  Apply
\rulename{DE} with a fresh $b$: the opened hypothesis is $\Q{y}E(y)$, because
$x$ does not occur in $E$, and the second premise
$\Gamma,\Q{y}E(y)\proves\Q{y}E(y)$ is \rulename{Hyp}.  The conclusion is
$\Gamma\proves\Q{y}E(y)$.

\emph{\rulename{DP} in $\Nat$.}  Let $\Gamma,E(t,z)\proves E(t,\Suc z)$ be the
translated premise.  The formula $E(t,y)$ is a literal equation in $y$ with
$t$ as a parameter, and $z$ is not free in $t$, so \rulename{QI} gives
$\Gamma\proves\Q{y}E(t,y)$.  Let $w$ be fresh.  Over $\Gamma,E(t,0)$,
weakening and \rulename{QE} at $\Suc w$ give $E(t,\Suc w)$.  This is the
premise of \rulename{DI} with the witness $t$, which concludes
$\Gamma\proves\D{x}{y}E(x,y)$.  No equality rule, no computation equation and
no substitution into the given premise is used.
\end{proof}

\begin{proposition}[Node counts]
\label{kern:prop:nodes}
Count every instance of a rule as one node.  A derivation of $\Nat$ with $n$
nodes, $r$ of them instances of \rulename{QI}, translates to a derivation of
$\Nat_{\mathrm{D}}$ with $n+2r\le 3n$ nodes.  A derivation of
$\Nat_{\mathrm{D}}$ with $n$ nodes, $r$ of them instances of \rulename{DP},
translates to a derivation of $\Nat$ with $n+3r\le 4n$ nodes.
\end{proposition}

\begin{proof}
An instance of \rulename{QI} becomes \rulename{DP}, \rulename{DE} and
\rulename{Hyp} over the same premise.  An instance of \rulename{DP} becomes
\rulename{QI}, \rulename{QE}, \rulename{Hyp} and \rulename{DI} over the same
premise.  Each premise derivation is used once, and weakening and renaming do
not change the number of nodes.
\end{proof}

The bounds count nodes of the inference tree; they do not bound the size of
formulas, of serialized certificates or of a checker.

\begin{remark}[Sources]
\label{kern:rem:direct-sources}
\pkg{direct\_closure} and \pkg{direct\_closure\_qprogress} are two
independent accounts of \cref{kern:thm:direct,kern:prop:nodes}, with the same
two constructions and the same counts; the second contains the proof twice,
once as a proof note and once as an audit.  Both state the calculi with
transport in a formula template in place of \rulename{Tr}, which the theorem
permits.  Both stress that the rule \rulename{QI} was a derived rule of the
calculus before it was taken as primitive, so that the change is one of the
choice of primitives.  Their example is $F(n)\defeq\It{n}{0}{u}{\Suc u}$ with
the equation $F(y)=y$: with transport in a formula template the derivation of
$F(n)=n$ by \rulename{QI} and \rulename{QE} has six nodes, its translation to
$\Nat_{\mathrm{D}}$ has eight, and the translation of that back has eleven.
In Lean the rule \rulename{QI} is the constructor \Lean{qIInd} of the kernel,
adopted from these packages, and \rulename{DP} is the derived rule
\Lean{NPrf.dIProg}; the equivalence of the calculus of the reports, with
\rulename{DP} and transport in a formula template, is \Lean{iprf\_iff\_nprf}.
\end{remark}

\subsection{One closure introduction with a mode read off the matrix}
\label{kern:sec:det}

For a matrix $B(y)$ and a variable $z$ put $\eta_B(z)\defeq z$ if $B$ is a
literal equation and $\eta_B(z)\defeq 0$ if $B$ is compound.

\begin{definition}
\label{kern:def:Ndet}
$\Nat_{\det}$ is $\Nat$ with \rulename{QI} replaced by
\[
 \infrule{\Gamma,\ B(\eta_B(z))\proves B(\Suc z)}
         {\Gamma\proves\Q{y}B(y)}\ \rulename{QI}_{\eta},
\]
where $B$ is an arbitrary native formula and $z$ is not free in $\Gamma$ and
$\Q{y}B$.  Its rules are \rulename{Hyp}, \rulename{I0}, \rulename{IS},
\rulename{Tr}, \rulename{QE}, \rulename{DE}, \rulename{DI} and
$\rulename{QI}_{\eta}$.
\end{definition}

A certificate for $\rulename{QI}_{\eta}$ carries no choice between cases and
progress: the root of the matrix decides.

\begin{corollary}
\label{kern:cor:det}
$\Nat_{\det}$, $\Nat$ and $\Nat_{\mathrm{D}}$ are equivalent, and the bounds
of \cref{kern:prop:nodes} hold between $\Nat_{\det}$ and $\Nat_{\mathrm{D}}$.
\end{corollary}

\begin{proof}
At a literal equation $\rulename{QI}_{\eta}$ is \rulename{QI}.  At a compound
matrix it is the introduction of $\Q{}$ by cases, which is derivable in $\Nat$
by \cref{kern:lem:local}(i): \rulename{DI} with an unused coordinate and the
witness $0$, then \rulename{DE} and \rulename{Hyp}.  This is the same
three-node construction as in \cref{kern:thm:direct}, with \rulename{DI} in
place of \rulename{DP}.  Conversely $\Nat_{\det}$ contains \rulename{QI}.
\end{proof}

The introduction of $\Q{y}E$ by cases for a literal equation $E$ is not an
instance of $\rulename{QI}_{\eta}$, and it remains derivable through
\rulename{DI} and \rulename{DE}.  For example $\Q{y}(y=1)$ is derivable in
all three calculi from the premise $0=1\proves\Suc z=1$ of the introduction
by cases, while the premise $z=1\proves\Suc z=1$ of progress is not
derivable.

\subsection{What cannot be removed from the introductions}
\label{kern:sec:introlimits}

\begin{proposition}[The guard of \rulename{QI}]
\label{kern:prop:guard}
If \rulename{QI} is allowed at every matrix, the resulting calculus derives
induction for every native formula: from $\Gamma\proves B(0)$ and
$\Gamma,B(z)\proves B(\Suc z)$ with $z$ fresh to $\Gamma\proves B(t)$.
\end{proposition}

\begin{proof}
The unrestricted rule gives $\Gamma\proves\Q{y}B(y)$, and \rulename{QE} gives
$B(t)$.
\end{proof}

Through the native connectives of \cref{base:sec:structural} this is
induction for every first-order formula, and the bridge of
\cref{base:thm:bridge}, whose target has induction for quantifier-free
formulas only, no longer applies.  The rule $\rulename{QI}_{\eta}$ avoids
this by using the instance at zero at a compound matrix.

\begin{proposition}[\rulename{DI} cannot be dropped]
\label{kern:prop:nodi}
Let $\Nat^{-}$ be $\Nat$ without \rulename{DI}, or $\Nat_{\det}$ without
\rulename{DI}.  If no formula of $\Gamma$ has the root $\D{}{}$, then
$\Nat^{-}$ derives from $\Gamma$ no formula with the root $\D{}{}$.  In
particular $\Nat^{-}$ does not derive $\D{x}{y}(0=0)$, which $\Nat$ derives.
\end{proposition}

\begin{proof}
Suppose a derivation of $\Gamma\proves_{\Nat^{-}}A$ with $A$ of root
$\D{}{}$ and $\Gamma$ as stated, and take one with the least number of nodes.
Its last rule is not \rulename{Hyp}, by the assumption on $\Gamma$; not a
computation equation or \rulename{Tr}, whose conclusions are equations; and
not an introduction of $\Q{}$, whose conclusion has the root $\Q{}$.  If it
is \rulename{QE} with the conclusion $B(t)$, then its second premise
$\Gamma\proves B(0)$ has the same root $\D{}{}$ and the same context, and a
smaller derivation.  If it is \rulename{DE}, then its first premise
$\Gamma\proves\D{x}{y}B$ has the root $\D{}{}$ and the same context, and a
smaller derivation.  Both contradict minimality.  In $\Nat$ the formula
$\D{x}{y}(0=0)$ follows by \rulename{DI} with the witness $0$ from the
premise $0=0\proves 0=0$.
\end{proof}

The assumptions inside a subderivation of an introduction of $\Q{}$ do not
affect the argument, because such an introduction is never the last rule of
a derivation of a formula with the root $\D{}{}$.

\begin{proposition}[\rulename{DI} cannot be restricted to compound matrices]
\label{kern:prop:compoundonly}
Restrict \rulename{DI} to compound matrices $B$.
\begin{enumerate}[label=\textup{(\roman*)}]
\item The calculus with \rulename{Hyp}, \rulename{I0}, \rulename{IS},
 \rulename{Tr}, \rulename{QE}, \rulename{DE}, \rulename{QI} and the
 restricted \rulename{DI} does not derive $\D{x}{y}(0=0)$.
\item The calculus with \rulename{Hyp}, \rulename{I0}, \rulename{IS},
 \rulename{Tr}, \rulename{QE}, \rulename{DE}, \rulename{DP} and the
 restricted \rulename{DI} does not derive $\D{x}{y}(y=1)$.
\end{enumerate}
Both formulas are derivable in $\Nat$.  The statements remain true with
transport in a formula template in place of \rulename{Tr}.
\end{proposition}

\begin{proof}
Both parts interpret terms in the standard natural numbers, equations as
equality and $\Q{y}B$ by its ordinary reading, and give a formula
$\D{x}{y}B$ with a compound matrix its ordinary reading.  The interpretation
of a formula $\D{x}{y}E$ with a literal equation $E$ is changed.  It commutes
with substitution, because substitution preserves the root of a matrix, and
truth of a sequent means that every assignment satisfying the context
satisfies the conclusion.

(i) Let $\D{x}{y}E$ be false for every literal equation $E$.  The rules
\rulename{Hyp}, \rulename{I0}, \rulename{IS}, transport and \rulename{QE}
preserve truth as they do for the ordinary reading.  \rulename{QI} preserves
truth by induction on the natural numbers.  \rulename{DE} at an atomic matrix
has a false first premise, and at a compound matrix it is the elimination of
an existential quantifier.  The restricted \rulename{DI} preserves truth
because every natural number is zero or a successor.  The formula
$\D{x}{y}(0=0)$ is false.

(ii) Let $\D{x}{y}E$ mean
$\exists x\,\forall y\,\bigl(E(x,y)\to E(x,\Suc y)\bigr)$ for a literal
equation $E$.  \rulename{DP} preserves truth: its premise states the
progress of $E(t,y)$, and $t$ is a witness.  \rulename{DE} at an atomic
matrix preserves truth: if $\forall y\,(E(b,y)\to E(b,\Suc y))$ holds, then
$E(b,0)\to\forall y\,E(b,y)$ holds by induction on the natural numbers, and
this is the reading of the opened hypothesis $\Q{y}E(b,y)$.  The other rules
preserve truth as in (i).  The formula $\D{x}{y}(y=1)$ is false, because at
$y=1$ the equation $1=1$ holds and $2=1$ does not.

In $\Nat$, $\D{x}{y}(0=0)$ is derived in \cref{kern:prop:nodi}, and
$\D{x}{y}(y=1)$ follows by \rulename{DI} with the witness $0$ from
$0=1\proves\Suc z=1$, which holds because $0=1$ derives every equation
(\cref{base:sec:structural}).
\end{proof}

\begin{remark}[Sources]
\label{kern:rem:limits-sources}
\Cref{kern:prop:guard,kern:prop:nodi} are from
\pkg{direct\_closure\_qprogress}; \pkg{direct\_closure} states the first.
\Cref{kern:prop:compoundonly}(i) is the obstruction section of
\pkg{reversal\_core}.  Part (ii) is stated in
\pkg{direct\_closure\_qprogress} with the interpretation and a reference to
an earlier package for the verification of the rules.
% ed.: the verification of the rules under the progressive interpretation is
% written out here; the package gives the interpretation and the falsified
% instance only.
\pkg{direct\_closure\_qprogress} also recalls why the ordinary theory $\Teq$
needs its coverage axiom: the product of two copies of the standard model,
with all operations coordinatewise, satisfies the computation equations,
separation and induction for equations, because the solution set of an
equation is a product $U_0\times U_1$, the base puts $0$ into both factors,
and the step applied at $(r,0)$ and at $(0,r)$ closes each factor under
successor; and the element $(0,1)$ is neither zero nor a successor.
\end{remark}

\subsection{One introduction with a mode read off the matrix and the witness}
\label{kern:sec:witness}

The two introductions \rulename{DI} and \rulename{QI} can be replaced by one
introduction of $\D{}{}$.  For a matrix $B(x,y)$, a witness $t$ and a
variable $z$ put
\[
 \eta_{B,t}(z)\defeq
 \begin{cases}
  z, & B\text{ is a literal equation and }t\equiv 0,\\
  0, & \text{otherwise.}
 \end{cases}
\]

\begin{definition}[The calculus $\mathsf{W}$]
\label{kern:def:W}
$\mathsf{W}$ has the seven rules \rulename{Hyp}, \rulename{I0},
\rulename{IS}, \rulename{Tr}, \rulename{QE}, \rulename{DE} and
\[
 \infrule{\Gamma,\ B(t,\eta_{B,t}(z))\proves B(t,\Suc z)}
         {\Gamma\proves\D{x}{y}B(x,y)}\ \rulename{DI}_{\mathrm{w}},
\]
where $z$ is not free in $\Gamma$, in $t$ or in $\D{x}{y}B$.
\end{definition}

The rule selects progress exactly when the matrix is a literal equation and
the witness is the literal $0$, and cases otherwise.  A witness such as
$\It{0}{0}{u}{0}$, which equals $0$, selects cases.  The calculus has no rule
that introduces $\Q{}$.

\begin{lemma}[A second name for zero]
\label{kern:lem:alias}
Let $c\defeq\It{0}{0}{u}{0}$.  Every calculus containing \rulename{I0} and
\rulename{Tr} derives $c=0$ and $0=c$, and $c\not\equiv 0$.
\end{lemma}

\begin{proof}
$c=0$ is an instance of \rulename{I0}, and $0=c$ follows by symmetry
(\cref{kern:lem:local}(iii)).
\end{proof}

\begin{theorem}[Witness-selected introduction]
\label{kern:thm:witness}
$\mathsf{W}$ and $\Nat$ are equivalent.
\end{theorem}

\begin{proof}
\emph{From $\mathsf{W}$ to $\Nat$.}  An instance of
$\rulename{DI}_{\mathrm{w}}$ in the mode of cases is an instance of
\rulename{DI}.  An instance in the mode of progress has a literal equation
$E(x,y)$ as its matrix, the witness $0$ and the premise
$\Gamma,E(0,z)\proves E(0,\Suc z)$; it is an instance of \rulename{DP} with
the witness $0$, which $\Nat$ derives by \cref{kern:thm:direct}.

\emph{From $\Nat$ to $\mathsf{W}$.}  The rule \rulename{QI} with the premise
$\Gamma,E(z)\proves E(\Suc z)$: choose $x$ not free in $E$ and $\Gamma$.  The
rule $\rulename{DI}_{\mathrm{w}}$ with the matrix $E(y)$ and the witness $0$
is in the mode of progress, has this premise and concludes
$\D{x}{y}E(y)$; \rulename{DE} and \rulename{Hyp} give $\Q{y}E(y)$ as in
\cref{kern:thm:direct}.

The rule \rulename{DI} with a compound matrix, or with a witness
$t\not\equiv 0$, is an instance of $\rulename{DI}_{\mathrm{w}}$.  There
remains \rulename{DI} with a literal equation $E(x,y)$ as its matrix and the
witness $0$.  Let
\[
 \pi\colon\quad\Gamma,\ E(0,0)\proves_{\mathsf{W}}E(0,\Suc z)
\]
be the translated premise, with $z$ fresh.  Work over $\Gamma,E(c,0)$ with
$c$ as in \cref{kern:lem:alias}.  Transport in the equation template $E(v,0)$
along $c=0$ gives $E(0,0)$ (\cref{kern:lem:local}(iii)).  Cut with $\pi$,
which is admissible in $\mathsf{W}$ by \cref{kern:lem:struct}, gives
$E(0,\Suc z)$.  Transport in the equation template $E(v,\Suc z)$ along $0=c$
gives $E(c,\Suc z)$.  Hence
$\Gamma,E(c,0)\proves_{\mathsf{W}}E(c,\Suc z)$.  Since $c\not\equiv 0$, the
rule $\rulename{DI}_{\mathrm{w}}$ with the witness $c$ is in the mode of
cases and has exactly this premise.  It concludes $\D{x}{y}E(x,y)$; the
conclusion does not record the witness.
\end{proof}

The construction changes an assumption by a derived transport and a cut.  It
does not substitute a term into the derivation $\pi$, so it cannot change the
mode of an introduction inside $\pi$.

\begin{remark}[Substitution in $\mathsf{W}$]
\label{kern:rem:wsubst}
A substitution can turn a witness that is a variable into the literal $0$,
and then an instance of $\rulename{DI}_{\mathrm{w}}$ in the mode of cases is
mapped to a sequent that is not the premise of an instance in the mode of
progress.  Substitution is admissible in $\mathsf{W}$ by
\cref{kern:prop:transfer}, through $\Nat$, and is not obtained by a direct
induction on the derivations of $\mathsf{W}$.  For this reason the Lean
kernel keeps the two introductions of $\Nat$.
\end{remark}

\begin{remark}[Sources and variants]
\label{kern:rem:witness-sources}
\Cref{kern:thm:witness} is the theorem of \pkg{witness\_rule}; in Lean it is
\Lean{wprf\_iff\_nprf}, for the calculus \Lean{WPrf}.  The package compares
$\mathsf{W}$ with the calculus that has \rulename{DI} and \rulename{DP}, so
its proof has four cases: \rulename{DI} repeated; \rulename{DI} at an
equation with the witness $0$, as above; \rulename{DP} with the witness $0$
repeated; and \rulename{DP} with a witness $t\not\equiv 0$, which is the
second construction of \cref{kern:thm:direct} with the inner closure
introduced through an unused coordinate.  The package works with a
restricted inventory of computation instances and therefore uses
$\It{2}{1}{u}{0}$ as its second name for zero; that inventory is treated in
\cref{kern:lem:comprepair}.

The argument does not use any property of the number zero in the witness
coordinate.  If $p$ is a closed term and $q\not\equiv p$ is a closed term
with $p=q$ and $q=p$ derivable from the common rules, then progress may be
selected exactly at the witness $p$: an introduction by cases at $p$ is
repaired at the witness $q$, and every auxiliary introduction with an unused
coordinate takes the witness $p$.  The coordinate $y$ of the induction still
starts at $0$.  The literal $0$ is chosen because it is the cheapest witness
to recognize.

\pkg{one\_context} states the same calculus with the two-template transport
of \cref{kern:sec:templates} replaced by \rulename{Tr} and calls it
$\mathcal{O}_7$; since $\Nat$ already has \rulename{Tr}, its calculus is
$\mathsf{W}$ as defined here.
\end{remark}

% ---------------------------------------------------------------------
\section{The transport rule}
\label{kern:sec:transport}

This section draws on \pkg{one\_context}, \pkg{one\_sided\_transport},
\pkg{reversal\_core}, \pkg{source\_dispatch}, \pkg{template\_free},
\pkg{backward\_replacement} and the note on single occurrences in
\pkg{kernel\_refinement}.

Throughout the section $J(a)\defeq\It{0}{a}{u}{0}$ and $c\defeq J(0)$.  The
equation $J(a)=a$ is an instance of \rulename{I0}, and $J(a)\not\equiv 0$ for
every term $a$; in particular $c=0$ is an instance of \rulename{I0} and
$c\not\equiv 0$.

\subsection{One template, two templates, a formula template}
\label{kern:sec:templates}

The rule \rulename{Tr} of $\Nat$ has one term template and a crossed
conclusion: from $s=t$ and $h(s)=r$ it infers $r=h(t)$.  The reports use two
stronger forms,
\[
 \infrule{\Gamma\proves s=t\qquad\Gamma\proves a(s)=b(s)}
         {\Gamma\proves a(t)=b(t)}\ \rulename{Tr}_{2}
 \qquad\qquad
 \infrule{\Gamma\proves s=t\qquad\Gamma\proves F(s)}
         {\Gamma\proves F(t)}\ \rulename{Tr}_{F},
\]
with two term templates $a(v)$, $b(v)$ and with a native formula template
$F(v)$.

\begin{definition}
\label{kern:def:N2}
$\Nat_{2}$ is $\Nat$ with \rulename{Tr} replaced by $\rulename{Tr}_{2}$, and
$\Nat_{F}$ is $\Nat$ with \rulename{Tr} replaced by $\rulename{Tr}_{F}$.  Each
has the rules \rulename{Hyp}, \rulename{I0}, \rulename{IS}, its transport
rule, \rulename{QE}, \rulename{DE}, \rulename{DI} and \rulename{QI}.
\end{definition}

\begin{theorem}[One term template suffices]
\label{kern:thm:onecontext}
$\Nat$, $\Nat_{2}$ and $\Nat_{F}$ are equivalent.  An instance of
$\rulename{Tr}_{2}$ is simulated in $\Nat$ by two instances of \rulename{Tr}
with the same first premise.  An instance of \rulename{Tr} is simulated in
$\Nat_{2}$ by one instance of $\rulename{Tr}_{2}$ with a constant second
template, followed by symmetry.
\end{theorem}

\begin{proof}
\emph{$\rulename{Tr}_{2}$ in $\Nat$.}  \rulename{Tr} with the template $a$,
applied to $s=t$ and $a(s)=b(s)$, gives $b(s)=a(t)$.  \rulename{Tr} with the
template $b$, applied to $s=t$ and $b(s)=a(t)$, gives $a(t)=b(t)$; the term
$a(t)$ is the fixed side of the second step.

\emph{\rulename{Tr} in $\Nat_{2}$.}  First, $\Nat_{2}$ derives reflexivity and
symmetry: $\rulename{Tr}_{2}$ with the templates $v$ and $a$, applied to
$J(a)=a$ and $J(a)=a$, gives $a=a$; and $\rulename{Tr}_{2}$ with the
templates $v$ and $s$, applied to $s=t$ and $s=s$, gives $t=s$.  Now let
$s=t$ and $h(s)=r$ be given.  $\rulename{Tr}_{2}$ with the templates $h(v)$
and the constant $r$ gives $h(t)=r$, and symmetry gives $r=h(t)$.

\emph{$\Nat_{F}$.}  An instance of $\rulename{Tr}_{2}$ is the instance of
$\rulename{Tr}_{F}$ at the template $a(v)=b(v)$, so $\Nat_{F}$ derives every
sequent of $\Nat_{2}$.  Conversely $\rulename{Tr}_{F}$ is a derived rule of
$\Nat$, by the Leibniz rule of \cref{base:sec:structural}.

In each direction the seven other rules are copied, so the induction on
derivations applies to open sequents and to premises under both binders.
\end{proof}

The crossing of the conclusion is what makes two successive replacements in
one template reproduce a simultaneous replacement in both sides of an
equation.  The rule \rulename{Tr} has no side condition, so substitution maps
its instances to instances (\cref{kern:lem:struct}).

\begin{remark}[Size of the simulation]
\label{kern:rem:dag}
The simulation of $\rulename{Tr}_{2}$ uses the derivation of $s=t$ twice.  A
proof represented as a directed acyclic graph shares it; a proof represented
as a tree copies it, and repeated translation can enlarge a tree
substantially.  No linear bound on the size of translated trees is claimed
in \pkg{one\_context} or \pkg{one\_sided\_transport}.
\end{remark}

\begin{remark}[Sources]
\label{kern:rem:onecontext-sources}
The comparison of $\Nat$ with $\Nat_{2}$ is the theorem of
\pkg{one\_context}, whose reference calculus is $\Nat_{2}$.  In Lean the rule
\rulename{Tr} is the constructor \Lean{NPrf.trOne} of the kernel, adopted
from that package, and $\rulename{Tr}_{2}$ is the derived rule
\Lean{NPrf.trEq}.  The comparison with $\Nat_{F}$ is the remark on the
calculus of the reports in \cref{base:sec:kernel}.
\end{remark}

\subsection{Sources other than the literal zero}
\label{kern:sec:onesided}

Let $\rulename{Tr}^{-}$ be the rule \rulename{Tr} with the side condition
$s\not\equiv 0$ on the left side of its first premise, the \emph{source}:
\[
 \infrule{\Gamma\proves s=t\qquad\Gamma\proves h(s)=r}
         {\Gamma\proves r=h(t)}\ \rulename{Tr}^{-}
 \qquad(s\not\equiv 0).
\]
Every variable and every term with the root $\Suc$ or $I$ is a permitted
source, including terms whose value is zero.

\begin{definition}[The calculi $\mathsf{O}$ and $\mathsf{M}$]
\label{kern:def:OM}
$\mathsf{O}$ is $\Nat$ with \rulename{Tr} replaced by $\rulename{Tr}^{-}$;
its rules are \rulename{Hyp}, \rulename{I0}, \rulename{IS},
$\rulename{Tr}^{-}$, \rulename{QE}, \rulename{DE}, \rulename{DI} and
\rulename{QI}.  $\mathsf{M}$ is $\mathsf{W}$ with \rulename{Tr} replaced by
$\rulename{Tr}^{-}$; its seven rules are \rulename{Hyp}, \rulename{I0},
\rulename{IS}, $\rulename{Tr}^{-}$, \rulename{QE}, \rulename{DE} and
$\rulename{DI}_{\mathrm{w}}$.
\end{definition}

\begin{lemma}[Equality from \rulename{I0} and $\rulename{Tr}^{-}$]
\label{kern:lem:oseq}
The following are derivable from \rulename{I0} and $\rulename{Tr}^{-}$ alone,
over an arbitrary context.
\begin{enumerate}[label=\textup{(\roman*)}]
\item Reflexivity: $a=a$.
\item Symmetry: from $a=b$ to $b=a$, also when $a$ or $b$ is the literal
 $0$.
\item Restricted transport in an equation template: if $s\not\equiv 0$,
 from $s=t$ and $a(s)=b(s)$ to $a(t)=b(t)$.
\item Joining through zero: if $p\not\equiv 0$ and $q\not\equiv 0$, from
 $p=0$ and $q=0$ to $p=q$.
\item Transitivity: from $r=s$ and $s=t$ to $r=t$.
\item Congruence: from $s=t$ to $f(s)=f(t)$, for every term template $f$.
\item The rule \rulename{Tr} without the side condition.
\end{enumerate}
\end{lemma}

\begin{proof}
(i) $\rulename{Tr}^{-}$ with the template $v$, applied to $J(a)=a$ and
$J(a)=a$, gives $a=a$; the source is $J(a)\not\equiv 0$.

(ii) $\rulename{Tr}^{-}$ with the constant template $a$, applied to $c=0$
and $a=b$, gives $b=a$; the source is $c\not\equiv 0$.

(iii) As in \cref{kern:thm:onecontext}: the two steps have the same first
premise $s=t$.

(iv) $\rulename{Tr}^{-}$ with the template $J(v)$, applied to $p=0$ and
$J(p)=p$, gives $p=J(0)$, that is $p=c$.  In the same way $q=c$.  By (ii),
$c=p$ and $c=q$.  $\rulename{Tr}^{-}$ with the template $v$, applied to
$c=q$ and $c=p$, gives $p=q$.  The sources are $p$, $q$ and $c$.

(v) If $s\not\equiv 0$, reverse $r=s$ by (ii) and apply $\rulename{Tr}^{-}$
with the template $v$ to $s=t$ and $s=r$.  Let $s\equiv 0$.  If $r\equiv 0$,
the second premise is the conclusion; if $t\equiv 0$, the first premise is.
Otherwise $r\not\equiv 0$ and $t\not\equiv 0$, the premises are $r=0$ and,
after (ii), $t=0$, and (iv) applies.

(vi) If $s\not\equiv 0$, apply $\rulename{Tr}^{-}$ with the template $f$ to
$s=t$ and $f(s)=f(s)$.  If $s\equiv 0$ and $t\not\equiv 0$, reverse the
premise, obtain $f(t)=f(0)$ and reverse again.  If both sides are the
literal $0$, use (i).

(vii) From $s=t$ and $h(s)=r$: (vi) gives $h(s)=h(t)$, (ii) gives $r=h(s)$,
and (v) gives $r=h(t)$.
\end{proof}

The case distinctions in this proof are on the syntax of the given terms and
belong to the construction of the derivation.  The proof uses no induction,
no numerical case distinction, no successor computation and no substitution
into a derivation.  It uses the constant template in (ii) and the crossing
of the conclusion in (ii) and (iv).

\begin{proposition}[Retraction criterion]
\label{kern:prop:retraction}
Let a term language have a literal $0$ and a term context $J(v)$ such that
$J(a)\not\equiv 0$ for every term $a$, and let $J(a)=a$ be available as an
axiom scheme in every context.  Then the crossed rule restricted to sources
other than the literal $0$ derives the crossed rule at all sources, and with
it reflexivity, symmetry, transitivity, congruence and transport in an
equation template.
\end{proposition}

\begin{proof}
The proof of \cref{kern:lem:oseq} uses nothing else.
\end{proof}

\begin{theorem}[Transport at sources other than the literal zero]
\label{kern:thm:onesided}
$\mathsf{O}$, $\mathsf{M}$ and $\Nat$ are equivalent.
\end{theorem}

\begin{proof}
Every instance of $\rulename{Tr}^{-}$ is an instance of \rulename{Tr}, and
\rulename{Tr} is derivable in $\mathsf{O}$ and in $\mathsf{M}$ by
\cref{kern:lem:oseq}(vii).  This proves the equivalence of $\mathsf{O}$ and
$\Nat$.  For $\mathsf{M}$, the rule $\rulename{DI}_{\mathrm{w}}$ is derivable
in $\Nat$ by \cref{kern:thm:witness}.  Conversely the proof of
\cref{kern:thm:witness} derives \rulename{QI} and \rulename{DI} in
$\mathsf{M}$: it uses $\rulename{DI}_{\mathrm{w}}$, \rulename{DE},
\rulename{Hyp}, cut, which is admissible in $\mathsf{M}$ by
\cref{kern:lem:struct}, the equations $c=0$ and $0=c$, and transport in an
equation template along them, which $\mathsf{M}$ derives from \rulename{Tr}
by \cref{kern:lem:local}(iii).
\end{proof}

\begin{remark}[Substitution and the choice of the kernel]
\label{kern:rem:ossubst}
Substituting $0$ for a variable can turn a permitted source into the literal
$0$, so substitution does not map instances of $\rulename{Tr}^{-}$ to
instances.  Substitution is admissible in $\mathsf{O}$ and $\mathsf{M}$ as an
operation on derivable sequents (\cref{kern:prop:transfer}); what fails is
the preservation of each rule instance.  \pkg{reversal\_core} makes the same
correction to an earlier wording about the Lean calculus \Lean{ZPrf}.  The
Lean kernel keeps unrestricted sources so that its substitution theorem is a
direct induction.
\end{remark}

\begin{remark}[Sources]
\label{kern:rem:onesided-sources}
\Cref{kern:lem:oseq,kern:thm:onesided} are from \pkg{one\_sided\_transport},
which calls $\mathsf{M}$ the calculus $\mathcal{C}$ and reserves the letter
$\mathcal{M}$ for the same seven rules with $\rulename{Tr}_{2}$ at sources
other than the literal $0$.  The package proves the equivalence twice: by
simulating that two-template rule with \cref{kern:lem:oseq}(iii) and quoting
an earlier equivalence, and directly as above.  In Lean the statement that
\rulename{I0} and $\rulename{Tr}^{-}$ recover \rulename{Tr} is
\Lean{OKernel.trOneF}, for the class \Lean{OKernel}; the equivalences are
\Lean{oprf\_iff\_nprf} for $\mathsf{O}$ and \Lean{mprf\_iff\_nprf},
\Lean{mprf\_iff\_iprf} for $\mathsf{M}$.  \Cref{kern:prop:retraction} is
stated in \pkg{reversal\_core}.
\end{remark}

\subsection{Reversal of the root equation}
\label{kern:sec:reversal}

Define an operation on the root of a native formula:
\[
 (s=t)^{\dagger}\defeq(t=s),\qquad
 (\Q{y}B)^{\dagger}\defeq\Q{y}B,\qquad
 (\D{x}{y}B)^{\dagger}\defeq\D{x}{y}B .
\]
It does not descend into a matrix, it satisfies $A^{\dagger\dagger}=A$, and
it commutes with substitution.

\begin{definition}[The calculus $\mathsf{X}$]
\label{kern:def:X}
$\mathsf{X}$ has the eight rules \rulename{Hyp}, \rulename{I0},
\rulename{IS}, \rulename{DE}, \rulename{DI}, \rulename{QI} and
\begin{gather*}
 \infrule{\Gamma\proves s=t\qquad\Gamma\proves F(s)}
         {\Gamma\proves F(t)^{\dagger}}\ \rulename{TrX}
 \quad(s\not\equiv 0),
 \\[1.2ex]
 \infrule{\Gamma\proves\Q{y}B(y)\qquad\Gamma\proves B(0)}
         {\Gamma\proves B(t)^{\dagger}}\ \rulename{QX},
\end{gather*}
where $F(v)$ is an arbitrary native formula template, constant templates
included.
\end{definition}

\begin{lemma}[Reversal]
\label{kern:lem:reverse}
If $\Gamma\proves_{\mathsf{X}}A$, then $\Gamma\proves_{\mathsf{X}}A^{\dagger}$.
\end{lemma}

\begin{proof}
Apply \rulename{TrX} with the first premise $c=0$ and the constant template
$F(v)\defeq A$.
\end{proof}

\begin{theorem}[Root-reversing eliminations]
\label{kern:thm:reversal}
$\mathsf{X}$ and $\Nat$ are equivalent.  The derivation of the rules of $\Nat$
in $\mathsf{X}$ uses neither \rulename{IS} nor an introduction rule nor
\rulename{DE}.
\end{theorem}

\begin{proof}
\emph{From $\mathsf{X}$ to $\Nat$.}  For \rulename{TrX}, transport in the
formula template $F$ (\cref{base:sec:structural}) gives $F(t)$; if $F(t)$ is
an equation, symmetry gives $F(t)^{\dagger}$, and otherwise
$F(t)^{\dagger}$ is $F(t)$.  For \rulename{QX}, use \rulename{QE} and the
same step.

\emph{From $\Nat$ to $\mathsf{X}$.}  The rule $\rulename{Tr}^{-}$ is the
instance of \rulename{TrX} at the template $F(v)\defeq\bigl(h(v)=r\bigr)$,
since $F(t)^{\dagger}$ is $r=h(t)$.  Hence \rulename{Tr} is derivable in
$\mathsf{X}$ by \cref{kern:lem:oseq}(vii).  The rule \rulename{QE} is
\rulename{QX} followed by \cref{kern:lem:reverse}; the two reversals cancel.
The six other rules are common.  The constructions hold over arbitrary
contexts, so they apply inside the premises of \rulename{DE} and of the two
introductions.
\end{proof}

Only the equational instances of transport have to be recovered in
$\mathsf{X}$; transport in a compound template is used in the other
direction, where it is available.

\begin{remark}[Sources]
\label{kern:rem:reversal-sources}
\Cref{kern:thm:reversal} is the theorem of \pkg{reversal\_core}, which proves
the items of \cref{kern:lem:oseq} directly for \rulename{TrX}, with the
templates $v=a$, $J(v)=p$, $r=v$ and $f(s)=f(v)$, and derives
$\rulename{Tr}_{2}$ from them.  \pkg{one\_sided\_transport} observes that
$\rulename{Tr}^{-}$ is the restriction of \rulename{TrX} to equations with
one constant side, needs no reversal of arbitrary formulas and leaves
\rulename{QE} unchanged.  In Lean the calculus $\mathsf{X}$ is the module
\Lean{NativeReversal}, adopted from the package, with
\Lean{rprf\_iff\_nprf}, \Lean{reversal\_iff\_pra} and
\Lean{pra\_iff\_reversal}.  The package reports one comparison of proof
sizes: a certificate of the same closed goal, symmetry of an equation with
the source $0$, has 114 instruction packets and 5435 bytes when produced by
the method of \cref{kern:sec:dispatch} and 37 packets and 1259 bytes when
produced with reversal.
\end{remark}

\subsection{The same restriction for the uncrossed rule}
\label{kern:sec:dispatch}

For transport in a formula template with an uncrossed conclusion the
restriction of the source is also harmless, but the proof needs an induction.
Let
\[
 \infrule{\Gamma\proves s=t\qquad\Gamma\proves F(s)}
         {\Gamma\proves F(t)}\ \rulename{Tr}_{F}^{-}
 \qquad(s\not\equiv 0).
\]

\begin{definition}[The calculus $\mathsf{Z}$]
\label{kern:def:Z}
$\mathsf{Z}$ is $\Nat_{F}$ with $\rulename{Tr}_{F}$ replaced by
$\rulename{Tr}_{F}^{-}$; its rules are \rulename{Hyp}, \rulename{I0},
\rulename{IS}, $\rulename{Tr}_{F}^{-}$, \rulename{QE}, \rulename{DE},
\rulename{DI} and \rulename{QI}.
\end{definition}

The sans-serif letter $\mathsf{Z}$ names the calculus; the italic letters
$Z$ and $N$ name the zero test $Z(c)=\sel(c,1,0)$ and the parity term
$N(n)=\It{n}{0}{u}{Z(u)}$ of \cref{base:sec:terms}.

In $\mathsf{Z}$ the two eliminations with a template can be merged into one
rule whose first premise is selected by the source.  For a template $F(v)$
put $M(s,t;F)\defeq\Q{v}F(v)$ if $s\equiv 0$ and $M(s,t;F)\defeq(s=t)$
otherwise; then
\[
 \infrule{\Gamma\proves M(s,t;F)\qquad\Gamma\proves F(s)}
         {\Gamma\proves F(t)}\ \rulename{SE}
\]
is \rulename{QE} when $s$ is the literal $0$ and $\rulename{Tr}_{F}^{-}$
otherwise.  This is the form in which \pkg{source\_dispatch} states the
calculus, and the reason for its name.

\begin{theorem}[Literal-zero sources of uncrossed transport]
\label{kern:thm:dispatch}
$\mathsf{Z}$ and $\Nat$ are equivalent.
\end{theorem}

Every instance of $\rulename{Tr}_{F}^{-}$ is derivable in $\Nat$ by the
Leibniz rule.  For the converse it suffices to derive \rulename{Tr} in
$\mathsf{Z}$; the seven other rules are common.  The proof does not assume
that substitution is admissible in $\mathsf{Z}$: every auxiliary equation is
derived at an arbitrary term by a derived rule applied at that term.  By
\cref{kern:lem:local}(i) and (ii), numerical cases and equation induction
are derivable in $\mathsf{Z}$, with no use of transport.

\begin{lemma}[Equality at permitted sources]
\label{kern:lem:zsafe}
$\mathsf{Z}$ derives: $r=r$; from $p=q$ with $p\not\equiv 0$, the equations
$q=p$ and $f(p)=f(q)$ for every term template $f$; from $p=q$ with
$p\not\equiv 0$ and $p=r$, the equation $q=r$; from $r=s$ and $s=t$ with
$s\not\equiv 0$, the equation $r=t$.
\end{lemma}

\begin{proof}
Transport $J(r)=r$ along itself in the template $v=r$; the source is
$J(r)$.  For the other statements transport, along $p=q$ or along $s=t$,
the equations $p=p$ in the template $v=p$, $f(p)=f(p)$ in the template
$f(p)=f(v)$, $p=r$ in the template $v=r$, and $r=s$ in the template $r=v$.
\end{proof}

\begin{lemma}[Joining through zero]
\label{kern:lem:zjoin}
If $p\not\equiv 0$ and $q\not\equiv 0$, then $\mathsf{Z}$ derives $p=q$ from
$p=0$ and $q=0$.
\end{lemma}

\begin{proof}
Transport $J(p)=J(p)$ along $p=0$ in the template $J(v)=J(p)$; the result is
$c=J(p)$.  Symmetry at the source $c$ gives $J(p)=c$, and transport along
$J(p)=p$ in the template $v=c$ gives $p=c$.  In the same way $q=c$.
Symmetry at the source $q$ gives $c=q$, and transport of $p=c$ along $c=q$
in the template $p=v$ gives $p=q$.  The sources are $p$, $q$, $c$, $J(p)$
and $J(q)$.
\end{proof}

Define, with the selector of \cref{base:sec:terms},
\[
 K(v)\defeq\sel(v,0,1),\qquad H(v)\defeq\sel(v,1,2),\qquad
 R(n)\defeq\It{n}{0}{u}{H(u)},\qquad L(n)\defeq N(R(n)).
\]
In the standard model $R$ takes the values $0,1,2,2,\dots$ and $L$ is $1$ at
$1$ and $0$ elsewhere; the proof derives what it needs and does not use this
description.

\begin{lemma}[A test for the argument one]
\label{kern:lem:ztest}
For every term $v$, $\mathsf{Z}$ derives
\[
 H(v)=\Suc K(v),\quad R(1)=1,\quad R(\Suc\Suc v)=2,\quad
 N(1)=1,\quad N(2)=0,\quad L(1)=1,\quad L(\Suc\Suc v)=0,\quad
 L(\Suc v)=Z(v).
\]
\end{lemma}

\begin{proof}
All chains of equations below pass through middle terms with the root $I$ or
$\Suc$, so each step is one of the operations of \cref{kern:lem:zsafe}.

$H(v)=\Suc K(v)$: by numerical cases on the matrix $H(y)=\Suc K(y)$ at the
term $v$.  At $0$: $H(0)=1$ and $K(0)=0$ by \rulename{I0}, hence
$\Suc K(0)=1$ by congruence at the source $K(0)$, and $H(0)=\Suc K(0)$ by
symmetry at the source $\Suc K(0)$ and transitivity through $1$.  At
$\Suc z$: $H(\Suc z)=2$ and $K(\Suc z)=1$ by \rulename{IS}, hence
$\Suc K(\Suc z)=2$, and the two equations are joined through $2$.

$N(1)=Z(N(0))$ by \rulename{IS}, $N(0)=0$ by \rulename{I0}, hence
$Z(N(0))=Z(0)$ by congruence at the source $N(0)$, and $Z(0)=1$ by
\rulename{I0}; so $N(1)=1$.  Likewise $N(2)=Z(N(1))=Z(1)=0$, the last by
\rulename{IS}.  In the same way $R(1)=H(R(0))=H(0)=1$.

\rulename{IS} and the first equation at the term $R(v)$ give
$R(\Suc v)=H(R(v))=\Suc K(R(v))$.  Then
$R(\Suc\Suc v)=H(R(\Suc v))=H(\Suc K(R(v)))=2$, by \rulename{IS}, congruence
at the source $R(\Suc v)$, and \rulename{IS} for the selector.

$L(1)=N(R(1))=N(1)=1$ and $L(\Suc\Suc v)=N(R(\Suc\Suc v))=N(2)=0$, by
congruence at the sources $R(1)$ and $R(\Suc\Suc v)$.

$L(\Suc v)=Z(v)$: by numerical cases on the matrix $L(\Suc y)=Z(y)$ at the
term $v$.  At $0$, $L(1)=1$ and $Z(0)=1$ are joined through $1$.  At
$\Suc z$, $L(\Suc\Suc z)=0$ and $Z(\Suc z)=0$, the latter by \rulename{IS},
are joined by \cref{kern:lem:zjoin}; both left sides have the root $I$.
\end{proof}

\begin{lemma}[Filter]
\label{kern:lem:zfilter}
For every term $t$, $\mathsf{Z}$ derives $t=\sel(Z(t),t,0)$ in every context.
\end{lemma}

\begin{proof}
By numerical cases on the matrix $y=\sel(Z(y),y,0)$ at the term $t$.  At
$0$: $Z(0)=1$, so $\sel(Z(0),0,0)=\sel(1,0,0)$ by congruence at the source
$Z(0)$, and $\sel(1,0,0)=0$ by \rulename{IS}; symmetry at the left side,
which has the root $I$, gives $0=\sel(Z(0),0,0)$.  At $\Suc z$:
$Z(\Suc z)=0$, so $\sel(Z(\Suc z),\Suc z,0)=\sel(0,\Suc z,0)=\Suc z$ by
congruence at the source $Z(\Suc z)$ and \rulename{I0}, and symmetry at the
left side gives the instance.
\end{proof}

\begin{lemma}[Symmetry at the source zero]
\label{kern:lem:zsym}
For every term $t$ and every context $\Gamma$ containing $0=t$, $\mathsf{Z}$
derives $\Gamma\proves t=0$.
\end{lemma}

\begin{proof}
If $t\equiv 0$, use reflexivity.  Otherwise $t$ is a permitted source.  Put
$A(n)\defeq\It{n}{t}{u}{\Suc u}$ and prove $n=A(n)$ by equation induction on
$n$, with $t$ as a parameter and $0=t$ in the context.

Base.  $A(0)=t$ by \rulename{I0}, and $t=A(0)$ by symmetry at the source
$A(0)$.  Transport of $0=t$ along $t=A(0)$ in the template $0=v$ gives
$0=A(0)$; the source is $t$.

Step.  Assume $z=A(z)$ with $z$ fresh.  The source $z$ is a variable, so
congruence gives $\Suc z=\Suc A(z)$.  \rulename{IS} gives
$A(\Suc z)=\Suc A(z)$, and symmetry at the source $A(\Suc z)$ gives
$\Suc A(z)=A(\Suc z)$.  Transport of $\Suc z=\Suc A(z)$ along this equation
in the template $\Suc z=v$ gives $\Suc z=A(\Suc z)$.

At the term $1$ the induction gives $1=A(1)$.  \rulename{IS} gives
$A(1)=\Suc A(0)$, and congruence at the source $A(0)$ gives
$\Suc A(0)=\Suc t$; so $A(1)=\Suc t$, and transport at the source $A(1)$
gives $1=\Suc t$.  Congruence at the source $1$ gives $L(1)=L(\Suc t)$.  With
$L(1)=1$ and $L(\Suc t)=Z(t)$ from \cref{kern:lem:ztest}, the third
operation of \cref{kern:lem:zsafe} at the source $L(1)$ gives $1=L(\Suc t)$,
and transport at the source $L(\Suc t)$ gives $1=Z(t)$.  Symmetry at the
source $1$ gives $Z(t)=1$.

Transport of the equation $t=\sel(Z(t),t,0)$ of \cref{kern:lem:zfilter}
along $Z(t)=1$ in the template $t=\sel(v,t,0)$ gives $t=\sel(1,t,0)$, and
transport along the instance $\sel(1,t,0)=0$ of \rulename{IS} in the
template $t=v$ gives $t=0$.
\end{proof}

The chain of the proof is
$0=t\ \Rightarrow\ 1=\Suc t\ \Rightarrow\ Z(t)=1\ \Rightarrow\ t=0$.
Numerical cases are used only for equations that hold in every context, and
the one induction that keeps the assumption $0=t$ has the literal equation
$n=A(n)$ as its matrix.

\begin{proof}[Proof of \cref{kern:thm:dispatch}]
Symmetry holds in $\mathsf{Z}$ without restriction: at a source other than
the literal $0$ by \cref{kern:lem:zsafe}, and at the source $0$ by
\cref{kern:lem:zsym} and cut.  Congruence from $s=t$: if $s\not\equiv 0$
use \cref{kern:lem:zsafe}; if $s\equiv t\equiv 0$ use reflexivity; if
$s\equiv 0$ and $t\not\equiv 0$, derive $t=0$, then $f(t)=f(0)$ at the
source $t$, and reverse.  Transitivity from $r=s$ and $s=t$: if
$s\not\equiv 0$ use \cref{kern:lem:zsafe}; if $s\equiv 0$ and one of $r$,
$t$ is the literal $0$, a premise is the conclusion; otherwise reverse
$0=t$ and apply \cref{kern:lem:zjoin} to $r=0$ and $t=0$.  These
constructions are not circular, since symmetry at the source $0$ was proved
before unrestricted congruence and transitivity.  Finally, from $s=t$ and
$h(s)=r$, congruence gives $h(s)=h(t)$, symmetry gives $r=h(s)$ and
transitivity gives $r=h(t)$, which is \rulename{Tr}.
\end{proof}

\begin{remark}[Sources and comparison]
\label{kern:rem:dispatch-sources}
\Cref{kern:thm:dispatch} and its lemmas are from \pkg{source\_dispatch},
whose calculus has \rulename{DP} in place of \rulename{QI}
(\cref{kern:thm:direct}).  The package also derives unrestricted transport
in a formula template inside its calculus, by induction on the template;
here that step is the Leibniz rule of $\Nat$, since the equivalence needs
\rulename{Tr} only.  In Lean the theorem is \Lean{zprf\_iff\_nprf}, and the
recovery of unrestricted transport is \Lean{ZKernel.trF} for a class
\Lean{ZKernel} that has equation induction and no logical rule.  The
contrast with \cref{kern:lem:oseq} is the point of the two later packages:
for the crossed rule, and for the reversing rule, a constant template
reverses an equation in one step, and neither the test terms nor an
induction is needed.
\end{remark}

\subsection{Replacement with no template}
\label{kern:sec:tfree}

A transport rule can be given no template at all.  For terms $s$, $t$, $u$
let $\operatorname{Rep}_{s}^{t}(u)$ be the result of one traversal of $u$
from the root: a subterm that is identical with $s$ is replaced by $t$ and
is not traversed further; at any other subterm the constructor is kept and
its immediate subterms are traversed; a variable or a $0$ that is not
replaced is unchanged.  Under the binder of an iterator the bound variable is
renamed away from $s$ and $t$, so that only free occurrences of $s$ are
matched and no variable of $t$ is captured.  The test at a subterm precedes
the descent, inserted copies of $t$ are not traversed, and the operation is
total even when $t$ contains $s$.  It is one pass and not rewriting to a
normal form.  In particular $\operatorname{Rep}_{s}^{t}(s)=t$.

\begin{definition}[The calculi $\mathsf{G}$ and $\mathsf{G}_7$]
\label{kern:def:G}
The rule of global replacement is
\[
 \infrule{\Gamma\proves s=t\qquad\Gamma\proves u=r}
         {\Gamma\proves r=\operatorname{Rep}_{s}^{t}(u)}\ \rulename{Rep}.
\]
$\mathsf{G}$ is $\Nat$ with \rulename{Tr} replaced by \rulename{Rep}; its
rules are \rulename{Hyp}, \rulename{I0}, \rulename{IS}, \rulename{Rep},
\rulename{QE}, \rulename{DE}, \rulename{DI} and \rulename{QI}.
$\mathsf{G}_7$ is $\mathsf{W}$ with \rulename{Tr} replaced by \rulename{Rep};
its seven rules are \rulename{Hyp}, \rulename{I0}, \rulename{IS},
\rulename{Rep}, \rulename{QE}, \rulename{DE} and
$\rulename{DI}_{\mathrm{w}}$.
\end{definition}

The rule cannot choose which occurrences of $s$ to replace.  A chosen
template is recovered by planting a marker that cannot occur by accident.

\begin{lemma}[Equality from \rulename{I0} and \rulename{Rep}]
\label{kern:lem:geq}
From \rulename{I0} and \rulename{Rep} alone the following are derivable:
from $s=t$ and $s=u$ the equation $u=t$; reflexivity; symmetry;
transitivity.
\end{lemma}

\begin{proof}
\rulename{Rep} applied to $s=t$ and $s=u$ concludes
$u=\operatorname{Rep}_{s}^{t}(s)$, which is $u=t$.  This joining, applied to
$J(a)=a$ and $J(a)=a$, gives $a=a$; applied to $a=a$ and $a=b$ it gives
$b=a$.  For transitivity let $a=b$ and $b=e$ be given: symmetry gives
$b=a$, and the joining applied to $b=e$ and $b=a$ gives $a=e$.
\end{proof}

Only matches at the root are used in this proof.  Let $|T|$ be the number of
nodes of a term $T$, the body of every iterator included; renaming of
variables does not change it.

\begin{lemma}[Exact marking]
\label{kern:lem:marker}
Let $h(v)$ be a term template, $s$ and $t$ terms, and
\[
 n\defeq 1+|h|+|s|+|t|,\qquad
 K\defeq\It{0}{s}{w}{\Suc^{n}0}.
\]
Then the occurrences of $K$ in $h(K)$ are exactly the copies inserted for
$v$, and
\[
 \operatorname{Rep}_{K}^{s}\bigl(h(K)\bigr)=h(s),\qquad
 \operatorname{Rep}_{K}^{t}\bigl(h(K)\bigr)=h(t).
\]
\end{lemma}

\begin{proof}
$K$ has more nodes than $h$.  Two subterm occurrences of a term are disjoint
or nested.  An occurrence in $h(K)$ of a subterm identical with $K$ that is
disjoint from all inserted copies lies in a part of $h$ and has at most
$|h|$ nodes.  One that lies properly inside an inserted copy has fewer nodes
than $K$.  One that properly contains an inserted copy has more nodes than
$K$.  None of these is identical with $K$, so the occurrences are the
inserted copies.  Below a binder of $h$ the inserted copy and the term
compared with it are subject to the same renaming, which preserves the
number of nodes.  The traversal replaces exactly these copies and stops at
each, which gives the two identities.
\end{proof}

The identities are facts about syntax; no rule of the calculus is applied in
them.  The term $K$ keeps the free variables of $s$, and $K=s$ is an
instance of \rulename{I0}.

\begin{lemma}[Chosen templates from global replacement]
\label{kern:lem:gcong}
From \rulename{I0} and \rulename{Rep} alone: from $s=t$ the equation
$h(s)=h(t)$, for every term template $h$; and the rule \rulename{Tr}.
\end{lemma}

\begin{proof}
Let $K$ be as in \cref{kern:lem:marker}.  From $K=s$ and $s=t$,
transitivity gives $K=t$.  Reflexivity gives $h(K)=h(K)$.  \rulename{Rep}
applied to $K=s$ and $h(K)=h(K)$ concludes $h(K)=h(s)$, and applied to $K=t$
and $h(K)=h(K)$ it concludes $h(K)=h(t)$, by \cref{kern:lem:marker}.  The
joining of \cref{kern:lem:geq}, applied to $h(K)=h(t)$ and $h(K)=h(s)$,
gives $h(s)=h(t)$.  For \rulename{Tr}, from $s=t$ and $h(s)=r$: symmetry
gives $r=h(s)$ and transitivity gives $r=h(t)$.
\end{proof}

For example let $s=t$ be given and let $h(v)\defeq\It{v}{s}{w}{s}$, in which
the two displayed copies of $s$ are to stay.  Replacement of $s$ in $h(s)$
would change all three copies.  In $h(K)=\It{K}{s}{w}{s}$ the marker is too
large to occur in the two fixed copies, and the two replacements produce
$\It{s}{s}{w}{s}$ and $\It{t}{s}{w}{s}$.

\begin{lemma}[\rulename{Rep} in $\Nat$]
\label{kern:lem:gsound}
Every instance of \rulename{Rep} is an instance of \rulename{Tr}.
\end{lemma}

\begin{proof}
The occurrences of $s$ that the traversal of $u$ replaces are disjoint.
Replace them by a fresh variable $v$ to obtain a term template $h(v)$; then
$u\equiv h(s)$ and $\operatorname{Rep}_{s}^{t}(u)\equiv h(t)$.  If no
occurrence is replaced, $h$ is constant.  \rulename{Tr} with this template,
applied to $s=t$ and $h(s)=r$, concludes $r=h(t)$.
\end{proof}

Thus \rulename{Rep} is \rulename{Tr} with the template computed from the two
premises.

\begin{theorem}[Replacement with no template]
\label{kern:thm:tfree}
$\mathsf{G}$, $\mathsf{G}_7$ and $\Nat$ are equivalent.
\end{theorem}

\begin{proof}
\rulename{Rep} is derivable in $\Nat$ by \cref{kern:lem:gsound}, and
\rulename{Tr} is derivable in $\mathsf{G}$ and in $\mathsf{G}_7$ by
\cref{kern:lem:gcong}.  This settles $\mathsf{G}$.  For $\mathsf{G}_7$ the
introductions are treated as in \cref{kern:thm:witness}: an instance of
$\rulename{DI}_{\mathrm{w}}$ is derivable in $\Nat$, and \rulename{QI} and
\rulename{DI} are derivable in $\mathsf{G}_7$ by the proof given there, with
the alias $c$, transport in an equation template derived from \rulename{Tr},
and cut, which is admissible in $\mathsf{G}_7$ by \cref{kern:lem:struct}.
\end{proof}

\begin{remark}[Substitution and certificates]
\label{kern:rem:tfree}
A substitution can create new occurrences of the source in the second
premise of \rulename{Rep}, so it does not map instances to instances;
injective renaming does.  Substitution is admissible through
\cref{kern:prop:transfer}.  An instance of \rulename{Rep} in a certificate
consists of its two premise derivations; the price is paid in the
derivations, which grow by the marker.
\end{remark}

\begin{remark}[Sources]
\label{kern:rem:tfree-sources}
The section is \pkg{template\_free}, whose calculus is $\mathsf{G}_7$.  In
Lean the calculus $\mathsf{G}$ is \Lean{TPrf}, with \Lean{tprf\_iff\_nprf};
the combination $\mathsf{G}_7$ with the single introduction is not
formalized.  The compact implementation of the package works with the
restricted inventory of \cref{kern:lem:comprepair} and uses the aliases
$C_s\defeq\It{2}{1}{u}{s}$ and $K\defeq\It{0}{C_s}{u}{\Suc^{n}0}$, so that
$C_s=s$ is a permitted instance of \rulename{IS} and $K=C_s$ a permitted
instance of \rulename{I0}, also when $s$ is the literal $0$.
\end{remark}

\subsection{Backward replacement}
\label{kern:sec:brep}

The rule \rulename{Rep} computes its conclusion from its two premises.  Its
goal-directed counterpart computes its second premise from its conclusion and
its first premise.  It uses the operation $\operatorname{Rep}$ of
\cref{kern:sec:tfree} with the two sides of the first premise exchanged:
$\operatorname{Rep}_{t}^{s}(u)$ is the result of one traversal of $u$ from
the root in which every subterm identical with $t$ is replaced by $s$ and is
not traversed further, the bound variables of $u$ being renamed away from $s$
and $t$.  Inserted copies of $s$ are not traversed, so the operation is total
also when $s$ contains $t$, and $\operatorname{Rep}_{t}^{s}(t)\equiv s$.

\begin{definition}[The calculi $\mathsf{B}$ and $\mathsf{B}_{7}$]
\label{kern:def:B}
The rule of backward replacement is
\[
 \infrule{\Gamma\proves s=t\qquad
          \Gamma\proves\operatorname{Rep}_{t}^{s}(u)=r}
         {\Gamma\proves r=u}\ \rulename{BRep}.
\]
$\mathsf{B}$ is $\Nat$ with \rulename{Tr} replaced by \rulename{BRep}; its
rules are \rulename{Hyp}, \rulename{I0}, \rulename{IS}, \rulename{BRep},
\rulename{QE}, \rulename{DE}, \rulename{DI} and \rulename{QI}.
$\mathsf{B}_{7}$ is $\mathsf{W}$ with \rulename{Tr} replaced by
\rulename{BRep}; its seven rules are \rulename{Hyp}, \rulename{I0},
\rulename{IS}, \rulename{BRep}, \rulename{QE}, \rulename{DE} and
$\rulename{DI}_{\mathrm{w}}$.
\end{definition}

For a conclusion $r=u$ and a first premise $s=t$ the second premise is
determined.  An instance carries no template, no selection of occurrences
and no term besides those of its conclusion and of its first premise.  The
source $s$ may be the literal $0$, and $t$ may occur in $s$.  Matching is a
test on syntax: a bound variable of $u$ never matches a free variable of $t$.

Among the transport rules of the section, \rulename{BRep} is closest to
\rulename{Rep}: both are instances of \rulename{Tr} whose template is
computed by one traversal (\cref{kern:lem:gsound,kern:lem:bsound}), and
neither can choose which occurrences to replace.  They differ in what is
matched and where the result goes.  \rulename{Rep} matches the source $s$ in
the left side $u$ of its second premise and puts
$\operatorname{Rep}_{s}^{t}(u)$ into the conclusion.  \rulename{BRep} matches
the target $t$ in the right side $u$ of its conclusion and puts
$\operatorname{Rep}_{t}^{s}(u)$ into the second premise.  A checker that works
from the goal backwards therefore computes the second premise of
\rulename{BRep} from the goal and the first premise, whereas for
\rulename{Rep} it has to be given the term $u$, which the conclusion
$r=\operatorname{Rep}_{s}^{t}(u)$ does not determine.  Neither rule is an
instance of the other: \rulename{Rep} applied to $s=t$ and
$\operatorname{Rep}_{t}^{s}(u)=r$ concludes
$r=\operatorname{Rep}_{s}^{t}(\operatorname{Rep}_{t}^{s}(u))$, which is $t$
and not $u$ when $u\equiv s$ and $t$ does not occur in $s$.  In the proofs
below the marker of \cref{kern:lem:marker} is used unchanged, with the first
premises $s=K$ and $t=K$ in place of the first premises $K=s$ and $K=t$ of
\cref{kern:lem:gcong}.

\begin{lemma}[\rulename{BRep} in $\Nat$]
\label{kern:lem:bsound}
Every instance of \rulename{BRep} is an instance of \rulename{Tr}.
\end{lemma}

\begin{proof}
The occurrences of $t$ that the traversal of $u$ replaces are disjoint, and
after the renaming of the bound variables of $u$ away from $s$ and $t$ they
are free occurrences.  Replace them by a fresh variable $v$ to obtain a term
template $h(v)$; then $u\equiv h(t)$ and
$\operatorname{Rep}_{t}^{s}(u)\equiv h(s)$, and no variable of $s$ is
captured in $h(s)$.  If no occurrence is replaced, $h$ is constant.
\rulename{Tr} with this template, applied to $s=t$ and $h(s)=r$, concludes
$r=h(t)$, which is $r=u$.
\end{proof}

Thus \rulename{BRep} is \rulename{Tr} with the template computed from the
conclusion, as \rulename{Rep} is \rulename{Tr} with the template computed from
the second premise.

\begin{lemma}[Equality from \rulename{I0} and \rulename{BRep}]
\label{kern:lem:beq}
From \rulename{I0} and \rulename{BRep} alone the following are derivable over
an arbitrary context: from $s=t$ and $s=r$ the equation $r=t$; reflexivity;
symmetry; transitivity.
\end{lemma}

\begin{proof}
\rulename{BRep} with the first premise $s=t$ and the conclusion $r=t$ has the
second premise $\operatorname{Rep}_{t}^{s}(t)=r$, which is $s=r$, since the
traversal matches at the root, also when one of $s$, $t$ contains the other.
This is the joining of \cref{kern:lem:geq}, from $s=t$ and $s=u$ to $u=t$,
with $r$ in place of $u$.  The proof of \cref{kern:lem:geq} uses
\rulename{Rep} only through this joining and otherwise only \rulename{I0},
so it gives reflexivity, symmetry and transitivity here as well.
\end{proof}

\begin{lemma}[Chosen templates from backward replacement]
\label{kern:lem:bcong}
From \rulename{I0} and \rulename{BRep} alone: from $s=t$ the equation
$h(s)=h(t)$, for every term template $h$; and the rule \rulename{Tr}.
\end{lemma}

\begin{proof}
Let $K$ be as in \cref{kern:lem:marker}.  $K=s$ is an instance of
\rulename{I0}, and symmetry gives $s=K$.  Transitivity of $K=s$ and $s=t$
gives $K=t$, and symmetry gives $t=K$ (\cref{kern:lem:beq}).
\rulename{BRep} with the first premise $s=K$ and the conclusion $h(s)=h(K)$
has the second premise $\operatorname{Rep}_{K}^{s}(h(K))=h(s)$, which is
$h(s)=h(s)$ by \cref{kern:lem:marker}; reflexivity provides it, and the
conclusion is $h(s)=h(K)$.  In the same way the first premise $t=K$ and
reflexivity of $h(t)$ give $h(t)=h(K)$, since
$\operatorname{Rep}_{K}^{t}(h(K))\equiv h(t)$.  Symmetry and transitivity
give $h(s)=h(t)$.  For \rulename{Tr}, from $s=t$ and $h(s)=r$: the joining of
\cref{kern:lem:beq}, applied to $h(s)=h(t)$ and $h(s)=r$, gives $r=h(t)$.
\end{proof}

The proof covers constant templates, templates with several occurrences of
$v$ and occurrences below any number of binders.  For the template
$h(v)\defeq\It{v}{s}{w}{s}$ of the example after \cref{kern:lem:gcong}, the
two steps with the first premises $s=K$ and $t=K$ act on
$h(K)=\It{K}{s}{w}{s}$ and leave the two fixed copies of $s$ alone.  The
derivation of \rulename{Tr} uses no logical rule, no successor computation
and no induction.

\begin{theorem}[Backward replacement]
\label{kern:thm:brep}
$\mathsf{B}$, $\mathsf{B}_{7}$ and $\Nat$ are equivalent.
\end{theorem}

\begin{proof}
\rulename{BRep} is derivable in $\Nat$ by \cref{kern:lem:bsound}, and
\rulename{Tr} is derivable in $\mathsf{B}$ and in $\mathsf{B}_{7}$ by
\cref{kern:lem:bcong}.  The six other rules of $\mathsf{B}$ are those of
$\Nat$, which settles $\mathsf{B}$.  For $\mathsf{B}_{7}$ the introductions
are treated as in the proof of \cref{kern:thm:tfree}: an instance of
$\rulename{DI}_{\mathrm{w}}$ is derivable in $\Nat$ by
\cref{kern:thm:witness}, and \rulename{QI} and \rulename{DI} are derivable in
$\mathsf{B}_{7}$ by the proof of \cref{kern:thm:witness}, with the alias $c$,
transport in an equation template derived from \rulename{Tr}
(\cref{kern:lem:local}(iii)), and cut, which is admissible in
$\mathsf{B}_{7}$ by \cref{kern:lem:struct}.  Each translation replaces one
inference by a derivation of the same sequent from the same premises over
the same context; the marker has no free variable outside $s$, so the
conditions on the eigenvariables of enclosing rules are kept, and no
admissibility of substitution in $\mathsf{B}$ is assumed.
\end{proof}

A substitution can create new occurrences of the target in the right side of
the conclusion of \rulename{BRep} that are not matched in its second premise,
so, as for \rulename{Rep} (\cref{kern:rem:tfree}), substitution does not map
instances to instances, and it is admissible through
\cref{kern:prop:transfer}.

\paragraph{Sources other than the literal zero.}  Let
$\rulename{BRep}^{-}$ be \rulename{BRep} with the side condition
$s\not\equiv 0$ on the source, the term that the rule inserts.
$\mathsf{B}^{-}$ is $\mathsf{B}$ and $\mathsf{B}^{-}_{7}$ is
$\mathsf{B}_{7}$ with \rulename{BRep} replaced by $\rulename{BRep}^{-}$.
For the crossed rule with a template the restriction is lifted by
\cref{kern:prop:retraction}, with a constant template and the first premise
$c=0$.  That argument is not available here: the first premise $c=0$ has the
target $0$, and $\operatorname{Rep}_{0}^{c}(u)$ replaces every leaf $0$ of
$u$, so the template is constant only when $u$ contains no $0$; and the
marker of \cref{kern:lem:bcong} needs the first premise $t=K$, which is
excluded when $t\equiv 0$.  The recovery below uses the logical rules.  The
terms $J(a)$ and $c=J(0)=\It{0}{0}{u}{0}$ are those of the beginning of the
section.

\begin{lemma}[Equality in $\mathsf{B}^{-}$]
\label{kern:lem:bmeq}
In $\mathsf{B}^{-}$, and in every calculus in which \rulename{I0},
\rulename{IS}, \rulename{QE}, \rulename{QI} and $\rulename{BRep}^{-}$ are
derivable, the following are derivable over an arbitrary context.
\begin{enumerate}[label=\textup{(\roman*)}]
\item Reflexivity; and joining from $s=t$ and $s=r$ to $r=t$ when
 $s\not\equiv 0$.
\item Symmetry, also when a side is the literal $0$.
\item $\It{p}{p}{v}{p}=p$ for every term $p$ and every $v$ not free in $p$.
\item From $p=0$ the equation $c=p$.
\item Joining and transitivity without restriction.
\item From $s=t$ with $s\not\equiv 0$ and $t\not\equiv 0$, the equation
 $h(s)=h(t)$ for every term template $h$.
\end{enumerate}
\end{lemma}

\begin{proof}
(i) $\rulename{BRep}^{-}$ with the first premise $J(a)=a$, an instance of
\rulename{I0} with the source $J(a)\not\equiv 0$, and the conclusion $a=a$ has
the second premise $\operatorname{Rep}_{a}^{J(a)}(a)=a$, which is $J(a)=a$
again.  Joining at a source $s\not\equiv 0$ is the joining of
\cref{kern:lem:beq}.

(ii) Let $u=r$ be given, and let $M\defeq\Suc^{k}0$ with $k\defeq|u|$, so
that $M$ is closed and has more nodes than $u$.  $J(M)=M$ is an instance of
\rulename{I0} with the source $J(M)\not\equiv 0$.  No subterm of $u$, also
below a binder, is identical with $M$, so
$\operatorname{Rep}_{M}^{J(M)}(u)\equiv u$, and $\rulename{BRep}^{-}$ with the
first premise $J(M)=M$ and the conclusion $r=u$ has the given equation $u=r$
as its second premise.

(iii) Let $y$ be fresh and $E(y)\defeq\bigl(\It{y}{p}{v}{p}=p\bigr)$, a
literal equation.  $E(\Suc z)$ is an instance of \rulename{IS}, because
$p\subst{\It{z}{p}{v}{p}}{v}\equiv p$; so $\Gamma,E(z)\proves E(\Suc z)$, the
assumption unused, and \rulename{QI} gives $\Q{y}E(y)$.  $E(0)$ is an
instance of \rulename{I0}, and \rulename{QE} at the term $p$ gives $E(p)$,
which is $\It{p}{p}{v}{p}=p$.

(iv) If $p\equiv 0$, then $c=p$ is an instance of \rulename{I0}.  Otherwise
$\rulename{BRep}^{-}$ with the first premise $p=0$, whose source is $p$, and
the conclusion $p=c$ has the second premise
$\operatorname{Rep}_{0}^{p}(c)=p$.  The traversal of $c=\It{0}{0}{u}{0}$ does
not match at the root, since $c\not\equiv 0$, and replaces the counter, the
seed and the body, each the literal $0$, by $p$, the bound variable being
renamed away from $p$.  So the second premise is $\It{p}{p}{u}{p}=p$, which is
(iii).  The conclusion is $p=c$, and (ii) gives $c=p$.

(v) Joining from $s=t$ and $s=r$: if $s\not\equiv 0$, use (i).  If
$s\equiv 0$, (ii) gives $t=0$ and $r=0$, (iv) gives $c=t$ and $c=r$, and
joining at the source $c\not\equiv 0$ gives $r=t$.  Transitivity from $a=b$
and $b=e$: (ii) gives $b=a$, and joining applied to $b=e$ and $b=a$ gives
$a=e$.

(vi) The proof of \cref{kern:lem:bcong}: its two applications of
\rulename{BRep} have the sources $s$ and $t$, and its other steps are
(i), (ii) and (v).
\end{proof}

\begin{lemma}[Congruence at a zero endpoint]
\label{kern:lem:bzero}
In $\mathsf{B}^{-}$ and in $\mathsf{B}^{-}_{7}$, from $0=t$ the equation
$h(0)=h(t)$ is derivable over an arbitrary context, for every term template
$h$.
\end{lemma}

\begin{proof}
If $t\equiv 0$, use reflexivity; let $t\not\equiv 0$.  By the proof of
\cref{kern:thm:witness}, the rule \rulename{QI} is derivable in
$\mathsf{B}^{-}_{7}$ with $\rulename{DI}_{\mathrm{w}}$, \rulename{DE} and
\rulename{Hyp} and without an equality rule, so \cref{kern:lem:bmeq} applies
to both calculi.

\emph{A variable equal to zero.}  Let $b$ be a variable not free in $h$, and
let $b=0$ be derivable over a context $\Delta$.  Put
$L\defeq\operatorname{Rep}_{0}^{b}(h(0))$.  Every occurrence of the term $0$
is a leaf.  In $h(0)$ the leaves $0$ are those of $h$ and the inserted
copies; in $h(b)$ they are those of $h$, since the variable $b$ contains no
$0$.  Hence $\operatorname{Rep}_{0}^{b}(h(b))\equiv L$ as well: both are $h$
with $v$ and every leaf $0$ replaced by $b$, under the same renaming of bound
variables away from $b$.  $\rulename{BRep}^{-}$ with the first premise $b=0$,
whose source is a variable, and the conclusion $L=h(0)$ has the second
premise $L=L$; with the conclusion $L=h(b)$ it has the second premise $L=L$
as well.  Reflexivity gives $L=h(0)$ and $L=h(b)$ over $\Delta$, and joining
and symmetry (\cref{kern:lem:bmeq}) give $h(0)=h(b)$.

\emph{A name for zero.}  Let $C(y,a,d)\defeq\It{y}{a}{v}{d}$ with $v$ not
free in $d$, and $E(x,y)\defeq\bigl(C(y,0,x)=C(y,0,0)\bigr)$.  The formula
$E(0,\Suc z)$ is an instance of reflexivity.  Hence \rulename{DI} with the
witness $0$ and the premise $\Gamma,E(0,0)\proves E(0,\Suc z)$, or in
$\mathsf{B}^{-}_{7}$ the rule $\rulename{DI}_{\mathrm{w}}$ with the witness
$0$, which is in the mode of progress and has the premise
$\Gamma,E(0,z)\proves E(0,\Suc z)$, gives $\Gamma\proves\D{x}{y}E(x,y)$; in
both the assumption is unused.  Let $b$ be fresh for $\Gamma$, $h$ and $t$,
and let $\Delta\defeq\Gamma,\Q{y}E(b,y)$.  Over $\Delta$:
$C(0,0,b)=0$ and $C(0,0,0)=0$ are instances of \rulename{I0}, so symmetry
and transitivity give $E(b,0)$; \rulename{Hyp} and \rulename{QE} at the term
$\Suc 0$ give $E(b,\Suc 0)$, that is $C(\Suc 0,0,b)=C(\Suc 0,0,0)$; the
instances $C(\Suc 0,0,b)=b$ and $C(\Suc 0,0,0)=0$ of \rulename{IS}, symmetry
and transitivity give $b=0$.  Only equations at the roots are compared;
no congruence is used.

\emph{Completion.}  Over $\Delta$, the equation $0=t$ is derivable by
weakening (\cref{kern:lem:struct}), and transitivity with $b=0$ gives $b=t$.
Since $b\not\equiv 0$ and $t\not\equiv 0$, \cref{kern:lem:bmeq}(vi) gives
$h(b)=h(t)$.  The first step gives $h(0)=h(b)$, and transitivity gives
$h(0)=h(t)$ over $\Delta$.  The variable $b$ occurs neither in $\Gamma$ nor in
$h(0)=h(t)$, so \rulename{DE} with the first premise $\D{x}{y}E(x,y)$
concludes $\Gamma\proves h(0)=h(t)$.
\end{proof}

\begin{theorem}[Backward replacement at sources other than the literal zero]
\label{kern:thm:bminus}
$\mathsf{B}^{-}$, $\mathsf{B}^{-}_{7}$, $\mathsf{B}$ and $\Nat$ are
equivalent.
\end{theorem}

\begin{proof}
Every instance of $\rulename{BRep}^{-}$ is an instance of \rulename{BRep}, so
$\mathsf{B}^{-}$ and $\mathsf{B}^{-}_{7}$ derive only sequents of $\Nat$ by
\cref{kern:thm:brep}.  Conversely, congruence from $s=t$ is derivable in both:
if $s\not\equiv 0$ and $t\not\equiv 0$ by \cref{kern:lem:bmeq}(vi); if
$s\equiv 0$ by \cref{kern:lem:bzero}; if $t\equiv 0\not\equiv s$, symmetry
gives $0=s$, \cref{kern:lem:bzero} gives $h(0)=h(s)$, and symmetry gives
$h(s)=h(0)$.  From $s=t$ and $h(s)=r$, congruence gives $h(s)=h(t)$ and
joining (\cref{kern:lem:bmeq}(v)) gives $r=h(t)$, which is \rulename{Tr}.  In
$\mathsf{B}^{-}$ the other rules are those of $\Nat$.  In
$\mathsf{B}^{-}_{7}$, \rulename{QI} was derived before any equality rule
was needed, and \rulename{DI} with a literal equation and the witness $0$ is
derived after \rulename{Tr} by the proof of \cref{kern:thm:witness}, with the
alias $c$, transport in an equation template and cut.  The order of the
derivations is not circular: \cref{kern:lem:bmeq,kern:lem:bzero} use
$\rulename{DI}_{\mathrm{w}}$ only in the mode of progress with the witness
$0$, which needs no equality rule, and the instance of \rulename{DI} that
needs transport is derived only after \rulename{Tr}.
\end{proof}

\begin{remark}[Sources]
\label{kern:rem:brep-sources}
The subsection is \pkg{backward\_replacement}, which presents itself as an
executable demonstrator of correctness for $\mathsf{B}$ and states
\cref{kern:thm:brep,kern:thm:bminus}; it takes the marker and the bridge to
$\iPRA$ from \pkg{template\_free} and \cref{base:thm:bridge}, and finds no
earlier statement of \rulename{BRep} in the corpus.  The Lean module
\Lean{NativeReplacementCalculi} formalizes $\mathsf{B}$ and $\mathsf{B}_{7}$
(\Lean{bprf\_iff\_nprf}, \Lean{b7prf\_iff\_nprf}); the equality rules and
\rulename{Tr} are derived from \rulename{I0} and \rulename{BRep} generically
(\Lean{BRKernel}), and the introductions come from \Lean{WKernel}.  The
module \Lean{NativeNonzeroSources} formalizes the restricted calculi
$\mathsf{B}^{-}$ and $\mathsf{B}^{-}_{7}$ (\Lean{bmprf\_iff\_nprf},
\Lean{bm7prf\_iff\_nprf}) with the bootstrap of
\cref{kern:lem:bmeq,kern:lem:bzero}, generically over \Lean{BMKernel}.  The
checker of the package implements $\mathsf{B}$ with unrestricted sources; in
the package $\mathsf{B}_{7}$, $\mathsf{B}^{-}$ and $\mathsf{B}^{-}_{7}$ are
written proofs only.  Its tests are finite: 18
accepted open or closed sequents, with instances of all eight rules, zero
and nonzero sources and five generated certificates of congruence by the
marker, whose templates have fixed copies of the source, a repeated
variable, an unused variable and one or two nested binders; 16 rejected
certificates, among them both kinds of capture, an escaping eigenvariable,
wrong computation equations, a wrong instance of a closure and progress at
a compound matrix; 7 direct tests of capture avoidance under $\D{}{}$,
$\Q{}$ and $I$ and of the absence of repeated replacement inside an inserted
source; and 18 runs of the driver.  An independent review in the package
adds four accepted sequents, among them a closed certificate that
introduces, uses and discharges the name for zero of
\cref{kern:lem:bzero}, two rejected certificates and four runs of the
driver.  These tests concern particular certificates; the equivalence rests
on the proofs above.
\end{remark}

\subsection{Replacement of one occurrence}
\label{kern:sec:onehole}

A \emph{term context with one hole} $C[\,]$ is a term with exactly one
occurrence of a hole; the hole may lie in the body of an iterator, and
$C[s]$ plugs $s$ into the hole after the bound variables of $C$ have been
renamed away from $s$.  The context consisting of the hole alone is allowed.

\begin{definition}[The calculi $\mathsf{L}$ and $\Nat_{\mathrm{u}}$]
\label{kern:def:L}
$\mathsf{L}$ is $\Nat$ with \rulename{Tr} replaced by
\[
 \infrule{\Gamma\proves s=t\qquad\Gamma\proves C[s]=r}
         {\Gamma\proves C[t]=r}\ \rulename{LOR},
\]
for a term context $C$ with one hole.  $\Nat_{\mathrm{u}}$ is $\Nat$ with
\rulename{Tr} replaced by the same rule for an arbitrary term template
$h(v)$ in place of $C$: from $s=t$ and $h(s)=r$ infer $h(t)=r$.  Each has
the rules \rulename{Hyp}, \rulename{I0}, \rulename{IS}, its replacement
rule, \rulename{QE}, \rulename{DE}, \rulename{DI} and \rulename{QI}.
\end{definition}

\begin{theorem}[One occurrence suffices]
\label{kern:thm:onehole}
$\mathsf{L}$, $\Nat_{\mathrm{u}}$ and $\Nat$ are equivalent.
\end{theorem}

\begin{proof}
An instance of \rulename{LOR} is an instance of the rule of
$\Nat_{\mathrm{u}}$, with the template obtained by putting a fresh variable
into the hole.  The rule of $\Nat_{\mathrm{u}}$ is derivable in $\Nat$:
\rulename{Tr} gives $r=h(t)$ and symmetry gives $h(t)=r$.

It remains to derive \rulename{Tr} in $\mathsf{L}$.  \rulename{LOR} with the
context consisting of the hole, applied to $J(a)=a$ and $J(a)=a$, gives
$a=a$; applied to $s=t$ and $s=s$ it gives $t=s$.  Now let $s=t$ and
$h(s)=r$ be given.  Rename the bound variables of $h$ away from $s$ and $t$
and number the occurrences of $v$ in $h$ from $1$ to $k$.  Let $h_j$ be $h$
with the first $j$ occurrences replaced by $t$ and the others by $s$, and
let $C_j$ be $h$ with the first $j-1$ occurrences replaced by $t$, the
$j$-th by the hole and the others by $s$.  Then $h_{j-1}\equiv C_j[s]$ and
$h_j\equiv C_j[t]$, and $C_j$ has one hole.  Hence $k$ applications of
\rulename{LOR} lead from $h(s)=r$, which is $h_0=r$, to $h_k=r$, which is
$h(t)=r$, and symmetry gives $r=h(t)$.  The occurrences are those of the
template variable; a subterm of $h$ that happens to be identical with $s$ is
part of every $C_j$ and is not replaced.
\end{proof}

\begin{remark}[Sources]
\label{kern:rem:onehole-sources}
The theorem for $\mathsf{L}$ is the auxiliary note of
\pkg{kernel\_refinement}, which states it against $\Nat_{2}$: the occurrences
in the left template are replaced one at a time, the equation is reversed,
the occurrences in the right template are replaced, and the equation is
reversed again.  It answers the question of \pkg{one\_context} whether a
restriction to one occurrence is adequate under the binder of the iterator.
In Lean, $\mathsf{L}$ is \Lean{LPrf} with \Lean{lprf\_iff\_nprf}; contexts
with one hole are \Lean{Tm.Hole}, and the chain $h_0,\dots,h_k$ is
\Lean{Reach.subst}, where below a binder the hole and the bound variable are
exchanged (\Lean{Reach.iterH}).  The calculus $\Nat_{\mathrm{u}}$ is
\Lean{UPrf} with \Lean{uprf\_iff\_nprf}.  The package measures no checker for
$\mathsf{L}$: a traversal by position has a cost of its own, and the compact
checkers share one opening routine between transport and \rulename{QE}.
\end{remark}

\subsection{A restricted inventory of computation instances}
\label{kern:sec:comprepair}

Several compact checkers of the packages accept a computation equation only
when its seed, and for \rulename{IS} also the predecessor of its counter, is
not the literal $0$.  Call an instance of \rulename{I0} \emph{permitted} if
$a\not\equiv 0$, and an instance of \rulename{IS} permitted if
$n\not\equiv 0$ and $a\not\equiv 0$.  Let $c_{*}\defeq\It{2}{1}{u}{0}$.

\begin{lemma}[Repair of computation instances]
\label{kern:lem:comprepair}
Let a calculus have the permitted instances of \rulename{I0} and
\rulename{IS} and a transport rule from which transport in an equation
template at sources other than the literal $0$ is derivable without a
computation equation.  Then it derives every instance of \rulename{I0} and
of \rulename{IS}.
\end{lemma}

\begin{proof}
$c_{*}=0$ is a permitted instance of \rulename{IS}: the predecessor of the
counter and the seed are $1$, and the body is the constant $0$.  Write an
instance that is not permitted as $P(0)$, where the equation template $P(v)$
has $v$ at the offending positions:
\begin{align*}
 &\It{0}{v}{u}{h}=v, &
 &\It{\Suc n}{v}{u}{h}=h\subst{\It{n}{v}{u}{h}}{u},\\
 &\It{\Suc v}{a}{u}{h}=h\subst{\It{v}{a}{u}{h}}{u}, &
 &\It{\Suc v}{v}{u}{h}=h\subst{\It{v}{v}{u}{h}}{u}.
\end{align*}
The variable $v$ is fresh, and other occurrences of $0$ are left alone.  In
each case $P(c_{*})$ is a permitted instance.  Transport in the template $P$
along $c_{*}=0$, at the source $c_{*}\not\equiv 0$, gives $P(0)$.
\end{proof}

The hypothesis on the transport rule holds for \rulename{Tr},
$\rulename{Tr}^{-}$ (\cref{kern:lem:oseq}(iii)), \rulename{TrX},
$\rulename{Tr}_{2}$ and $\rulename{Tr}_{F}^{-}$.  Hence every result of this
chapter that uses $c=\It{0}{0}{u}{0}$ or another computation instance with a
literal $0$ holds for the restricted inventory as well: the instance is
first derived by the lemma.  The repair is one transport on two computation
instances, and its only source is $c_{*}$.

\begin{remark}[Sources]
\label{kern:rem:comprepair-sources}
The lemma is proved in \pkg{witness\_rule}, where $c_{*}$ serves also as the
second name for zero in \cref{kern:thm:witness}, and in
\pkg{source\_dispatch}, which checks that the repair does not use transport
at the source $0$ before \cref{kern:thm:dispatch} is available.
\pkg{reversal\_core} and \pkg{template\_free} use it in their certificate
producers.  The restricted inventory itself belongs to an earlier package,
\pkg{certified\_zero}.
\end{remark}

% ---------------------------------------------------------------------
\section{Elimination of the closure by abstraction of the goal}
\label{kern:sec:aqe}

This section draws on \pkg{abstraction\_and\_refinements}.

The rule \rulename{QE} is given a matrix $B(y)$ and a term $t$ and concludes
$B(t)$.  Its goal-directed form is given the conclusion and the term $t$
only, and computes the matrix by abstracting every occurrence of $t$ in the
conclusion.  Together with \rulename{BRep} of \cref{kern:def:B} and
$\rulename{DI}_{\mathrm{w}}$ of \cref{kern:def:W} this gives a calculus of
seven rules in which neither the transport rule nor the elimination of a
closure is given a template: the premises of each are computed from its
conclusion and one further term or equation.  The difficulty is completeness: an
abstraction cannot leave chosen occurrences of $t$ in place, and the
elimination at an arbitrary matrix has to be recovered without assuming the
transport in compound formulas that it is used to prove.

\subsection{Replacement in formulas and abstraction}
\label{kern:sec:abstr}

The operation $\operatorname{Rep}_{t}^{s}$ of \cref{kern:sec:brep} extends
to native formulas by the same traversal.  The bound variables of $\Q{}$ and
$\D{}{}$ are renamed away from $\FV(s)\cup\FV(t)$, as those of iterators; an
equation $u_1=u_2$ becomes
$\operatorname{Rep}_{t}^{s}(u_1)=\operatorname{Rep}_{t}^{s}(u_2)$; and
$\Q{y}B$ and $\D{x}{y}B$ become $\Q{y}\operatorname{Rep}_{t}^{s}(B)$ and
$\D{x}{y}\operatorname{Rep}_{t}^{s}(B)$.  Only term nodes are compared with
$t$, a node identical with $t$ is replaced by $s$ and not traversed further,
and inserted copies of $s$ are not traversed.  The operation is total, also
when $s$ contains $t$.  A bound variable of the formula never matches a free
variable of $t$.

\begin{lemma}[Abstraction]
\label{kern:lem:abstr}
Let $C$ be a term or a native formula and $s$, $t$ terms, and let the bound
variables of $C$ be renamed away from $\FV(s)\cup\FV(t)$.
\begin{enumerate}[label=\textup{(\roman*)}]
\item The occurrences of $t$ that the traversal of $C$ replaces are pairwise
 disjoint free occurrences, and they depend on $C$ and $t$ only, not on $s$.
 Replacing them by a fresh variable $v$ gives a template $H(v)$ with
 $C\equiv H(t)$ and $\operatorname{Rep}_{t}^{s}(C)\equiv H(s)$.
\item If $y$ is not free in $C$ or in $t$, then
 \[
  \operatorname{Rep}_{t}^{y}(C)\subst{t}{y}\equiv C,\qquad
  \operatorname{Rep}_{t}^{y}(C)\subst{0}{y}\equiv\operatorname{Rep}_{t}^{0}(C).
 \]
\end{enumerate}
\end{lemma}

\begin{proof}
(i) At each node the traversal first compares the node with $t$ and descends
into its children exactly when the comparison fails.  Which nodes are
compared, and with what result, is therefore determined by $C$ and $t$; the
term $s$ enters only as the inserted copy, which is not traversed.  The
traversal does not descend into a replaced occurrence, so two replaced
occurrences are disjoint.  A replaced occurrence is free in $C$, because the
bound variables of $C$ are distinct from the variables of $t$.  Put $v$ at
the replaced occurrences to obtain $H(v)$.  Then $H(t)$ is $C$, and $H(s)$
is the result of the traversal, in which no variable of $s$ is captured
because the bound variables of $H$ avoid $\FV(s)$.

(ii) Rename the bound variables of $C$ away from $\FV(t)\cup\{y\}$ and let
$H$ be the template of (i); by (i) it serves for $s\defeq y$ and for
$s\defeq 0$, since $\FV(0)$ is empty.  Then
$\operatorname{Rep}_{t}^{y}(C)\equiv H(y)$ and
$\operatorname{Rep}_{t}^{0}(C)\equiv H(0)$.  The variable $y$ is not free in
$C$, so its free occurrences in $H(y)$ are exactly the positions of $v$.
Hence $H(y)\subst{t}{y}\equiv H(t)\equiv C$, with no capture because the bound
variables of $H$ avoid $\FV(t)$, and $H(y)\subst{0}{y}\equiv H(0)$.
\end{proof}

For terms, (i) is the argument of \cref{kern:lem:bsound}.  The identities of
(ii) are facts about syntax and not derivations of equations.

\begin{lemma}[Exact marking in a formula]
\label{kern:lem:fmarker}
Let $C(v)$ be a term template or a formula template, $p$ a term, and let
the bound variables of $C$ be renamed away from $\FV(p)\cup\{v\}$.  Let $|C|$
be the number of nodes of $C$, the nodes of its terms and the bodies of its
iterators included, and put
\[
 n\defeq 1+|C|,\qquad K\defeq\It{0}{p}{w}{\Suc^{n}0}.
\]
Then the occurrences of $K$ in $C(K)$ are exactly the copies inserted for
$v$, and $\operatorname{Rep}_{K}^{a}\bigl(C(K)\bigr)\equiv C(a)$ for every
term $a$.  Moreover $K=p$ is an instance of \rulename{I0} and
$\FV(K)=\FV(p)$.
\end{lemma}

\begin{proof}
The body $\Suc^{n}0$ alone has $n+1$ nodes, so $|K|>n>|C|$.  Only term
nodes are compared with $K$.  Two term
occurrences in $C(K)$ are disjoint or nested.  A term occurrence that neither
contains nor lies inside an inserted copy is an occurrence of a subterm of
$C$ and has at most $|C|$ nodes; one that lies properly inside an inserted
copy has fewer nodes than $K$; one that properly contains an inserted copy
has more.  None of these is identical with $K$, also after the renaming of
bound variables, which preserves the number of nodes.  So the traversal
reaches every inserted copy, since no node above it matches, replaces it by
$a$ and stops there, and changes nothing else.  The result is $C$ with $a$ at
the positions of $v$, which is $C(a)$.  The last sentence is immediate from
the form of $K$.
\end{proof}

\Cref{kern:lem:marker} is the same argument for term templates, with the
seed $s$ and a larger $n$.  The proof allows several occurrences of $v$, no
occurrence, and occurrences below any number of iterator, $\Q{}$ and $\D{}{}$
binders.  The bound is on the template $C$, not on $C(K)$.

\subsection{The rule and the calculus}
\label{kern:sec:aqerule}

\begin{definition}[The calculus $\mathsf{A}_7$]
\label{kern:def:A7}
The rule of elimination by abstraction is
\[
 \infrule{\Gamma\proves\Q{y}\operatorname{Rep}_{t}^{y}(A)\qquad
          \Gamma\proves\operatorname{Rep}_{t}^{0}(A)}
         {\Gamma\proves A}\ \rulename{AQE},
\]
where $y$ is a variable not free in $A$ or in $t$.
% ed.: the package also asks y fresh for Gamma; y is bound in the first premise, so its name is immaterial.
$\mathsf{A}_7$ is $\mathsf{B}_7$ with \rulename{QE} replaced by
\rulename{AQE}; its seven rules are \rulename{Hyp}, \rulename{I0},
\rulename{IS}, \rulename{BRep}, \rulename{AQE}, \rulename{DE} and
$\rulename{DI}_{\mathrm{w}}$, with \rulename{BRep} at unrestricted sources,
the literal $0$ included.
\end{definition}

For a conclusion $A$ and a term $t$ both premises of \rulename{AQE} are
determined, up to the name of the bound variable $y$; an instance carries no
matrix.  The variable $y$ is bound, so the condition on it only ensures that
abstraction does not identify $y$ with a variable of $A$ or $t$.

\begin{lemma}[\rulename{AQE} in $\Nat$]
\label{kern:lem:aqesound}
Every instance of \rulename{AQE} is an instance of \rulename{QE}.
\end{lemma}

\begin{proof}
Let $B(y)\defeq\operatorname{Rep}_{t}^{y}(A)$.  By
\cref{kern:lem:abstr}(ii), $B(0)\equiv\operatorname{Rep}_{t}^{0}(A)$ and
$B(t)\equiv A$, so the instance is \rulename{QE} with the matrix $B$, the
premises $\Q{y}B(y)$ and $B(0)$ and the term $t$.
\end{proof}

The converse fails.  \rulename{QE} with the matrix $y=x$, the premises
$\Q{y}(y=x)$ and $0=x$ and the term $x$ concludes $x=x$; \rulename{AQE} with
the conclusion $x=x$ and the term $x$ asks for $\Q{y}(y=y)$ and $0=0$, and no
other term gives the premise $0=x$, since both sides of $x=x$ are the same
term and are replaced together.  An abstraction cannot leave a chosen
occurrence in place; \cref{kern:lem:fmarker} recovers a chosen matrix.

\begin{remark}[Substitution in $\mathsf{A}_7$]
\label{kern:rem:aqesubst}
Weakening, injective renaming and cut are admissible in $\mathsf{A}_7$ by
\cref{kern:lem:struct}: weakening and cut change no term or formula of an
instance of \rulename{AQE}, and an injective renaming commutes with
$\operatorname{Rep}$ and so maps an instance to an instance.  Substitution does not map instances to
instances.  \rulename{AQE} with the conclusion $x=z$, $z\not\equiv x$, and
the term $x$ has the premises $\Q{y}(y=z)$ and $0=z$.  The substitution of
$x$ for $z$ maps it to the inference from $\Q{y}(y=x)$ and $0=x$ to $x=x$,
which by the preceding paragraph is not an instance of \rulename{AQE}.  A
substitution can also turn a witness of $\rulename{DI}_{\mathrm{w}}$ into the
literal $0$ (\cref{kern:rem:wsubst}).  For this reason the proofs below use
no substitution into a derivation, and substitution is admissible in
$\mathsf{A}_7$ only through \cref{kern:prop:transfer}, after
\cref{kern:thm:aqe}.
\end{remark}

\subsection{Introductions before elimination}
\label{kern:sec:aqeintro}

\begin{lemma}[Rules of $\mathsf{A}_7$ derived without \rulename{AQE}]
\label{kern:lem:aqeintro}
The following rules are derivable in $\mathsf{A}_7$ over an arbitrary
context by derivations that contain no instance of \rulename{AQE}.
\begin{enumerate}[label=\textup{(\roman*)}]
\item Reflexivity, symmetry, transitivity, congruence in every term template,
 \rulename{Tr}, and transport in an equation template: from $s=t$ and
 $a(s)=b(s)$ to $a(t)=b(t)$.
\item \rulename{QI}.
\item \rulename{DI}, for every matrix and every witness.
\item Closure by cases: from $\Gamma,B(0)\proves B(\Suc z)$, with $z$ not
 free in $\Gamma$ or in $\Q{y}B(y)$, to $\Gamma\proves\Q{y}B(y)$, for every
 native matrix $B$.
\end{enumerate}
\end{lemma}

\begin{proof}
(i) \Cref{kern:lem:beq,kern:lem:bcong} derive the equality rules,
congruence and \rulename{Tr} from \rulename{I0} and \rulename{BRep} alone,
and transport in an equation template is two applications of \rulename{Tr}
and the equality rules (\cref{kern:lem:local}(iii)).

(ii) and (iii) The proof of \cref{kern:thm:witness}, from $\Nat$ to
$\mathsf{W}$, derives \rulename{QI} from $\rulename{DI}_{\mathrm{w}}$ in the
mode of progress, \rulename{DE} and \rulename{Hyp}, and the remaining case of
\rulename{DI}, a literal equation with the witness $0$, from
$\rulename{DI}_{\mathrm{w}}$ in the mode of cases at the witness
$c=\It{0}{0}{u}{0}$, the equation $c=0$ of \rulename{I0} and its symmetric
form, transport in an equation template and cut.  These are available by
(i) and \cref{kern:rem:aqesubst}, and the proof uses no elimination of a
closure.

(iv) Choose $x$ not free in $B$ or in $\Gamma$, and regard $B(y)$ as a
matrix in $x$ and $y$ in which $x$ is not used.  At the closed witness $c$,
$\rulename{DI}_{\mathrm{w}}$ is in the mode of cases, and its premise
$\Gamma,B(0)\proves B(\Suc z)$ is the given one, since $x$ does not occur;
$z$ is not free in $c$.  It concludes $\Gamma\proves\D{x}{y}B(y)$.
\rulename{DE} with a fresh $b$ opens $\Q{y}B(y)$, and its second premise
$\Gamma,\Q{y}B(y)\proves\Q{y}B(y)$ is \rulename{Hyp}.
% ed.: the package derives (iv) from the general case introduction; at the witness c no transport is needed.
\end{proof}

\subsection{Recovery of the ordinary elimination}
\label{kern:sec:aqerank}

Let $r(A)$ be the binder rank of \cref{base:rem:single-binder-size}, the
number of occurrences of $\Q{}$ and $\D{}{}$ in $A$.  Substitution of terms
and replacement $\operatorname{Rep}_{t}^{s}$ leave it unchanged.

\begin{lemma}[Transport and elimination by rank]
\label{kern:lem:aqerank}
For every $r\ge 0$ the following hold in $\mathsf{A}_7$, for every context
$\Gamma$.
\begin{itemize}
\item[$(\mathrm L_r)$] If $F(v)$ is a formula template with $r(F)\le r$,
 $\Gamma\proves s=t$ and $\Gamma\proves F(s)$, then $\Gamma\proves F(t)$.
\item[$(\mathrm E_r)$] If $B(y)$ is a matrix with $r(B)\le r$,
 $\Gamma\proves\Q{y}B(y)$ and $\Gamma\proves B(0)$, then
 $\Gamma\proves B(t)$ for every term $t$.
\end{itemize}
\end{lemma}

\begin{proof}
By induction on $r$, proving $(\mathrm L_r)$ from $(\mathrm L_{r-1})$ and
$(\mathrm E_{r-1})$, and then $(\mathrm E_r)$ from $(\mathrm L_r)$.  The
bound variables of $F$ and $B$ and the variable $v$ are chosen fresh for the
other data.

\emph{$(\mathrm L_r)$, $F$ a literal equation.}  This is transport in an
equation template, \cref{kern:lem:aqeintro}(i).

\emph{$(\mathrm L_r)$, $F(v)\equiv\Q{y}B(v,y)$.}  Then $r\ge 1$ and
$r(B)\le r-1$.  Let $z$ be a variable not free in $\Gamma$, $s$, $t$ or $B$,
and $\Gamma'\defeq\Gamma,B(t,0)$.  Over $\Gamma'$, by weakening: $s=t$ and
$\Q{y}B(s,y)$; symmetry gives $t=s$; \rulename{Hyp} gives $B(t,0)$;
$(\mathrm L_{r-1})$ with the template $B(v,0)$ and $t=s$ gives $B(s,0)$;
$(\mathrm E_{r-1})$ with the matrix $B(s,y)$ and the term $\Suc z$ gives
$B(s,\Suc z)$; and $(\mathrm L_{r-1})$ with the template $B(v,\Suc z)$ and
$s=t$ gives $B(t,\Suc z)$.  Hence $\Gamma,B(t,0)\proves B(t,\Suc z)$, and
closure by cases (\cref{kern:lem:aqeintro}(iv)) gives
$\Gamma\proves\Q{y}B(t,y)$, which is $F(t)$.

\emph{$(\mathrm L_r)$, $F(v)\equiv\D{x}{y}B(v,x,y)$.}  Then $r(B)\le r-1$.
Let $b$ and $z$ be distinct variables not free in $\Gamma$, $s$, $t$ or $B$,
and $\Delta\defeq\Gamma,\Q{y}B(s,b,y)$.  Over $\Delta,B(t,b,0)$ the steps of
the preceding case, with the templates $B(v,b,0)$ and $B(v,b,\Suc z)$ and the
matrix $B(s,b,y)$, give $B(t,b,\Suc z)$.  $\rulename{DI}_{\mathrm{w}}$ with
the matrix $B(t,x,y)$ and the witness $b$, which is not the literal $0$, is
in the mode of cases and concludes $\Delta\proves\D{x}{y}B(t,x,y)$; $z$ is not
free in $\Delta$, in $b$ or in the conclusion.  \rulename{DE} with the first
premise $\Gamma\proves\D{x}{y}B(s,x,y)$ and the eigenvariable $b$, which is
not free in $\Gamma$, in $\D{x}{y}B(s,x,y)$ or in $\D{x}{y}B(t,x,y)$,
concludes $\Gamma\proves\D{x}{y}B(t,x,y)$, which is $F(t)$.

These three cases exhaust the formulas of rank at most $r$.  Every
transport and elimination invoked in the last two acts on a matrix of rank
at most $r-1$.

\emph{$(\mathrm E_r)$.}  Let $r(B)\le r$, $\Gamma\proves\Q{y}B(y)$,
$\Gamma\proves B(0)$, and let $t$ be a term.  Let the bound variables of $B$
avoid $\FV(t)$, let $v$ be fresh and $K\defeq\It{0}{t}{w}{\Suc^{n}0}$ the
marker of \cref{kern:lem:fmarker} for the template $C(v)\defeq B(v)$ and the
term $p\defeq t$.  Let $y'$ be a variable not free in $B$ or in $t$.  By
\cref{kern:lem:fmarker},
\[
 \operatorname{Rep}_{K}^{y'}\bigl(B(K)\bigr)\equiv B(y'),\qquad
 \operatorname{Rep}_{K}^{0}\bigl(B(K)\bigr)\equiv B(0).
\]
Hence \rulename{AQE} with the conclusion $B(K)$ and the term $K$ has the
premises $\Q{y'}B(y')$, which is $\Q{y}B(y)$ up to the name of the bound
variable, and $B(0)$.  It concludes $\Gamma\proves B(K)$.  \rulename{I0}
gives $\Gamma\proves K=t$, and $(\mathrm L_r)$ with the template $B(v)$
gives $\Gamma\proves B(t)$.  For $r=0$ this uses $(\mathrm L_0)$ only.
\end{proof}

The order $(\mathrm L_0),(\mathrm E_0),(\mathrm L_1),(\mathrm E_1),\dots$ is
not circular: \cref{kern:lem:aqeintro} uses no elimination of a closure,
transport at rank $r$ uses elimination at rank $r-1$ only, and elimination at
rank $r$ uses \rulename{AQE} once, at the marker, and transport at rank $r$.
No transport of a compound formula enters the derivation of the
introductions, and no substitution into a derivation is used.

\begin{theorem}[Elimination by abstraction]
\label{kern:thm:aqe}
$\mathsf{A}_7$, $\mathsf{B}_7$ and $\Nat$ are equivalent.
\end{theorem}

\begin{proof}
$\mathsf{B}_7$ and $\Nat$ are equivalent by \cref{kern:thm:brep}.
\emph{From $\mathsf{A}_7$ to $\Nat$.}  \rulename{Hyp}, \rulename{I0},
\rulename{IS} and \rulename{DE} are common; an instance of \rulename{BRep} is
an instance of \rulename{Tr} (\cref{kern:lem:bsound}), an instance of
\rulename{AQE} one of \rulename{QE} (\cref{kern:lem:aqesound}), and an instance
of $\rulename{DI}_{\mathrm{w}}$ is derivable in $\Nat$ (\cref{kern:thm:witness}).
\emph{From $\Nat$ to $\mathsf{A}_7$.}  By induction on a derivation in
$\Nat$: \rulename{Hyp}, \rulename{I0}, \rulename{IS} and \rulename{DE} are
common; \rulename{Tr}, \rulename{QI} and \rulename{DI} are derivable by
\cref{kern:lem:aqeintro}; and an instance of \rulename{QE} with the matrix
$B$ is derivable by $(\mathrm E_{r(B)})$ of \cref{kern:lem:aqerank}.  Each
derived rule acts over the context of the inference it replaces and
introduces only fresh variables, so the conditions on the eigenvariables of
enclosing inferences are kept.  Both translations leave terms, formulas,
contexts and conclusions unchanged.
\end{proof}

By \cref{kern:prop:transfer}, $\mathsf{A}_7$, $\Teq$ and $\iPRA$ are
bi-interpretable with whole domains, actual equality and identity numerical
composites, and $\proves_{\mathsf{A}_7}0=\Suc 0$ holds exactly when $\iPRA$
is inconsistent.  The translations between $\mathsf{A}_7$ and $\Nat$ are the
identity on terms and formulas, so the bridge of \cref{base:thm:bridge} is
used unchanged.

The proof needs \rulename{BRep} at the source $0$.  The restriction
$\rulename{BRep}^{-}$ of \cref{kern:thm:bminus} is enough as well.  Its
bootstrap (\cref{kern:lem:bmeq,kern:lem:bzero}) uses \rulename{QE}, which in
$\mathsf{A}_7$ is recovered only after transport; but it uses \rulename{QE}
only twice, at instances whose abstraction is the matrix itself.

\begin{theorem}[Elimination by abstraction with a restricted source]
\label{kern:thm:aqeminus}
$\mathsf{A}_7$ with \rulename{BRep} replaced by $\rulename{BRep}^{-}$ is
equivalent to $\Nat$.
\end{theorem}

\begin{proof}
Every instance of $\rulename{BRep}^{-}$ is an instance of \rulename{BRep}, so
the restricted calculus derives only sequents of $\mathsf{A}_7$, hence of
$\Nat$ by \cref{kern:thm:aqe}.  Conversely, the proofs of
\cref{kern:lem:bmeq,kern:lem:bzero} go through with each use of
\rulename{QE} replaced by \rulename{AQE}.

In \cref{kern:lem:bmeq}(iii), \rulename{QE} at the term $p$ concludes
$E(p)\equiv\bigl(\It{p}{p}{v}{p}=p\bigr)$ from $\Q{y}E(y)$ and $E(0)$.  The
left side of $E(p)$ is larger than $p$, so the traversal does not match at its
root, and its counter, seed and body are copies of $p$.  Hence the abstraction
of $p$ in $E(p)$ is $E(y)$, its replacement by $0$ is $E(0)$, and
\rulename{AQE} at the term $p$ draws the same conclusion from the same
premises.

In \cref{kern:lem:bzero}, \rulename{QE} at the term $\Suc 0$ concludes
$E(b,\Suc 0)$, that is $C(\Suc 0,0,b)=C(\Suc 0,0,0)$, from $\Q{y}E(b,y)$ and
$E(b,0)$.  Its only occurrences of $\Suc 0$ are the two counters: the seeds
$0$, the variable $b$ and the bodies are not $\Suc 0$.  So again the
abstraction is the matrix $E(b,y)$, the replacement by $0$ is $E(b,0)$, and
\rulename{AQE} at the term $\Suc 0$ applies.

The introductions in the two proofs are $\rulename{DI}_{\mathrm{w}}$ at the
witness $0$ in the mode of progress, which is \rulename{DI} at the witness $0$
of $\mathsf{A}_7$, and \rulename{QI} derived from it with \rulename{DE} and
\rulename{Hyp}; in both the assumption is unused.  The remaining steps use
\rulename{I0}, \rulename{IS}, $\rulename{BRep}^{-}$, \rulename{DE}, weakening
and injective renaming, which are admissible by induction on derivations as
for $\mathsf{A}_7$; renaming keeps the side condition, since a variable is
renamed to a variable.  So congruence and \rulename{Tr} are derived as in the
proof of \cref{kern:thm:bminus}, hence transport in an equation template, and
from there the proof of \cref{kern:thm:aqe} recovers \rulename{QE} by
induction on the number of binders (\cref{kern:lem:aqerank}) and the
introductions of $\Nat$.
\end{proof}

The Lean module \Lean{NativeNonzeroSources} formalizes
\cref{kern:thm:aqeminus} and \cref{kern:thm:bminus} for all three restricted
calculi (\Lean{amprf\_iff\_nprf}, \Lean{bmprf\_iff\_nprf},
\Lean{bm7prf\_iff\_nprf}).  The bootstrap of
\cref{kern:lem:bmeq,kern:lem:bzero} is proved once, generically over the class
\Lean{BMKernel}, whose only elimination of a closure is \rulename{AQE} and
whose only introduction is the one at the witness $0$ with the assumption
unused.  In $\mathsf{B}^{-}$ and $\mathsf{B}^{-}_{7}$, \rulename{AQE} is
\rulename{QE} followed by the identities of \cref{kern:lem:abstr}(ii)
(\Lean{NFm.aqE\_of\_qE}).  For the restricted $\mathsf{A}_7$ the recovery of
\rulename{QE} is that of \Lean{NativeAbstraction}, stated there over the class
\Lean{AKernel} of the rules of $\mathsf{A}_7$.  In the variable lemma of
\cref{kern:lem:bzero} the formal proof needs no freshness of $b$: if $b$ is a
variable, $\operatorname{Rep}_{0}^{b}(h(0))\equiv\operatorname{Rep}_{0}^{b}(h(b))$
holds for every template $h$ (\Lean{Tm.rpl\_zero\_subst}).

\begin{remark}[Sources]
\label{kern:rem:aqe-sources}
The section is the abstraction part of \pkg{abstraction\_and\_refinements},
whose calculus $\mathcal{A}_7$ is $\mathsf{A}_7$; the package takes
\rulename{BRep}, $\rulename{DI}_{\mathrm{w}}$ and the bridge from this
chapter and from \cref{base:thm:bridge}, and presents its proof twice, as a
section and as a bootstrap note.  The package proves the equality bootstrap
again from \rulename{I0} and \rulename{BRep}; the article cites
\cref{kern:lem:beq,kern:lem:bcong}.  Its marker has the seed $t$ and
$N\defeq|C|+1$, as in \cref{kern:lem:fmarker}.  It derives closure by cases
from the general introduction by cases; the article introduces at the
witness $c$ directly.  The Lean module \Lean{NativeAbstraction} formalizes
$\mathsf{A}_7$ and \cref{kern:thm:aqe} for it (\Lean{APrf},
\Lean{aprf\_iff\_nprf}), along the proof above: replacement and abstraction
are one operation that puts a term at the occurrences of $t$ and renames the
other variables, so the identities of \cref{kern:lem:abstr}(ii) are instances
of its commutation with substitution; weakening and injective renaming of
derivations are proved by induction on derivations; and
\cref{kern:lem:aqerank} is one induction on the number of binders
(\Lean{APrf.trEl}).  The
tests of the package are finite: 804 checks of abstraction and of open
inferences, 521 checks of the recovery by the marker, 3 further identities
of capture avoidance, 28 runs of the driver on closed certificates, and
1521 independent symbolic checks on 300 generated syntax trees, repeated
with assertions disabled.  The 45-node predecessor of the package is
\cref{base:prop:pred45}; its refinements of checkers for inherited
presentations are reported in the chapter on certificates.
\end{remark}

% ---------------------------------------------------------------------
\section{Induction and uniqueness of iteration}
\label{kern:sec:uniq}

This section draws on \pkg{scalar\_uniqueness}, \pkg{orbit\_closure},
\pkg{crossed\_uniqueness} and the rigidity section of \pkg{reductions}.

Induction for literal equations can be replaced by a rule that has no
induction hypothesis in its premise: a term that satisfies the recurrence of
an iterator is that iterator.  The three packages state the rule in three
forms.  \pkg{scalar\_uniqueness} concludes the equation at a supplied target;
this is the rule \rulename{UR} below, and
\crefrange{kern:sec:ur}{kern:sec:urreflect} prove that it replaces equation
induction.  \pkg{orbit\_closure} concludes the closure of the equation, with
no target, and isolates what the restriction of the rule to the body
$\Suc u$ already gives.  \pkg{crossed\_uniqueness} adds a premise for the
seed and crosses the conclusion; on fully normalized terms that rule also
replaces transport and the two computation equations.  The second and third
forms are treated in \cref{kern:sec:urforms}, which reduces each to
\rulename{UR}; the calculus of \pkg{crossed\_uniqueness} is the one calculus
of the chapter that is stated on normal forms.

\subsection{The rule and the two theorems}
\label{kern:sec:ur}

Let $f(z)$ and $h(u)$ be term templates, and write $h(f(z))$ for
$h\subst{f(z)}{u}$.  The rule of \emph{uniqueness of iteration} is
\[
 \infrule{\Gamma\proves f(\Suc z)=h(f(z))}
         {\Gamma\proves f(t)=\It{t}{f(0)}{u}{h(u)}}\ \rulename{UR},
\]
where $z$ is not free in $\Gamma$, in $t$ or in $h$.  The body $h$ may have
other free variables, which are parameters.  The premise is derived in
$\Gamma$ itself, with no added assumption; the seed $f(0)$ is read off the
template, so there is no premise for the base; and the body receives one
state and no counter.  The rule is a scheme over terms: $z$ is a numerical
variable and nothing ranges over functions.

\begin{definition}[The calculus $\mathsf{U}$ and the theory $T_{\mathrm{UR}}$]
\label{kern:def:U}
$\mathsf{U}$ is $\Nat$ with \rulename{QI} replaced by \rulename{UR}; its
eight rules are \rulename{Hyp}, \rulename{I0}, \rulename{IS}, \rulename{Tr},
\rulename{QE}, \rulename{DE}, \rulename{DI} and \rulename{UR}.  The ordinary
theory $T_{\mathrm{UR}}$ is $\Teq$ (\cref{base:sec:teq}) with its rule of
induction for literal equations replaced by \rulename{UR}, read for
contexts of ordinary formulas: intuitionistic logic with equality, the two
computation equations, coverage $x=0\vee\exists y\,x=\Suc y$, separation
$0=\Suc 0\to\bot$, and \rulename{UR}.
\end{definition}

In $\mathsf{U}$ the introduction \rulename{DI} treats atomic and compound
matrices alike, and no rule looks at a witness; substitution maps instances
of \rulename{UR} to instances after a renaming of $z$
(\cref{kern:lem:struct}).

\begin{definition}[Successor uniqueness]
\label{kern:def:SU}
\rulename{SU} is the rule \rulename{UR} restricted to the body
$h(u)\defeq\Suc u$: from $\Gamma\proves f(\Suc z)=\Suc f(z)$ infer
$\Gamma\proves f(t)=\It{t}{f(0)}{u}{\Suc u}$.  The theory
$T_{\mathrm{SU}}$ is $T_{\mathrm{UR}}$ with \rulename{UR} replaced by
\rulename{SU}.
\end{definition}

An instance of \rulename{SU} supplies no body.  Every derivation of
$T_{\mathrm{SU}}$ is a derivation of $T_{\mathrm{UR}}$, and the lemmas below
say which of them use \rulename{SU} only.

\begin{proposition}[Uniqueness from induction]
\label{kern:prop:uniq}
\rulename{UR} is a derived rule of $\Nat$ and of $\Teq$.
\end{proposition}

\begin{proof}
Let $g(y)\defeq\It{y}{f(0)}{u}{h(u)}$ and apply equation induction to
$f(y)=g(y)$.  Base: $g(0)=f(0)$ by \rulename{I0}, and symmetry.  Step:
assume $f(z)=g(z)$.  The premise gives $f(\Suc z)=h(f(z))$, congruence gives
$h(f(z))=h(g(z))$, and \rulename{IS} gives $g(\Suc z)=h(g(z))$; hence
$f(\Suc z)=g(\Suc z)$.
\end{proof}

In Lean this derived rule is \Lean{NPrf.uniq}.

\begin{theorem}[Uniqueness of iteration in place of equation induction]
\label{kern:thm:ur}
$T_{\mathrm{UR}}$ and $\Teq$ derive the same sequents of ordinary formulas.
Consequently $T_{\mathrm{UR}}$ is bi-interpretable with $\iPRA$ with whole
domains, actual equality and identity numerical composites.
\end{theorem}

\begin{theorem}[The native form]
\label{kern:thm:urnative}
$\mathsf{U}$ and $\Nat$ are equivalent.
\end{theorem}

One half of each theorem is \cref{kern:prop:uniq}.  The other half, the
derivation of induction for literal equations from \rulename{UR}, occupies
\crefrange{kern:sec:urfirst}{kern:sec:urreflect}.  Its order avoids a
circle: uniqueness gives induction for equations of the form $q(n)=0$ and,
separately, a predecessor; truncated subtraction and an explicit accumulator
give a term that reflects equality; and only then is induction for arbitrary
equations recovered.  Pairing and the compiler for primitive recursion
(\cref{base:thm:primrec}) are not used; they become available afterwards.
All derivations in
\crefrange{kern:sec:urfirst}{kern:sec:urreflect} are in $T_{\mathrm{UR}}$,
over an arbitrary fixed context $\Gamma$ of ordinary formulas, and variables
that are generalized or used as eigenvariables are not free in $\Gamma$.

\subsection{First consequences of uniqueness}
\label{kern:sec:urfirst}

\begin{lemma}[Common update, constancy, identity]
\label{kern:lem:cu}
\begin{enumerate}[label=\textup{(\roman*)}]
\item If $f(\Suc z)=h(f(z))$, $g(\Suc z)=h(g(z))$ and $f(0)=g(0)$, then
 $f(t)=g(t)$.
\item If $f(\Suc z)=f(z)$, then $f(t)=f(0)$.
\item $\It{n}{0}{u}{\Suc u}=n$.
\item If $h(c)=c$, then $\It{n}{c}{u}{h(u)}=c$.
\end{enumerate}
\end{lemma}

\begin{proof}
(i) \rulename{UR} for $f$ and for $g$ gives two iterators that differ only
in the equal seeds.  (ii) is (i) for $f$ and the constant template $f(0)$,
with $h(u)\defeq u$.  (iii) is \rulename{UR} for $f(z)\defeq z$ and
$h(u)\defeq\Suc u$, whose premise is reflexivity.  (iv) is \rulename{UR}
for the constant template $c$.
\end{proof}

Part (iii), and part (i) for the body $\Suc u$, are derivable in
$T_{\mathrm{SU}}$.  Part (ii) is uniqueness for the body $u$, and it is not
an instance of \rulename{SU}.

Put $a+b\defeq\It{b}{a}{u}{\Suc u}$, $2a\defeq a+a$ and $3a\defeq a+a+a$.
The computation equations give $a+0=a$ and $a+\Suc b=\Suc(a+b)$, so
$a+1=\Suc a$.

\begin{lemma}[Addition and iteration]
\label{kern:lem:uradd}
Addition is associative and commutative with the neutral element $0$, and
$\Suc a+n=\Suc(a+n)$.  For every body $h$,
\[
 \It{m+n}{a}{u}{h}=\It{n}{\It{m}{a}{u}{h}}{u}{h},\qquad
 \It{\Suc n}{a}{u}{h}=\It{n}{h(a)}{u}{h}.
\]
\end{lemma}

\begin{proof}
Each statement is \cref{kern:lem:cu}(i) or \rulename{UR} for functions of
$n$.  $0+n=n$ is \cref{kern:lem:cu}(iii).  $\Suc(a+n)$ has the seed
$\Suc a$ and the update $\Suc$, hence equals $\Suc a+n$.  The term $n+a$
has the seed $0+a=a$ and satisfies $\Suc n+a=\Suc(n+a)$, so it has the
update $\Suc$ and equals $a+n$.  The term $a+(b+n)$ has the seed $a+b$ and
the update $\Suc$, hence equals $(a+b)+n$.  The term $\It{m+n}{a}{u}{h}$
has the seed $\It{m}{a}{u}{h}$ and the update $h$.  The term
$\It{\Suc n}{a}{u}{h}$ has the seed $\It{1}{a}{u}{h}=h(a)$ and the update
$h$.
\end{proof}

The laws of addition in this lemma compare terms with the update $\Suc$
only, so they are derivable in $T_{\mathrm{SU}}$; the two equations for an
arbitrary body $h$ use \rulename{UR} for $h$.

\subsection{Retractions and induction for zero equations}
\label{kern:sec:urretract}

\begin{proposition}[A right inverse from successor uniqueness]
\label{kern:prop:rinv}
Let $e(x)$ be a term template such that
\[
 \Gamma\proves e(0)=0,\qquad \Gamma\proves e(\Suc e(x))=\Suc e(x)
\]
for a variable $x$ not free in $\Gamma$, and put
\[
 r(n)\defeq\It{n}{0}{u}{\Suc e(u)} .
\]
Then for every term $t$
\[
 \Gamma\proves e(r(t))=t,\qquad
 \Gamma\proves r(\Suc t)=\Suc t,\qquad
 \Gamma\proves e(\Suc t)=\Suc t .
\]
The derivation uses the two computation equations, equality reasoning and
two instances of \rulename{SU}; it uses neither coverage, nor separation,
nor idempotence of $e$.
\end{proposition}

\begin{proof}
Let $z$ be a variable not free in $\Gamma$, in $t$ or in $e$, and put
$f(n)\defeq e(r(n))$.  By \rulename{I0}, $r(0)=0$, so $f(0)=e(0)=0$.  By
\rulename{IS}, $r(\Suc z)=\Suc e(r(z))$, and the second hypothesis at the
argument $r(z)$, obtained by substituting $r(z)$ for $x$, gives
\[
 f(\Suc z)=e(r(\Suc z))=e\bigl(\Suc e(r(z))\bigr)=\Suc e(r(z))=\Suc f(z).
\]
\rulename{SU} for the template $f$ gives
$f(t)=\It{t}{f(0)}{u}{\Suc u}$; congruence with $f(0)=0$ and
\cref{kern:lem:cu}(iii), the second instance of \rulename{SU}, give
$f(t)=t$, that is, $e(r(t))=t$.  Then \rulename{IS} gives
$r(\Suc t)=\Suc e(r(t))=\Suc t$.  Finally, the second hypothesis at the
argument $r(t)$ is $e(\Suc e(r(t)))=\Suc e(r(t))$, and replacing $e(r(t))$
by $t$ on both sides gives $e(\Suc t)=\Suc t$.
\end{proof}

The first conclusion says that the term $r$ is a right inverse of $e$: it
is a term witness for the surjectivity of $e$, which is more than closure of
the image of $e$ under successor.  The trajectory $r$ applies $\Suc$ after
$e$, and \rulename{SU} is applied to the projected orbit $e\circ r$; this is
what makes the computation of $f(\Suc z)$ immediate.

\begin{lemma}[Rigidity of retractions]
\label{kern:lem:rigid}
Let $e(x)$ be a term template such that
\[
 e(0)=0,\qquad e(e(x))=e(x),\qquad e(\Suc e(x))=\Suc e(x)
\]
are derivable for a variable $x$ not free in $\Gamma$.  Then $e(t)=t$ for
every term $t$.  The derivation is in $T_{\mathrm{SU}}$ and uses neither
coverage nor separation.
\end{lemma}

\begin{proof}
By \cref{kern:prop:rinv}, $e(r(t))=t$.  Idempotence at the argument $r(t)$
gives
\[
 e(t)=e\bigl(e(r(t))\bigr)=e(r(t))=t . \qedhere
\]
\end{proof}

\begin{corollary}[Rigidity without idempotence]
\label{kern:cor:rigidcov}
Let $e(x)$ be a term template such that $e(0)=0$ and
$e(\Suc e(x))=\Suc e(x)$ are derivable for a variable $x$ not free in
$\Gamma$.  Then coverage and \rulename{SU} derive $e(t)=t$ for every term
$t$, with the computation equations and equality reasoning.  Neither
separation nor idempotence is used.
\end{corollary}

\begin{proof}
Coverage gives $t=0\vee\exists y\,t=\Suc y$.  If $t=0$, then
$e(t)=e(0)=0=t$.  If $t=\Suc y$ for a variable $y$ not free in $\Gamma$ or
$t$, then \cref{kern:prop:rinv} gives $e(\Suc y)=\Suc y$, hence $e(t)=t$.
Elimination of the disjunction and of the existential quantifier concludes
$e(t)=t$.
\end{proof}

Of the three hypotheses of \cref{kern:lem:rigid}, therefore, coverage may
replace idempotence, and idempotence makes coverage unnecessary; the two
remaining hypotheses alone already give $e(\Suc t)=\Suc t$ for every $t$.

\begin{lemma}[Induction from a numerical split]
\label{kern:lem:split}
Let $A(n)$ be an ordinary formula with $\Gamma\proves A(0)$ and
$\Gamma,A(z)\proves A(\Suc z)$.  Suppose that
$\Gamma\proves A(n)\vee\neg A(n)$ and that a term template $e(n)$ satisfies
$\Gamma,A(n)\proves e(n)=n$ and $\Gamma,\neg A(n)\proves e(n)=0$, with $n$
not free in $\Gamma$.  Then $\Gamma\proves A(t)$ for every term $t$.
\end{lemma}

\begin{proof}
By cases, $A(e(n))$: if $A(n)$, then $e(n)=n$; otherwise $e(n)=0$ and
$A(0)$ holds.  The first property of $e$ at the argument $e(n)$ gives
$e(e(n))=e(n)$.  The step at $e(n)$ gives $A(\Suc e(n))$, hence
$e(\Suc e(n))=\Suc e(n)$.  From $A(0)$, $e(0)=0$.
\Cref{kern:lem:rigid} gives $e(n)=n$, and $A(e(n))$ becomes $A(n)$.
\end{proof}

The lemma is a statement about a split that has been justified; it gives no
induction for formulas that are not decided by a term.

\begin{lemma}[Induction for zero equations]
\label{kern:lem:zi}
Let $q(n)$ be a term template.  If $\Gamma\proves q(0)=0$ and
$\Gamma,q(z)=0\proves q(\Suc z)=0$ with $z$ fresh, then
$\Gamma\proves q(t)=0$.
\end{lemma}

\begin{proof}
Separation gives $\Suc v\neq 0$: from $\Suc v=0$, the template
$\sel(w,0,1)$ gives $1=0$.  With coverage, $q(n)=0$ is decided.  The
template $e(n)\defeq\sel(q(n),n,0)$ equals $n$ if $q(n)=0$ and equals $0$
if $q(n)$ is a successor.  \Cref{kern:lem:split} applies.
\end{proof}

\Cref{kern:lem:split,kern:lem:zi} use uniqueness only through
\cref{kern:lem:rigid}.  Hence:

\begin{theorem}[Zero equations from successor uniqueness]
\label{kern:thm:su}
$T_{\mathrm{SU}}$ derives induction for zero equations, in the form of
\cref{kern:lem:zi}, over arbitrary contexts of ordinary formulas.
\end{theorem}

\begin{remark}[Sources]
\label{kern:rem:scopekernel}
\Cref{kern:lem:rigid,kern:lem:split,kern:lem:zi} are reproduced in
\pkg{scalar\_uniqueness} from an earlier package, \pkg{scope\_kernel}, which
leaves the predecessor and the reflection of equality as an open bootstrap.
The next three subsections close it.  The three packages prove
\cref{kern:lem:rigid} in three ways.  In \pkg{scope\_kernel} and
\pkg{scalar\_uniqueness} the trajectory is $r'(n)\defeq\It{n}{0}{u}{h(u)}$
with the unprojected body $h(u)\defeq e(\Suc e(u))$, and the first step,
$e(r'(n))=r'(n)$, is an instance of \rulename{UR} for the body $h$: by
idempotence $e(h(u))=h(u)$ and $h(e(u))=h(u)$, so the template $e(r'(n))$
has the seed $0$ and satisfies $e(r'(\Suc n))=h(r'(n))=h(e(r'(n)))$.  The
third hypothesis then gives $r'(\Suc n)=\Suc r'(n)$, \rulename{SU} gives
$r'(n)=n$, and $e(n)=e(r'(n))=r'(n)=n$.  The same proof is
\cref{obs:lem:retraction}.  \pkg{orbit\_closure} keeps $r'$ and replaces the
instance of \rulename{UR} by coverage on $n$: $r'(0)=0=e(0)$, and
$r'(\Suc m)=e(\Suc e(r'(m)))$ is a value of $e$, so idempotence gives
$e(r'(\Suc m))=r'(\Suc m)$; this places the derivation in $T_{\mathrm{SU}}$,
and \cref{kern:thm:su} is the theorem of that package.  The proof given
above, through \cref{kern:prop:rinv} with the body $\Suc e(u)$, is from
\pkg{reductions}; it needs neither the instance of \rulename{UR} nor
coverage, and \cref{kern:cor:rigidcov} is the other trade-off that it
permits.  In \cref{kern:lem:split,kern:lem:zi,kern:thm:su} coverage and
separation are now used only in \cref{kern:lem:zi}, to decide $q(n)=0$ and
so to justify the split $e(n)=\sel(q(n),n,0)$; the right inverse does not
remove them from that step.
\end{remark}

\begin{remark}[Successor rigidity in the clause chapter]
\label{kern:rem:rinv-claus}
\Cref{claus:thm:rigidity} derives induction for zero equations from the rule
$\mathrm{SR}_0$ (from $f(0)=0$ and $f(\Suc z)=\Suc f(z)$ infer $f(t)=t$),
which \rulename{SU} derives, with the trajectory $\It{n}{0}{u}{e(\Suc u)}$
and with coverage.  Its hypotheses on a formula $E$ and a term $e$,
$\Gamma\proves E(e(x))$ and $\Gamma,E(x)\proves e(x)=x$ with $x$ fresh,
together with the base $E(0)$ and the step $E(z)\proves E(\Suc z)$, imply the
three hypotheses of \cref{kern:lem:rigid}: substituting $0$ for $x$ in the
second hypothesis and using the base gives $e(0)=0$; substituting $e(x)$ for
$x$ and using the first gives $e(e(x))=e(x)$; and the step at $e(x)$ gives
$E(\Suc e(x))$, whence $e(\Suc e(x))=\Suc e(x)$.  In the proof of
\cref{kern:prop:rinv} the two instances of \rulename{SU} combine into one
instance of $\mathrm{SR}_0$ for the template $e(r(n))$, whose premises
$e(r(0))=0$ and $e(r(\Suc z))=\Suc e(r(z))$ are derived there.  Hence the
general statement of \cref{claus:thm:rigidity}, $\Gamma\proves E(t)$, holds
over the computation equations and $\mathrm{SR}_0$ without coverage:
\cref{kern:lem:rigid} gives $e(t)=t$, and $E(e(t))$ transports to $E(t)$.
Coverage remains in the derivation of $\mathrm{ZI}$ there, where it proves
$E(e(x))$ for $e(x)=\sel(q(x),x,0)$.
% ed.: this comparison combines kern:prop:rinv with the hypotheses of
% claus:thm:rigidity; neither package states it.
\end{remark}

\begin{remark}[Coverage is not derivable from the other principles]
\label{kern:rem:rinv-product}
\Cref{kern:lem:rigid} removes coverage from the step of rigidity only, and
coverage cannot be removed from the whole scalar presentation on the
strength of \rulename{SU}.  The product $\mathbb N\times\mathbb N$ of two
copies of the standard model, with $0=(0,0)$ and every operation
coordinatewise, satisfies the computation equations, separation and
induction for literal equations over arbitrary contexts, and refutes
coverage at $(0,1)$ (\cref{kern:rem:limits-sources}).  The derivation of
\rulename{UR} in \cref{kern:prop:uniq} uses equation induction, the
computation equations and congruence only, so every instance of
\rulename{UR}, and in particular of \rulename{SU}, is valid in the product.
Hence intuitionistic logic with equality, the computation equations,
separation and \rulename{SU} do not derive coverage.  \pkg{reductions}
checks the validity of \rulename{SU} in the product directly: its premise
gives the successor recurrence in each coordinate, and external induction in
each coordinate gives its conclusion.  The model has the standard domain
and the fixed language; it says nothing about an interpretation on a changed
domain (\cref{kern:q:coverdomain}).
\end{remark}

\subsection{A predecessor from uniqueness}
\label{kern:sec:urpred}

The construction follows the predecessor of \cref{base:thm:pred} (parity,
three-cycle, ramp, comparator, saturating counter), with a ramp modified so
that every induction of the original becomes an instance of
\cref{kern:lem:cu}.  It uses coverage and does not use separation.

Coverage and the selector equations give, by cases on $e$,
\begin{equation}
 \sel(e,\Suc a,\Suc b)=\Suc\,\sel(e,a,b). \label{kern:eq:selS}
\end{equation}
Recall $Z(c)=\sel(c,1,0)$ and $N(n)=\It{n}{0}{u}{Z(u)}$, and define
\[
 j(u)\defeq Z(u)+2N(u),\qquad M(n)\defeq\It{n}{0}{u}{j(u)},\qquad
 g(u)\defeq u+N(M(u))+1,\qquad G(n)\defeq\It{n}{0}{u}{g(u)}.
\]

\begin{lemma}[Two clocks]
\label{kern:lem:clocks}
\begin{enumerate}[label=\textup{(\roman*)}]
\item $N(n)=0\vee N(n)=1$; $Z(Z(N(n)))=N(n)$; $N(n)+Z(N(n))=1$;
 $N(N(n))=N(n)$; $N(n+2)=N(n)$.
\item $N(2n)=0$, $N(2n+1)=1$ and $N(3n)=N(n)$.
\item $M(n)=0\vee M(n)=1\vee M(n)=2$; $j(0)=1$, $j(1)=2$, $j(2)=0$;
 $M(n+3)=M(n)$.
\item $M(3n)=0$ and $M(3n+1)=1$.
\end{enumerate}
\end{lemma}

\begin{proof}
(i) By coverage on $n$: $N(0)=0$, and $N(\Suc m)=Z(N(m))$ is a value of
$Z$, which is $1$ at $0$ and $0$ at a successor.  The next three equations
hold in the two cases $N(n)=0$ and $N(n)=1$, by $Z(0)=1$, $Z(1)=0$,
$N(0)=0$ and $N(1)=Z(N(0))=1$.  Two computations give
$N(n+2)=Z(Z(N(n)))$.

(ii) The template $N(2n)$ satisfies $N(2(\Suc n))=N(2n+2)=N(2n)$, so
\cref{kern:lem:cu}(ii) gives $N(2n)=N(0)=0$; then $N(2n+1)=Z(N(2n))=1$.
The template $N(3n)$ has the seed $0$ and satisfies
$N(3(\Suc n))=N(3n+3)=Z(N(3n+2))=Z(N(3n))$, so \rulename{UR} gives
$N(3n)=\It{n}{0}{u}{Z(u)}=N(n)$.

(iii) $j(0)=Z(0)+2N(0)=1$.  At a successor $u$, $Z(u)=0$ and $2N(u)$ is $0$
or $2$.  So $j$ takes values in $\{0,1,2\}$, and so does $M$, by coverage
on $n$.  $j(1)=0+2N(1)=2$ and $j(2)=0+2N(2)=0$, since $N(2)=Z(N(1))=0$.
Three computations give $M(n+3)=j(j(j(M(n))))$, and $j$ permutes $0,1,2$
cyclically.

(iv) $M(3(\Suc n))=M(3n+3)=M(3n)$, so \cref{kern:lem:cu}(ii) gives
$M(3n)=M(0)=0$, and $M(3n+1)=j(M(3n))=1$.
\end{proof}

\begin{lemma}[The ramp]
\label{kern:lem:ramp}
\begin{enumerate}[label=\textup{(\roman*)}]
\item $M(g(u))=Z(M(u))$.
\item $M(G(n))=N(n)$.
\item $G(\Suc n)=G(n)+N(n)+1$, $G(n+2)=G(n)+3$ and
 $G(n+3)=G(n)+N(n)+4$.
\item $G(2n)=3n$ and $G(3n)=3G(n)+N(n)$.
\end{enumerate}
\end{lemma}

\begin{proof}
(i) By the three cases for $M(u)$.  If $M(u)=0$, then $g(u)=\Suc u$ and
$M(\Suc u)=j(0)=1=Z(0)$.  If $M(u)=1$, then $g(u)=\Suc\Suc u$ and
$M(\Suc\Suc u)=j(j(1))=0=Z(1)$.  If $M(u)=2$, then $N(2)=0$,
$g(u)=\Suc u$ and $M(\Suc u)=j(2)=0=Z(2)$.

(ii) The template $M(G(n))$ has the seed $M(0)=0$ and, by (i), satisfies
$M(G(\Suc n))=M(g(G(n)))=Z(M(G(n)))$.  \rulename{UR} gives
$M(G(n))=\It{n}{0}{u}{Z(u)}=N(n)$.

(iii) $G(\Suc n)=g(G(n))=G(n)+N(M(G(n)))+1=G(n)+N(N(n))+1$, and
$N(N(n))=N(n)$.  Two steps give $G(n+2)=G(n)+N(n)+1+Z(N(n))+1=G(n)+3$,
by $N(n)+Z(N(n))=1$.  A third step gives
$G(n+3)=G(n)+3+N(n+2)+1=G(n)+N(n)+4$.

(iv) $G(2n)$ and $3n$ have the seed $0$ and the update $v\mapsto v+3$, by
(iii); \cref{kern:lem:cu}(i) gives $G(2n)=3n$.  For the second equation put
\[
 f(n)\defeq G(3n),\qquad k(n)\defeq 3G(n)+N(n),\qquad
 \kappa(v)\defeq v+M(v)+4 .
\]
By (ii) and \cref{kern:lem:clocks}(ii), $M(f(n))=N(3n)=N(n)$.  By
\cref{kern:lem:clocks}(iv) and the two cases for $N(n)$, $M(k(n))=N(n)$.
By (iii), $f(\Suc n)=G(3n+3)=G(3n)+N(3n)+4=\kappa(f(n))$.  By (iii) and
\cref{kern:lem:clocks}(i),
\[
 k(\Suc n)=3G(n)+3N(n)+3+Z(N(n))=3G(n)+2N(n)+4=\kappa(k(n)).
\]
Both templates have the seed $0$, and \cref{kern:lem:cu}(i) gives
$f(n)=k(n)$.
\end{proof}

The update $g$ differs from the update $u+M(u)+1$ of the ramp in
\cref{base:thm:pred} away from the orbit of $0$; the difference is what
makes (i) hold at every argument, so that (ii) is an instance of
\rulename{UR} with no induction hypothesis.

Now define the comparison step, the comparator, the saturating step and the
predecessor:
\begin{gather*}
 \chi(u)\defeq\sel(N(u),G(u),u),\qquad
 \varepsilon(x)\defeq\It{x}{1}{u}{2u},\qquad
 \lambda(x,y)\defeq N\bigl(\It{y}{\varepsilon(x)}{u}{\chi(u)}\bigr),\\
 \beta_m(u)\defeq\sel(\lambda(\Suc u,m),u,\Suc u),\qquad
 \sigma(t,m)\defeq\It{t}{0}{u}{\beta_m(u)},\\
 P(m)\defeq\It{m}{0}{u}{\sel(\lambda(\Suc\Suc u,m),u,\Suc u)} .
\end{gather*}

\begin{lemma}[The comparator]
\label{kern:lem:comparator}
\begin{enumerate}[label=\textup{(\roman*)}]
\item $\chi(2u)=3u$, $\chi(1)=1$ and $\chi(3u)=3\chi(u)$.
\item $\It{y}{3a}{u}{\chi(u)}=3\,\It{y}{a}{u}{\chi(u)}$ and
 $\It{y}{1}{u}{\chi(u)}=1$.
\item $\lambda(0,y)=1$, $\lambda(\Suc x,0)=0$ and
 $\lambda(\Suc x,\Suc y)=\lambda(x,y)$.
\end{enumerate}
\end{lemma}

\begin{proof}
(i) $N(2u)=0$ and $G(2u)=3u$ give $\chi(2u)=3u$; $N(1)=1$ gives
$\chi(1)=1$.  For the third equation, $N(3u)=N(u)$.  If $N(u)=0$, then
$\chi(3u)=G(3u)=3G(u)+N(u)=3G(u)=3\chi(u)$.  If $N(u)=1$, then
$\chi(3u)=3u=3\chi(u)$.

(ii) The template $3\,\It{y}{a}{u}{\chi(u)}$ has the seed $3a$ and, by (i),
the update $\chi$; \rulename{UR} gives the first equation.  The second is
\cref{kern:lem:cu}(iv).

(iii) $\varepsilon(0)=1$, so $\lambda(0,y)=N(1)=1$ by (ii).
$\lambda(\Suc x,0)=N(\varepsilon(\Suc x))=N(2\varepsilon(x))=0$.  By the
second equation of \cref{kern:lem:uradd}, by (i) and by (ii),
\[
 \It{\Suc y}{2\varepsilon(x)}{u}{\chi(u)}
 =\It{y}{\chi(2\varepsilon(x))}{u}{\chi(u)}
 =\It{y}{3\varepsilon(x)}{u}{\chi(u)}
 =3\,\It{y}{\varepsilon(x)}{u}{\chi(u)},
\]
and $N(3v)=N(v)$ gives $\lambda(\Suc x,\Suc y)=\lambda(x,y)$.
\end{proof}

\begin{theorem}[Predecessor from uniqueness]
\label{kern:thm:urpred}
The computation equations, coverage and \rulename{UR} derive $P(0)=0$ and
$P(\Suc m)=m$, and hence that $\Suc x=\Suc y$ implies $x=y$.
\end{theorem}

\begin{proof}
By \eqref{kern:eq:selS} and \cref{kern:lem:comparator}(iii),
\[
 \beta_{\Suc m}(\Suc u)=\sel(\lambda(\Suc u,m),\Suc u,\Suc\Suc u)
 =\Suc\beta_m(u),\qquad
 \beta_{\Suc m}(0)=\sel(\lambda(0,m),0,1)=1 .
\]
As functions of $t$, the terms $\sigma(\Suc t,\Suc m)$ and
$\Suc\sigma(t,m)$ have the seed $1$ and the update $\beta_{\Suc m}$:
for the first by \rulename{IS}, with
$\sigma(1,\Suc m)=\beta_{\Suc m}(0)=1$, and for the second by
$\Suc\beta_m(\sigma(t,m))=\beta_{\Suc m}(\Suc\sigma(t,m))$.
\Cref{kern:lem:cu}(i) gives
\begin{equation}
 \sigma(\Suc t,\Suc m)=\Suc\sigma(t,m). \label{kern:eq:sat}
\end{equation}
The template $\lambda(\Suc n,n)$ satisfies
$\lambda(\Suc\Suc n,\Suc n)=\lambda(\Suc n,n)$ and has the value
$\lambda(1,0)=0$ at $0$, so $\lambda(\Suc n,n)=0$ by
\cref{kern:lem:cu}(ii).  The template $\sigma(n,n)$ has the seed $0$ and,
by \eqref{kern:eq:sat}, the update $\Suc$, so $\sigma(n,n)=n$ by
\rulename{UR} and \cref{kern:lem:cu}(iii).  Hence
\begin{equation}
 \sigma(\Suc n,n)=\beta_n(\sigma(n,n))=\sel(\lambda(\Suc n,n),n,\Suc n)=n.
 \label{kern:eq:satdiag}
\end{equation}
$P(0)=0$ is an instance of \rulename{I0}.  For a variable $m$ put
$\sigma'(t,m)\defeq\It{t}{0}{u}{\sel(\lambda(\Suc\Suc u,\Suc m),u,\Suc u)}$,
so that $P(\Suc m)$ is $\sigma'(\Suc m,m)$.  By
\cref{kern:lem:comparator}(iii) the body of $\sigma'$ equals $\beta_m(u)$
at every $u$, so $\sigma'(t,m)$ and $\sigma(t,m)$ have the same seed and
the same update, and \cref{kern:lem:cu}(i) gives
$\sigma'(t,m)=\sigma(t,m)$.  With \eqref{kern:eq:satdiag},
$P(\Suc m)=\sigma(\Suc m,m)=m$.  Applying $P$ to $\Suc x=\Suc y$ gives
$x=y$.
\end{proof}

Every appeal to induction in the construction of \cref{base:thm:pred} has
been replaced by an instance of \cref{kern:lem:cu}, and every case
distinction uses coverage and a finite range.

\subsection{Truncated subtraction and diagonal descent}
\label{kern:sec:urmonus}

Put $d(a,b)\defeq\It{b}{a}{u}{P(u)}$ and $\operatorname{sg}(a)\defeq
\sel(a,0,1)$.

\begin{lemma}[Truncated subtraction]
\label{kern:lem:monus}
$d(a,0)=a$, $d(a,\Suc b)=P(d(a,b))$, $d(\Suc a,\Suc b)=d(a,b)$,
$d(a,a)=0$, $d(0,a)=0$, $P(d(\Suc n,a))=d(n,a)$ and $d(a+n,n)=a$.
Addition is cancellative: $a+n=b+n$ implies $a=b$.
\end{lemma}

\begin{proof}
The first two are the computation equations.  As functions of $b$, the
terms $d(\Suc a,\Suc b)$ and $d(a,b)$ have the seed $a$, by
$d(\Suc a,1)=P(\Suc a)=a$, and the update $P$.  The template $d(a,a)$
satisfies $d(\Suc a,\Suc a)=d(a,a)$ and is $0$ at $0$;
\cref{kern:lem:cu}(ii) applies.  The template $d(0,a)$ and the constant $0$
have the seed $0$ and the update $P$.  For the sixth equation use coverage
on $a$: at $0$ it is $P(\Suc n)=n$, and at $\Suc a'$ it is
$P(d(n,a'))=d(n,\Suc a')$.  The template $d(a+n,n)$ satisfies
$d(a+\Suc n,\Suc n)=d(a+n,n)$ and is $a$ at $0$.  Cancellation follows by
applying $d(\,\cdot\,,n)$.
\end{proof}

\begin{lemma}[Diagonal descent]
\label{kern:lem:dd}
Let $q(x_1,\dots,x_r)$ be a term template, $r\ge 1$, such that
\[
 q(\Suc x_1,\dots,\Suc x_r)=q(x_1,\dots,x_r)
\]
and such that $q(x_1,\dots,x_r)=0$ whenever some $x_i$ is $0$.  Then
$q(x_1,\dots,x_r)=0$.
\end{lemma}

\begin{proof}
Coverage in each argument gives
$q(P(x_1),\dots,P(x_r))=q(x_1,\dots,x_r)$: if some $x_i$ is $0$, both sides
are $0$, and otherwise the first hypothesis applies.  The template
$f(k)\defeq q(d(x_1,k),\dots,d(x_r,k))$ therefore satisfies
$f(\Suc k)=f(k)$, and \cref{kern:lem:cu}(ii) gives $f(k)=f(0)$.  At
$k=x_1$ the first argument is $d(x_1,x_1)=0$, so $f(x_1)=0$; and $f(0)$ is
$q(x_1,\dots,x_r)$.
\end{proof}

The lemma uses uniqueness for one scalar state; it does not quantify over
tuples and uses no rule for simultaneous recursion.

A term is \emph{Boolean} if it is derivably $0$ or $1$.  Put
$\operatorname{or}(p,q)\defeq\sel(p,q,1)$ and
$\operatorname{xor}(p,q)\defeq\sel(p,q,Z(q))$.  For Boolean $p$ and $q$,
four cases show that $\operatorname{xor}(p,q)=0$ implies $p=q$.

\begin{lemma}[Three Boolean identities]
\label{kern:lem:bool}
\begin{enumerate}[label=\textup{(\roman*)}]
\item $Z(d(\Suc n,a))=\operatorname{sg}(d(a,n))$.
\item $Z\bigl(d(d(x,a),d(b,a))\bigr)
 =\operatorname{or}\bigl(Z(d(x,a)),Z(d(x,b))\bigr)$.
\item $\operatorname{or}(\operatorname{sg}(u),\operatorname{sg}(v))
 =\operatorname{sg}(u+v)$.
\end{enumerate}
\end{lemma}

\begin{proof}
(i) Let $q(n,a)$ be the $\operatorname{xor}$ of the two sides, which are
Boolean.  By \cref{kern:lem:monus}, $q(\Suc n,\Suc a)=q(n,a)$.  If $a=0$,
the sides are $Z(\Suc n)=0$ and $\operatorname{sg}(0)=0$.  If $n=0$ and
$a=\Suc b$, the sides are $Z(d(1,\Suc b))=Z(d(0,b))=1$ and
$\operatorname{sg}(\Suc b)=1$.  \Cref{kern:lem:dd} gives $q(n,a)=0$.

(ii) Let $q(x,a,b)$ be the $\operatorname{xor}$ of the two sides.
Simultaneous successors cancel in every subtraction, so
$q(\Suc x,\Suc a,\Suc b)=q(x,a,b)$.  If $x=0$, both sides are $1$.  If
$a=0$, the sides are $Z(d(x,b))$ and
$\operatorname{or}(Z(x),Z(d(x,b)))$; they agree by coverage on $x$, since
at $x=0$ both are $1$ and at a successor $Z(x)=0$.  If $b=0$, the sides
are $Z(d(x,a))$ and $\operatorname{or}(Z(d(x,a)),Z(x))$; they agree by
coverage on $x$ and the two values of the Boolean term $Z(d(x,a))$.
\Cref{kern:lem:dd} gives $q(x,a,b)=0$.

(iii) By coverage on $u$: at $0$ both sides are $\operatorname{sg}(v)$,
and at a successor both are $1$, by $\Suc u'+v=\Suc(u'+v)$.
\end{proof}

\subsection{An accumulator and the commutativity of the maximum}
\label{kern:sec:uracc}

Define
\[
 \mu(n,a)\defeq d(n,d(n,a)),\qquad
 \varphi(n,a,b)\defeq\mu(n,a)+\mu\bigl(d(n,a),d(b,a)\bigr),\qquad
 m(a,b)\defeq a+d(b,a).
\]
These terms are compositions of $d$ and addition; none is introduced by a
recursion whose body sees the counter.

\begin{lemma}[The accumulator]
\label{kern:lem:acc}
\begin{enumerate}[label=\textup{(\roman*)}]
\item $\mu(\Suc n,a)=\mu(n,a)+Z(d(\Suc n,a))$.
\item $\varphi(0,a,b)=0$ and
 $\varphi(\Suc n,a,b)=\varphi(n,a,b)+\operatorname{sg}\bigl(d(a,n)+d(b,n)\bigr)$.
\item $m(a,b)=m(P(a),P(b))+\operatorname{sg}(a+b)$.
\end{enumerate}
\end{lemma}

\begin{proof}
(i) Split $d(\Suc n,a)$ by coverage.  If it is $0$, then
$d(n,a)=P(0)=0$ by \cref{kern:lem:monus}, so $\mu(\Suc n,a)=\Suc n$ and
$\mu(n,a)=n$, and $Z(0)=1$.  If it is $\Suc r$, then $d(n,a)=r$ and
$\mu(\Suc n,a)=d(\Suc n,\Suc r)=d(n,r)=\mu(n,a)$, and $Z(\Suc r)=0$.

(ii) $\varphi(0,a,b)=\mu(0,a)+\mu(0,d(b,a))=0$, because $d(0,x)=0$.  Put
$w\defeq d(b,a)$ and split $d(\Suc n,a)$ again.  If it is $0$, then
$d(n,a)=0$, $\varphi(\Suc n,a,b)=\Suc n$, $\varphi(n,a,b)=n$ and
$Z(d(0,w))=1$.  If it is $\Suc r$, then $d(n,a)=r$, the first summand is
unchanged by (i), and (i) for the second summand gives
$\mu(\Suc r,w)=\mu(r,w)+Z(d(\Suc r,w))$.  In both cases
\[
 \varphi(\Suc n,a,b)=\varphi(n,a,b)+Z\bigl(d(d(\Suc n,a),d(b,a))\bigr).
\]
\Cref{kern:lem:bool}(ii) at $x=\Suc n$, then (i) of that lemma for each of
the two Boolean terms, then (iii), turn the last summand into
$\operatorname{sg}(d(a,n)+d(b,n))$.

(iii) By coverage on $a$ and $b$.  At $(0,0)$ both sides are $0$.  At
$(0,\Suc b)$ and at $(\Suc a,0)$ both sides are the successor of
$m(0,b)=b$ and of $m(a,0)=a$.  At $(\Suc a,\Suc b)$,
$m(\Suc a,\Suc b)=\Suc a+d(b,a)=\Suc m(a,b)$ and
$\operatorname{sg}(\Suc a+\Suc b)=1$.
\end{proof}

\begin{lemma}[Commutativity of the maximum]
\label{kern:lem:max}
$a+d(b,a)=b+d(a,b)$.
\end{lemma}

\begin{proof}
For fixed $a$, $b$ put
\[
 F(n)\defeq m(d(a,n),d(b,n))+\varphi(n,a,b),\qquad
 F'(n)\defeq m(d(b,n),d(a,n))+\varphi(n,a,b).
\]
By \cref{kern:lem:acc}(iii) at the arguments $d(a,n)$, $d(b,n)$, by
\cref{kern:lem:acc}(ii) and by the laws of addition, $F(\Suc n)=F(n)$, and
likewise $F'(\Suc n)=F'(n)$.  \Cref{kern:lem:cu}(ii) gives
$F(n)=F(0)=m(a,b)$ and $F'(n)=F'(0)=m(b,a)$.  At $n=b$, $d(b,b)=0$ and
$m(x,0)=x=m(0,x)$, so $F(b)$ and $F'(b)$ both equal
$d(a,b)+\varphi(b,a,b)$.  Hence $m(a,b)=m(b,a)$.
\end{proof}

The familiar proof of this identity is an induction whose hypothesis is used
at the predecessor.  Here the two sides are joined by a term $\varphi$
whose recurrence depends on the counter, although its definition is a
composition; the recurrence is \cref{kern:lem:acc}(ii).

\subsection{Reflection of equality and the proofs of the theorems}
\label{kern:sec:urreflect}

\begin{lemma}[Reflection of equality]
\label{kern:lem:urreflect}
Let $\Delta(a,b)\defeq\operatorname{sg}(d(a,b)+d(b,a))$.  Then
$\Delta(a,b)=0\vee\Delta(a,b)=1$, and $\Delta(a,b)=0$ holds if and only if
$a=b$.
\end{lemma}

\begin{proof}
By coverage and separation, $\operatorname{sg}(u)=0$ if and only if $u=0$,
and $u+v=0$ if and only if $u=0$ and $v=0$, since a sum with a successor
summand is a successor.  If $\Delta(a,b)=0$, then $d(a,b)=0=d(b,a)$, and
\cref{kern:lem:max} gives $a=a+d(b,a)=b+d(a,b)=b$.  If $a=b$, both
subtractions are $d(a,a)=0$.
\end{proof}

The equality reflected is the equality of the logic on the whole numerical
domain.

\begin{proof}[Proof of \cref{kern:thm:ur}]
By \cref{kern:prop:uniq}, every derivation of $T_{\mathrm{UR}}$ is a
derivation of $\Teq$ after each instance of \rulename{UR} is replaced by its
derivation.  Conversely let $\Gamma\proves s(0)=r(0)$ and
$\Gamma,s(z)=r(z)\proves s(\Suc z)=r(\Suc z)$ in $T_{\mathrm{UR}}$.  Put
$q(n)\defeq\Delta(s(n),r(n))$.  By \cref{kern:lem:urreflect} the two
premises give $\Gamma\proves q(0)=0$ and
$\Gamma,q(z)=0\proves q(\Suc z)=0$; \cref{kern:lem:zi} gives $q(t)=0$; and
\cref{kern:lem:urreflect} gives $s(t)=r(t)$.  The context $\Gamma$ is
arbitrary, so this is the contextual induction rule of $\Teq$.  The two
theories have the same language and the same derivable sequents, and
\cref{base:thm:bridge} applies to $T_{\mathrm{UR}}$.  In particular
induction for quantifier-free formulas (\cref{base:thm:qfi}) holds in
$T_{\mathrm{UR}}$.
\end{proof}

\begin{proof}[Proof of \cref{kern:thm:urnative}]
\rulename{UR} is derivable in $\Nat$ by \cref{kern:prop:uniq}.  Conversely
it suffices to derive equation induction in $\mathsf{U}$: given
$\Gamma,E(z)\proves_{\mathsf{U}}E(\Suc z)$, equation induction over the
context $\Gamma,E(0)$ gives $E(\Suc z')$ for a fresh $z'$, and the
introduction of $\Q{}$ by cases (\cref{kern:lem:local}(i)) gives
$\Gamma\proves_{\mathsf{U}}\Q{y}E(y)$.

The native connectives of \cref{base:sec:structural} obey their natural
deduction rules in $\mathsf{U}$, and the equality rules, the Leibniz rule
and numerical cases hold in $\mathsf{U}$: by \cref{base:rem:noprogress},
their derivations use the seven rules of $\Nat$ other than \rulename{QI}.
With these, $\mathsf{U}$ derives the reading of every axiom and rule of
$T_{\mathrm{UR}}$: the computation equations,
coverage by numerical cases, separation because the native falsity is the
equation $0=\Suc 0$, and \rulename{UR} itself.  Hence the derivation of the
induction rule in the proof of \cref{kern:thm:ur}, written with the native
connectives, is a derivation in $\mathsf{U}$.
\end{proof}

\begin{remark}[Sources and status]
\label{kern:rem:ur-sources}
\Crefrange{kern:sec:ur}{kern:sec:urreflect} are \pkg{scalar\_uniqueness},
except for \cref{kern:prop:rinv}, the proof of \cref{kern:lem:rigid} and
\cref{kern:cor:rigidcov}, which are from \pkg{reductions}, and
\cref{kern:thm:su}, which is from \pkg{orbit\_closure}
(\cref{kern:rem:scopekernel}).  \pkg{scalar\_uniqueness} attributes the
idea of replacing induction by uniqueness of recursion to the literature
(Okada and Scott study uniqueness of iteration in a typed setting) and
claims the completion for the scalar language only: one state, no counter
in the body, no pairing.  \Cref{kern:lem:monus} also yields the Mal'cev
term $d(x+z,y)$, which the package notes.  In Lean,
\cref{kern:prop:uniq} is proved, and so is \cref{kern:thm:ur}: the modules
\Lean{UrQuotient}, \Lean{UrArith}, \Lean{UrRetract}, \Lean{UrPred},
\Lean{UrMonus} and \Lean{UrReflect} state $T_{\mathrm{UR}}$ with
\rulename{UR} as an axiom scheme (\Lean{TurAx}) and prove
\Lean{tur\_iff\_teq}, that $T_{\mathrm{UR}}$ and $\Teq$ derive the same
sequents.  The formal proof follows
\crefrange{kern:sec:urfirst}{kern:sec:urreflect} step by step: every appeal
to \cref{kern:lem:cu}(i) is one application of a derived rule on classes of
terms whose step is a computation, the hypotheses on templates in
\cref{kern:prop:rinv,kern:lem:rigid,kern:lem:zi} are stated over all
extensions of the context, those of \cref{kern:lem:dd} over every context,
and equation induction is \cref{kern:lem:zi} on $\Delta(s(n),r(n))$
(\Lean{Tur.eqInd}).  The native form \cref{kern:thm:urnative} is not
formalized.
\end{remark}

\subsection{Other forms of the uniqueness rule}
\label{kern:sec:urforms}

Three variations of \rulename{UR} occur in \pkg{orbit\_closure} and
\pkg{crossed\_uniqueness}.  They are related to \rulename{UR} as follows.
The rule \rulename{Orb} concludes the closure of the equation of
\rulename{UR} and takes no target; it is interderivable with \rulename{UR}
by the rules for the closure alone (\cref{kern:prop:orbit}).  The rule
\rulename{SU} of \cref{kern:def:SU} is the instance of \rulename{UR} at one
body; it gives induction for zero equations (\cref{kern:thm:su}) and the
laws of addition, and whether it gives induction for all equations is open.
The rule \rulename{UX} has a premise for the seed in place of the seed
$f(0)$ and a crossed conclusion; \rulename{UR} is its instance at the seed
$f(0)$, and with transport it follows from \rulename{UR}.  Taken as the
only arithmetic rule on fully normalized terms, \rulename{UX} derives
transport as well, so that it replaces \rulename{I0}, \rulename{IS},
\rulename{Tr} and \rulename{QI} together (\cref{kern:thm:crossed}).

\subsubsection*{The closure of the orbit equation}

For term templates $f$ and $h$ as in \rulename{UR}, and a variable $y$ not
free in $h$, the ordinary and the native form of the rule are
\[
 \infrule{\Gamma\proves f(\Suc z)=h(f(z))}
         {\Gamma\proves\forall y\,\bigl(f(y)=\It{y}{f(0)}{u}{h(u)}\bigr)}
 \qquad\qquad
 \infrule{\Gamma\proves f(\Suc z)=h(f(z))}
         {\Gamma\proves\Q{y}\bigl(f(y)=\It{y}{f(0)}{u}{h(u)}\bigr)}\
 \rulename{Orb},
\]
where $z$ is not free in $\Gamma$ or in $h$.

\begin{definition}[The calculus $\mathsf{U}_{\mathrm{o}}$ and the theory
$T_{\mathrm{Orb}}$]
\label{kern:def:Uo}
$\mathsf{U}_{\mathrm{o}}$ is $\Nat$ with \rulename{QI} replaced by the
native form of \rulename{Orb}; its eight rules are \rulename{Hyp},
\rulename{I0}, \rulename{IS}, \rulename{Tr}, \rulename{QE}, \rulename{DE},
\rulename{DI} and \rulename{Orb}.  $T_{\mathrm{Orb}}$ is $T_{\mathrm{UR}}$
with \rulename{UR} replaced by the ordinary form of \rulename{Orb}.
\end{definition}

\begin{proposition}[Orbit uniqueness]
\label{kern:prop:orbit}
$T_{\mathrm{Orb}}$ and $T_{\mathrm{UR}}$ derive the same sequents of
ordinary formulas, and $\mathsf{U}_{\mathrm{o}}$ and $\mathsf{U}$ are
equivalent.  Hence $T_{\mathrm{Orb}}$ derives the sequents of $\Teq$, and
$\mathsf{U}_{\mathrm{o}}$ is equivalent to $\Nat$.
\end{proposition}

\begin{proof}
Let $E(y)$ be the equation $f(y)=\It{y}{f(0)}{u}{h(u)}$.  In the ordinary
theories, an instance of \rulename{UR} is \rulename{Orb} followed by the
elimination of $\forall$ at $t$; and \rulename{UR} at a fresh variable $w$
as its target, followed by the introduction of $\forall$, is \rulename{Orb}.

In the native calculi, \rulename{Orb} is derived in $\mathsf{U}$ thus:
\rulename{UR} at the target $\Suc w$, with $w$ fresh, gives $E(\Suc w)$,
and the introduction of $\Q{}$ by cases (\cref{kern:lem:local}(i)), whose
assumption $E(0)$ is not needed, gives $\Q{y}E(y)$.  Conversely
\rulename{UR} is derived in $\mathsf{U}_{\mathrm{o}}$ by \rulename{QE} at
$t$: its second premise $E(0)$ is the instance
$\It{0}{f(0)}{u}{h(u)}=f(0)$ of \rulename{I0}, reversed.  Both
constructions hold over arbitrary contexts.  The last sentence of the
proposition is \cref{kern:thm:ur,kern:thm:urnative}.
\end{proof}

The base $E(0)$ of the orbit equation is a theorem, so $\Q{y}E$ and the
unconditional $\forall y\,E$ say the same; in a calculus with clause
contexts (\cref{base:sec:clauses}) \rulename{Orb} adds the clause of $E$
with a trivial antecedent, and the term at which it is used is chosen when
the clause is used.  The rule removes a parameter of the certificate and
not the principle.

\begin{remark}[The body must not see the counter]
\label{kern:rem:orbcounter}
In \rulename{Orb} the variable $y$ of the conclusion is not free in $h$.
A checker that reads $f$ and $h$ off a proposed closure has to test this,
also below nested binders.  Without the test, $f(y)\defeq y$ and the body
$\Suc y$ have the premise $\Suc z=\Suc z$ and the conclusion
$\Q{y}\bigl(y=\It{y}{0}{u}{\Suc y}\bigr)$, whose instance at $1$ is
$1=2$.
\end{remark}

\subsubsection*{The successor body}

\begin{proposition}[Addition from successor uniqueness]
\label{kern:prop:suadd}
$T_{\mathrm{SU}}$ derives that addition is associative and commutative with
the neutral element $0$.  If $\Gamma\proves k(\Suc z)=k(z)$ for a term
template $k$ and a fresh $z$, then $T_{\mathrm{SU}}$ derives
$\Gamma\proves n+k(n)=n+k(0)$.
\end{proposition}

\begin{proof}
The laws of addition are those of \cref{kern:lem:uradd}, whose proofs
compare terms with the update $\Suc$.  For the second statement let
$f(n)\defeq n+k(n)$.  Then $f(0)=0+k(0)=k(0)$ and
\[
 f(\Suc n)=\Suc n+k(\Suc n)=\Suc n+k(n)=\Suc(n+k(n))=\Suc f(n),
\]
so \rulename{SU} gives $f(n)=\It{n}{k(0)}{u}{\Suc u}=k(0)+n=n+k(0)$.
\end{proof}

If addition were cancellative in $T_{\mathrm{SU}}$, the proposition would
give $k(n)=k(0)$, which is \cref{kern:lem:cu}(ii), uniqueness for the body
$u$.  Cancellation is proved in \cref{kern:lem:monus} from the predecessor,
whose construction uses other bodies.

\begin{remark}[What successor uniqueness is not known to give]
\label{kern:rem:sugap}
To pass from \cref{kern:thm:su} to induction for all equations by the route
of \cref{kern:sec:urreflect}, a term $\Delta$ with
$\Delta(a,b)=0\leftrightarrow a=b$ is needed.  The construction of
\crefrange{kern:sec:urpred}{kern:sec:uracc} uses \rulename{UR} for the
bodies $u$, $Z(u)$, $j(u)$, $g(u)$, $\chi(u)$, $\beta_m(u)$ and $P(u)$,
among others, and cannot be repeated in $T_{\mathrm{SU}}$ as it stands.
The right inverse of \cref{kern:prop:rinv} does not supply such a term
either: it concerns a term $e$ whose values are already known at $0$ and
along successors of its values, and \pkg{reductions} records that it does
not construct $\Delta$ from \rulename{SU}.  \pkg{orbit\_closure} does not claim that $T_{\mathrm{SU}}$ is equivalent to
$\Teq$.  It also records that a model used in another package,
\pkg{canonical\_recursors}, to separate induction for zero equations from
induction for all equations does not decide the question: in that model,
for a nonstandard counter $p$, the term $e(a)\defeq\It{p}{a}{u}{u}$
satisfies $e(0)=0$ and $e(\Suc a)=\Suc e(a)$ and is not the identity,
whereas \rulename{SU} with \cref{kern:lem:cu}(iii) gives $e(a)=a$.  The
model therefore refutes \rulename{SU}, and it is not a model of
$T_{\mathrm{SU}}$ in which equation induction fails.  The model is not
reproduced here.
\end{remark}

\subsubsection*{A seed premise and a crossed conclusion, on normal forms}

Let $\nf(t)$ be the full normal form of a term $t$ for the two computation
equations, read from left to right, and let $\nf(A)$ and $\nf(\Gamma)$ be
obtained by normalizing every term.  The following facts are taken from the
packages \pkg{full\_normalization} and \pkg{computation\_normal}, which
belong to another chapter: $\nf$ is total; $\nf(\nf(t)[\theta])=\nf(t[\theta])$
for every substitution $\theta$; and $\Nat$ derives $t=\nf(t)$.  A term, a
formula or a context is \emph{normal} if $\nf$ fixes it.  For normal terms
the selector satisfies $\nf(\sel(0,a,b))=a$ and $\nf(\sel(\Suc c,a,b))=b$,
and $\nf(\It{1}{a}{u}{h})=\nf(h(a))$.

\begin{definition}[The calculus $\mathsf{U}_{\mathrm{c}}$]
\label{kern:def:Uc}
The sequents of $\mathsf{U}_{\mathrm{c}}$ consist of normal formulas, and
every instance $B(t)$, $f(n)$ or $h(f(z))$ in a rule denotes the normal
form of the substitution.  Its six rules are \rulename{Hyp}, \rulename{QE},
\rulename{DE}, \rulename{DI} in their forms of \cref{base:sec:kernel},
reflexivity,
\[
 \infrule{}{\Gamma\proves t=t}\ \rulename{Ref},
\]
and \emph{seeded crossed uniqueness},
\[
 \infrule{\Gamma\proves f(0)=a\qquad\Gamma\proves f(\Suc z)=h(f(z))}
         {\Gamma\proves\It{n}{a}{u}{h(u)}=f(n)}\ \rulename{UX},
\]
where $f(z)$ and $h(u)$ are term templates and $z$ is not free in $\Gamma$,
$a$, $n$ or $h$.
\end{definition}

There is no computation rule, no transport rule and no rule for progress:
computation is in the syntax, and normalization is part of the
specification of the calculus and of any checker for it.  The introduction
\rulename{DI} is by cases at every matrix.  The body $h$ receives the state
and no counter, and the second premise has no induction hypothesis.  This
is the one calculus of the chapter without the rules \rulename{I0} and
\rulename{IS}.

\begin{lemma}[Structural rules on normal forms]
\label{kern:lem:ucstruct}
Weakening and cut are admissible in $\mathsf{U}_{\mathrm{c}}$, and if
$\Gamma\proves_{\mathsf{U}_{\mathrm{c}}}A$ then
$\nf(\Gamma[\theta])\proves_{\mathsf{U}_{\mathrm{c}}}\nf(A[\theta])$.
\end{lemma}

\begin{proof}
By induction on derivations, with eigenvariables renamed.  No rule has a
mode or a test on a term, and
$\nf(\nf(t)[\theta])=\nf(t[\theta])$ identifies the substituted premises
and conclusions of each rule with an instance of the same rule.
\end{proof}

\begin{lemma}[Equality from seeded crossed uniqueness]
\label{kern:lem:uxeq}
$\mathsf{U}_{\mathrm{c}}$ derives, over an arbitrary context:
\begin{enumerate}[label=\textup{(\roman*)}]
\item numerical cases and the introduction of $\Q{}$ by cases;
\item symmetry: from $s=t$ to $t=s$;
\item congruence: from $s=t$ to $H(s)=H(t)$, for every term template $H$;
\item joining: from $s=t$ and $s=r$ to $t=r$;
\item transitivity, and the rule \rulename{Tr}: from $s=t$ and $H(s)=r$ to
 $r=H(t)$;
\item the rule \rulename{UR}.
\end{enumerate}
None of these derivations uses an induction, a predecessor or a test of
equality.
\end{lemma}

\begin{proof}
(i) \Cref{kern:lem:local}(i); its derivations use \rulename{Hyp},
\rulename{QE}, \rulename{DE} and \rulename{DI}.

(ii) Apply \rulename{UX} with the constant template $f(z)\defeq s$, the
seed $a\defeq t$, the body $h(u)\defeq u$ and the counter $n\defeq 0$.  The
first premise is $s=t$, the second is $s=s$, and the conclusion
$\It{0}{t}{u}{u}=s$ is, in normal form, $t=s$.

(iii) Apply \rulename{UX} with $f(z)\defeq\It{z}{s}{u}{H(u)}$,
$a\defeq t$, $h\defeq H$ and $n\defeq 1$.  The first premise is $s=t$,
because $f(0)$ has the normal form $s$.  The two sides of the second
premise, $\It{\Suc z}{s}{u}{H(u)}$ and $H(\It{z}{s}{u}{H(u)})$, have the
same normal form, so it is an instance of \rulename{Ref}.  The conclusion is
$H(t)=H(s)$ in normal form, and (ii) reverses it.  The template may use its
variable several times or not at all, also below binders.

(iv) Apply \rulename{UX} with $f(z)\defeq\sel(z,s,r)$, $a\defeq t$,
$h(u)\defeq u$ and $n\defeq 1$.  The first premise is $s=t$.  The second
premise is $r=\sel(z,s,r)$, since $f(\Suc z)$ has the normal form $r$; it
follows by numerical cases on $z$, from $r=s$ at $0$, which is (ii) applied
to $s=r$, and from $r=r$ at a successor.  The conclusion is $t=r$ in normal
form.

(v) From $s=t$ and $t=r$: (ii) gives $t=s$, and (iv) gives $s=r$.  From
$s=t$ and $H(s)=r$: (iii) gives $H(s)=H(t)$, and (iv) gives $r=H(t)$.

(vi) \rulename{UX} with the seed $a\defeq f(0)$ and the first premise by
\rulename{Ref}, followed by (ii).
\end{proof}

The counter $0$ makes \rulename{UX} a replacement at the root, the counter
$1$ makes it an application of a template, and the selector joins two
equations; each conclusion is read after normalization, which is why the
calculus is stated on normal forms.

\begin{theorem}[Seeded crossed uniqueness]
\label{kern:thm:crossed}
For every context $\Gamma$ and every native formula $A$,
\[
 \Gamma\proves A\ \text{in }\Nat
 \quad\Longleftrightarrow\quad
 \nf(\Gamma)\proves_{\mathsf{U}_{\mathrm{c}}}\nf(A).
\]
Hence the bridge of \cref{base:thm:bridge} holds for
$\mathsf{U}_{\mathrm{c}}$, with the translations composed with $\nf$.
\end{theorem}

\begin{proof}
\emph{From right to left.}  Every rule of $\mathsf{U}_{\mathrm{c}}$ is
derivable in $\Nat$ on normal formulas.  \rulename{Ref} is reflexivity.  For
\rulename{UX} apply equation induction to $\It{y}{a}{u}{h(u)}=f(y)$: the
base follows from \rulename{I0} and the first premise, and in the step
\rulename{IS}, congruence under $h$ applied to the induction hypothesis,
and the second premise give
$\It{\Suc z}{a}{u}{h(u)}=h(\It{z}{a}{u}{h(u)})=h(f(z))=f(\Suc z)$.  The
normalization of an instance is justified by $t=\nf(t)$ and the Leibniz
rule.  The same two facts make $A$ and $\nf(A)$ interderivable in $\Nat$,
so $\nf(\Gamma)\proves\nf(A)$ gives $\Gamma\proves A$.

\emph{From left to right.}  First, if $\Gamma\proves_{\mathsf{U}}A$ in the
calculus $\mathsf{U}$ of \cref{kern:def:U}, then
$\nf(\Gamma)\proves_{\mathsf{U}_{\mathrm{c}}}\nf(A)$, by induction on the
derivation.  \rulename{Hyp} is clear.  The two sides of an instance of
\rulename{I0} or \rulename{IS} have the same normal form, so the instance
becomes \rulename{Ref}.  \rulename{Tr} becomes
\cref{kern:lem:uxeq}(v) with the template $\nf(h)$, because
$\nf(h(s))=\nf(\nf(h)(\nf(s)))$.  \rulename{QE}, \rulename{DE} and
\rulename{DI} become the same rules, by the same identity.  \rulename{UR}
becomes \cref{kern:lem:uxeq}(vi).  Second, $\Gamma\proves A$ in $\Nat$
implies $\Gamma\proves_{\mathsf{U}}A$ by \cref{kern:thm:urnative}.
\end{proof}

\begin{remark}[Sources, overlaps and status]
\label{kern:rem:urforms-sources}
\Cref{kern:prop:orbit,kern:prop:suadd,kern:rem:orbcounter,kern:rem:sugap}
are from \pkg{orbit\_closure}, which states \rulename{Orb} for the ordinary
theory and describes the native form with its checker.  Its translator maps
a derivation with \rulename{UR} to one with \rulename{Orb} by the second
construction in the proof of \cref{kern:prop:orbit}.
\Cref{kern:def:Uc,kern:lem:uxeq,kern:thm:crossed} are from
\pkg{crossed\_uniqueness}.  That package also derives transport in a
formula template inside $\mathsf{U}_{\mathrm{c}}$ by induction on the
template, and restates the native connectives through $\D{}{}$ with an
unused coordinate; here both follow from \cref{kern:thm:crossed} and
\cref{base:sec:structural}.  Both packages take the passage from
\rulename{UR} to equation induction from \pkg{scalar\_uniqueness}, as the
proofs above do, and neither claims it as new.  Neither package is
formalized in Lean.
\end{remark}

% ---------------------------------------------------------------------
\section{The progressive reading of the closure}
\label{kern:sec:prog}

This section draws on \pkg{uniform\_progress}, \pkg{target\_closure} and the
note on mixed closure in \pkg{scope\_and\_closure}.

The single introduction of \cref{kern:sec:witness} tests the witness.  A
single introduction whose mode depends on the root of the matrix alone is
possible if the meaning of $\Q{y}E$ at a literal equation $E$ is changed
from $E(0)\to\forall y\,E(y)$ to progressivity,
$\forall y\,(E(y)\to E(\Suc y))$.  The calculi of this section have seven
rules, keep all native formulas, and are stable under substitution or under
renaming; they are equivalent to $\Nat$ through two translations and not on
the same formulas.

\subsection{The calculus with orbit elimination}
\label{kern:sec:R}

Extend the selector of modes of \cref{kern:sec:det} to terms:
$\eta_B(a)\defeq a$ if the matrix $B$ is a literal equation and
$\eta_B(a)\defeq 0$ if $B$ is compound.  The \emph{orbit term} is
\[
 \omega(a,n)\defeq\It{a}{n}{u}{\Suc u},
\]
with the counter $a$ and the seed $n$.

\begin{definition}[The calculus $\mathsf{R}$]
\label{kern:def:R}
$\mathsf{R}$ has the seven rules \rulename{Hyp}, \rulename{I0},
\rulename{IS}, $\rulename{Tr}_{2}$, \rulename{DE} and
\begin{gather*}
 \infrule{\Gamma\proves\Q{y}B(y)\qquad\Gamma\proves B(\eta_B(a))}
         {\Gamma\proves B\bigl(\omega(\eta_B(a),n)\bigr)}\
 \rulename{QE}_{\omega}
 \\[1.2ex]
 \infrule{\Gamma,\ B(t,\eta_B(z))\proves B(t,\Suc z)}
         {\Gamma\proves\D{x}{y}B(x,y)}\ \rulename{DI}_{\eta},
\end{gather*}
where $a$ and $n$ are arbitrary terms, and $z$ is not free in $\Gamma$, in
$t$ or in $\D{x}{y}B$.  The designated contradiction is $0=\Suc 0$.
\end{definition}

At a compound matrix $\rulename{QE}_{\omega}$ passes from $B(0)$ to
$B(\omega(0,n))$, and $\omega(0,n)=n$ is an instance of \rulename{I0}.  At a
literal equation $E$ it passes from $E(a)$ to $E(\omega(a,n))$, the instance
$n$ steps further.  Transport in a compound formula is not a rule, and there
is no rule of numerical cases.

\begin{definition}[The progressive reading]
\label{kern:def:preading}
The ordinary formula $F^{\mathrm{p}}$ is defined by
\begin{align*}
 (s=t)^{\mathrm{p}} &\defeq (s=t),\\
 (\Q{y}E)^{\mathrm{p}} &\defeq \forall y\,\bigl(E(y)\to E(\Suc y)\bigr)
  && (E\text{ a literal equation}),\\
 (\Q{y}B)^{\mathrm{p}} &\defeq
  B^{\mathrm{p}}(0)\to\forall y\,B^{\mathrm{p}}(y)
  && (B\text{ compound}),\\
 (\D{x}{y}B)^{\mathrm{p}} &\defeq \exists x\,(\Q{y}B)^{\mathrm{p}} .
\end{align*}
\end{definition}

The reading commutes with substitution, because substitution preserves the
root of a matrix.  The elimination and the reading have to change together:

\begin{proposition}[The two readings differ]
\label{kern:prop:parity}
Let $E(y)$ be the equation $N(y)=1$ for the parity term $N$ of
\cref{base:sec:terms}.
\begin{enumerate}[label=\textup{(\roman*)}]
\item In the standard model $E(0)\to\forall y\,E(y)$ is true and
 $\forall y\,(E(y)\to E(\Suc y))$ is false.
\item The rule $\rulename{QE}_{\omega}$ does not preserve truth for the
 ordinary reading of \cref{base:sec:reading}.
\item $\Nat$ derives $\Q{y}E(y)$, and $\mathsf{R}$ does not.
\end{enumerate}
\end{proposition}

\begin{proof}
$N(0)=0$, $N(1)=1$ and $N(2)=0$.  (i) The premise $E(0)$ is false, and
progress fails at $y=1$.  (ii) With $a\defeq 1$ and $n\defeq 1$ the two
premises $\Q{y}E$ and $E(1)$ are true in the ordinary reading and the
conclusion $E(\omega(1,1))$, which is $E(2)$, is false.  (iii) In $\Nat$ the
closure follows by the introduction by cases from $N(0)=1\proves
N(\Suc z)=1$, which holds because $N(0)=0$ and $0=1$ derives every equation.
In $\mathsf{R}$ it is not derivable by \cref{kern:thm:rsound} below and (i).
\end{proof}

\subsection{Derived rules of the orbit calculus}
\label{kern:sec:Rrules}

\begin{lemma}[Structural rules]
\label{kern:lem:rstruct}
Weakening, cut and substitution of terms are admissible in $\mathsf{R}$, and
substitution maps every rule instance to an instance of the same rule.
\end{lemma}

\begin{proof}
\Cref{kern:lem:struct}: no rule of $\mathsf{R}$ tests a term, and
$\eta_B(a)[\theta]=\eta_{B[\theta]}(a[\theta])$.
\end{proof}

\begin{lemma}[Equality, closure introduction, equation induction]
\label{kern:lem:rbasic}
$\mathsf{R}$ derives:
\begin{enumerate}[label=\textup{(\roman*)}]
\item reflexivity, symmetry, transitivity, congruence in a term template,
 and \rulename{Tr};
\item the introduction $\rulename{QI}_{\eta}$ of \cref{kern:def:Ndet}: from
 $\Gamma,B(\eta_B(z))\proves B(\Suc z)$ with $z$ fresh to
 $\Gamma\proves\Q{y}B(y)$;
\item equation induction: from $\Gamma\proves E(0)$ and
 $\Gamma,E(z)\proves E(\Suc z)$, for a literal equation $E$ and $z$ fresh,
 to $\Gamma\proves E(t)$.
\end{enumerate}
The derivations of (i) and (ii) use \rulename{Hyp}, \rulename{I0},
$\rulename{Tr}_{2}$, \rulename{DE} and $\rulename{DI}_{\eta}$ only.
\end{lemma}

\begin{proof}
(i) Reflexivity and symmetry are derived from \rulename{I0} and
$\rulename{Tr}_{2}$ in the proof of \cref{kern:thm:onecontext}.
Transitivity and congruence are $\rulename{Tr}_{2}$ applied to $r=s$ in the
template $r=v$ and to $f(s)=f(s)$ in the template $f(s)=f(v)$.
\rulename{Tr} follows as in that proof.

(ii) Apply $\rulename{DI}_{\eta}$ with the matrix $B(y)$, a coordinate $x$
not free in $B$ and the witness $0$; then \rulename{DE}, whose opened
hypothesis is $\Q{y}B(y)$, and \rulename{Hyp}.

(iii) By (ii) the step gives $\Q{y}E(y)$.  $\rulename{QE}_{\omega}$ with
$a\defeq 0$ and $n\defeq t$ gives $E(\omega(0,t))$ from the base, and
$\omega(0,t)=t$ is an instance of \rulename{I0}; transport in the equation
template $E$ gives $E(t)$.
\end{proof}

The order of the arguments of $\omega$ is what makes (iii) work:
$\omega(0,t)=t$ is a computation, whereas $\omega(t,0)=t$ is a theorem that
needs the induction under construction.

\begin{lemma}[Orbit identities]
\label{kern:lem:orbit}
In every calculus or theory with the equality rules, the computation
equations and equation induction, in particular in $\mathsf{R}$, $\Nat$ and
$\Teq$,
\begin{gather*}
 \omega(0,n)=n,\qquad \omega(\Suc a,n)=\Suc\,\omega(a,n),\qquad
 \omega(a,0)=a,\\
 \omega(a,\Suc n)=\Suc\,\omega(a,n),\qquad
 \omega(a,n)=\omega(n,a),\qquad \omega(a,1)=\Suc a .
\end{gather*}
\end{lemma}

\begin{proof}
The first two are \rulename{I0} and \rulename{IS}.  The third is by
induction on $a$: $\omega(0,0)=0$, and
$\omega(\Suc a,0)=\Suc\,\omega(a,0)=\Suc a$.  The fourth is by induction on
$a$ with the parameter $n$: both sides are $\Suc n$ at $0$, and
$\omega(\Suc a,\Suc n)=\Suc\,\omega(a,\Suc n)=\Suc\Suc\,\omega(a,n)
=\Suc\,\omega(\Suc a,n)$.  The fifth is by induction on $a$: the base is
$n=\omega(n,0)$, and
$\omega(\Suc a,n)=\Suc\,\omega(a,n)=\Suc\,\omega(n,a)=\omega(n,\Suc a)$.
The last is $\omega(a,1)=\Suc\,\omega(a,0)=\Suc a$.
\end{proof}

\begin{lemma}[The two interfaces of a closure, and transport]
\label{kern:lem:rtransport}
In $\mathsf{R}$, over an arbitrary context:
\begin{enumerate}[label=\textup{(\roman*)}]
\item one step: from $\Q{y}E(y)$ and $E(a)$ to $E(\Suc a)$, for a literal
 equation $E$;
\item transport: from $s=t$ and $F(s)$ to $F(t)$, for every native formula
 template $F$;
\item instance from the base: from $\Q{y}B(y)$ and $B(0)$ to $B(t)$, for
 every matrix $B$;
\item homogeneous congruence: if $B$ and $C$ are interderivable over
 $\Gamma$ and are both literal equations or both compound, then $\Q{y}B$
 and $\Q{y}C$ are interderivable over $\Gamma$, and so are $\D{x}{y}B$ and
 $\D{x}{y}C$, where the bound variables are not free in $\Gamma$.
\end{enumerate}
\end{lemma}

\begin{proof}
(i) $\rulename{QE}_{\omega}$ with $n\defeq 1$ gives $E(\omega(a,1))$, and
$\omega(a,1)=\Suc a$ by \cref{kern:lem:orbit}.

(ii) and (iii) are proved together by induction on the number of binders of
the formula, a measure that substitution does not change; (iii) for $B$ uses
(ii) for $B$.  For an equation, (ii) is $\rulename{Tr}_{2}$.  For (iii): at
a literal equation, $\rulename{QE}_{\omega}$ with $a\defeq 0$ and
$n\defeq t$ gives $B(\omega(0,t))$; at a compound matrix it gives the same
from $B(0)$; and (ii) for $B$ along $\omega(0,t)=t$ gives $B(t)$.

(ii) for $\Q{y}B(v,y)$.  Derive the target $\Q{y}B(t,y)$ by
$\rulename{QI}_{\eta}$: under the assumption $B(t,\eta_B(z))$, transport in
$B$ along $t=s$ gives $B(s,\eta_B(z))$; then (i), if $B$ is a literal
equation, or (iii) at $\Suc z$, if $B$ is compound, gives $B(s,\Suc z)$ from
the given closure $\Q{y}B(s,y)$; and transport in $B$ along $s=t$ gives
$B(t,\Suc z)$.

(ii) for $\D{x}{y}B(v,x,y)$.  Apply \rulename{DE} to the given formula with
a fresh $b$, and derive the target by $\rulename{DI}_{\eta}$ with the
witness $b$: the premise
$B(t,b,\eta_B(z))\proves B(t,b,\Suc z)$ follows from the opened closure
$\Q{y}B(s,b,y)$ exactly as in the previous case.  The conclusion does not
contain $b$.

Every use of transport or of (iii) inside these two cases is for the matrix
$B$, which has fewer binders.

(iv) Repeat the two cases of (ii) with the given derivations between $B$
and $C$, instantiated by substitution (\cref{kern:lem:rstruct}), in place
of transport.  Because $B$ and $C$ have the same kind of root,
$\eta_B=\eta_C$, so the assumptions of the two introductions match.
\end{proof}

The restriction in (iv) is necessary.  An equation $E$ and a compound
formula $C$ can be interderivable while $\Q{y}E$, which is progressive, and
$\Q{y}C$, which is conditional, are not; the padding of
\cref{kern:lem:pad} gives such pairs.  Transport is unaffected, because it
preserves the root of the matrix.

\subsection{Soundness for the progressive reading}
\label{kern:sec:Rsound}

\begin{theorem}[Soundness]
\label{kern:thm:rsound}
If $\Gamma\proves_{\mathsf{R}}F$, then
$\Gamma^{\mathrm{p}}\proves_{\Teq}F^{\mathrm{p}}$.
\end{theorem}

\begin{proof}
By induction on the derivation.  \rulename{Hyp}, the computation equations,
$\rulename{Tr}_{2}$ and \rulename{DE} have their ordinary meanings; for
\rulename{DE} the opened hypothesis reads $(\Q{y}B(b,y))^{\mathrm{p}}$, the
matrix of the existential quantifier.

$\rulename{DI}_{\eta}$ at a literal equation: the premise gives
$\forall z\,(E(t,z)\to E(t,\Suc z))$, and $t$ is a witness.  At a compound
matrix: the premise gives $B^{\mathrm{p}}(t,\Suc z)$ under
$B^{\mathrm{p}}(t,0)$, so coverage gives
$B^{\mathrm{p}}(t,0)\to\forall y\,B^{\mathrm{p}}(t,y)$.

$\rulename{QE}_{\omega}$ at a compound matrix: the first premise reads
$B^{\mathrm{p}}(0)\to\forall y\,B^{\mathrm{p}}(y)$, the second gives
$B^{\mathrm{p}}(0)$, and $\omega(0,n)=n$.  At a literal equation: assume
$\forall y\,(E(y)\to E(\Suc y))$ and $E(a)$.  Induction on $k$ for the
literal equation $E(\omega(k,a))$, with these two assumptions in the
context, gives $E(\omega(k,a))$ for every $k$: the base is $E(a)$ by
$\omega(0,a)=a$, and the step is the assumed progress at $\omega(k,a)$ with
$\omega(\Suc k,a)=\Suc\,\omega(k,a)$.  At $k\defeq n$, the equation
$\omega(n,a)=\omega(a,n)$ of \cref{kern:lem:orbit} gives $E(\omega(a,n))$.
\end{proof}

The use of the commutativity of $\omega$ in this proof is harmless: it is a
theorem of $\Teq$, whose induction is given.  It is separate from the
recovery of induction inside $\mathsf{R}$, which uses $\omega(0,t)=t$ only.

\subsection{The embedding of the core}
\label{kern:sec:Rup}

For a literal equation $E$ and fresh variables $r$, $v$ let
\[
 \operatorname{pad}(E)\defeq\D{r}{v}\operatorname{gate}(v,E),
 \qquad\text{that is,}\qquad
 \D{r}{v}\bigl(\sel(v,0,s)=\sel(v,0,t)\bigr)\quad\text{for }E\equiv(s=t),
\]
and for a native formula $F$ let $F^{\sharp}\defeq\operatorname{pad}(F)$ if
$F$ is a literal equation and $F^{\sharp}\defeq F$ if $F$ is compound.  The
formula $F^{\sharp}$ is always compound.  The translation $A^{\uparrow}$ of a
native formula is
\[
 (s=t)^{\uparrow}\defeq(s=t),\qquad
 (\Q{y}B)^{\uparrow}\defeq\Q{y}(B^{\uparrow})^{\sharp},\qquad
 (\D{x}{y}B)^{\uparrow}\defeq\D{x}{y}(B^{\uparrow})^{\sharp}.
\]
It fixes terms and equations, preserves the kind of the root, commutes with
substitution, and puts a compound matrix below every binder.

\begin{lemma}[Padding]
\label{kern:lem:pad}
For a literal equation $E$, the formulas $\operatorname{pad}(E)$ and $E$
are interderivable in $\mathsf{R}$ over every context.  Hence $F^{\sharp}$
and $F$ are interderivable for every $F$.
\end{lemma}

\begin{proof}
Let $E$ be $s=t$.  The computation equations give $\operatorname{gate}(0,E)$,
both sides being $0$, and make $\operatorname{gate}(\Suc z,E)$ and $E$
interderivable.  From $E$: the matrix $\operatorname{gate}(v,E)$ is a
literal equation, so $\rulename{DI}_{\eta}$ asks for
$\operatorname{gate}(\Suc z,E)$ under the assumption
$\operatorname{gate}(z,E)$, which is not needed.  From
$\operatorname{pad}(E)$: \rulename{DE} opens
$\Q{v}\operatorname{gate}(v,E)$, and $\rulename{QE}_{\omega}$ with
$a\defeq 0$, $n\defeq 1$ and the second premise $\operatorname{gate}(0,E)$
gives $\operatorname{gate}(\omega(0,1),E)$.  \rulename{I0} and
$\rulename{Tr}_{2}$ give $\operatorname{gate}(1,E)$, hence $E$.
\end{proof}

\begin{theorem}[Embedding]
\label{kern:thm:rup}
If $\Gamma\proves A$ in $\Nat$, then
$\Gamma^{\uparrow}\proves_{\mathsf{R}}A^{\uparrow}$.
\end{theorem}

\begin{proof}
By induction on the derivation.  \rulename{Hyp}, \rulename{I0},
\rulename{IS} translate to themselves, and \rulename{Tr} is derivable in
$\mathsf{R}$ (\cref{kern:lem:rbasic}(i)).

\rulename{QE}.  The premises translate to $\Q{y}(B^{\uparrow})^{\sharp}$ and
$B^{\uparrow}(0)$.  \Cref{kern:lem:pad} gives $(B^{\uparrow})^{\sharp}(0)$,
\cref{kern:lem:rtransport}(iii) gives $(B^{\uparrow})^{\sharp}(t)$, and
\cref{kern:lem:pad} gives $B^{\uparrow}(t)$, which is $B(t)^{\uparrow}$.

\rulename{DE}.  The first premise translates to
$\D{x}{y}(B^{\uparrow})^{\sharp}$, and the hypothesis opened by
\rulename{DE} in $\mathsf{R}$ is $\Q{y}(B^{\uparrow})^{\sharp}(b,y)$, the
translation of $\Q{y}B(b,y)$.

\rulename{DI}.  The translated premise is
$\Gamma^{\uparrow},B^{\uparrow}(t,0)\proves B^{\uparrow}(t,\Suc z)$.  The
target has the compound matrix $(B^{\uparrow})^{\sharp}$, so
$\rulename{DI}_{\eta}$ asks for $(B^{\uparrow})^{\sharp}(t,\Suc z)$ under
$(B^{\uparrow})^{\sharp}(t,0)$; remove the padding of the assumption, cut
with the premise, and pad the conclusion.

\rulename{QI}.  The translated premise is
$\Gamma^{\uparrow},E(z)\proves E(\Suc z)$, and the target is
$\Q{y}\operatorname{pad}(E(y))$, with a compound matrix.  By
\cref{kern:lem:rbasic}(ii) it suffices to derive
$\operatorname{pad}(E(\Suc z'))$ from $\operatorname{pad}(E(0))$.  The
assumption gives $E(0)$; equation induction in $\mathsf{R}$
(\cref{kern:lem:rbasic}(iii)) with the translated premise as its step gives
$E(\Suc z')$; and \cref{kern:lem:pad} pads it.
\end{proof}

\subsection{The reading inside the core and the two round trips}
\label{kern:sec:Rdown}

The translation $F^{\downarrow}$ expresses the progressive reading with the
native connectives $\Rightarrow_*$, $\forall_*$, $\exists_*$ of
\cref{base:sec:structural}.  It fixes equations and is defined by
\begin{align*}
 (\Q{y}E)^{\downarrow} &\defeq
   \forall_* y\,\bigl(E(y)\Rightarrow_* E(\Suc y)\bigr),
 & (\Q{y}B)^{\downarrow} &\defeq \Q{y}B^{\downarrow},\\
 (\D{x}{y}E)^{\downarrow} &\defeq
   \exists_* x\,\forall_* y\,\bigl(E(x,y)\Rightarrow_* E(x,\Suc y)\bigr),
 & (\D{x}{y}B)^{\downarrow} &\defeq \D{x}{y}B^{\downarrow},
\end{align*}
where $E$ is a literal equation and $B$ is compound.  It maps compound
formulas to compound formulas and commutes with substitution.

\begin{theorem}[The reading is derivation-preserving]
\label{kern:thm:rdown}
If $\Gamma\proves_{\mathsf{R}}F$, then
$\Gamma^{\downarrow}\proves F^{\downarrow}$ in $\Nat$.
\end{theorem}

\begin{proof}
By \cref{base:thm:bridge} the ordinary reading of a native connective is
equivalent in $\Teq$ to the corresponding connective of the readings.  By
induction on $F$, therefore, $(F^{\downarrow})^{\circ}$ is equivalent in
$\Teq$ to $F^{\mathrm{p}}$.  \Cref{kern:thm:rsound} gives
$\Gamma^{\mathrm{p}}\proves_{\Teq}F^{\mathrm{p}}$, hence
$(\Gamma^{\downarrow})^{\circ}\proves_{\Teq}(F^{\downarrow})^{\circ}$, and
the faithfulness of the ordinary reading (\cref{base:thm:bridge}) gives
$\Gamma^{\downarrow}\proves F^{\downarrow}$.
\end{proof}

A formula of the form $A^{\uparrow}$ is an \emph{image}.  The rules of the
native connectives are available in $\mathsf{R}$ for images, over arbitrary
contexts of $\mathsf{R}$.

\begin{lemma}[Logic on images]
\label{kern:lem:image}
Let $\Delta$ be a context of $\mathsf{R}$, let $A$, $A_i$, $B$, $C$ be
native formulas, and let $F$ be a formula of $\mathsf{R}$.
\begin{enumerate}[label=\textup{(\roman*)}]
\item If $A_1,\dots,A_k\proves C$ in $\Nat$ and
 $\Delta\proves_{\mathsf{R}}A_i^{\uparrow}$ for every $i$, then
 $\Delta\proves_{\mathsf{R}}C^{\uparrow}$.
\item If $\Delta,C(0)^{\uparrow}\proves_{\mathsf{R}}C(\Suc z)^{\uparrow}$
 with $z$ fresh, then $\Delta\proves_{\mathsf{R}}(\Q{y}C(y))^{\uparrow}$.
\item If $\Delta,A^{\uparrow}\proves_{\mathsf{R}}B^{\uparrow}$, then
 $\Delta\proves_{\mathsf{R}}(A\Rightarrow_* B)^{\uparrow}$.
\item If $\Delta\proves_{\mathsf{R}}C(z)^{\uparrow}$ with $z$ not free in
 $\Delta$, then $\Delta\proves_{\mathsf{R}}(\forall_* y\,C(y))^{\uparrow}$.
\item If $\Delta\proves_{\mathsf{R}}(\exists_* x\,C(x))^{\uparrow}$ and
 $\Delta,C(b)^{\uparrow}\proves_{\mathsf{R}}F$ with $b$ not free in $\Delta$,
 $F$ and $C$, then $\Delta\proves_{\mathsf{R}}F$.
\end{enumerate}
\end{lemma}

\begin{proof}
(i) \Cref{kern:thm:rup}, weakening and cut.  Every rule of the native
connectives without a discharged assumption is of this form.

(ii) The target is $\Q{y}(C^{\uparrow})^{\sharp}(y)$, with a compound matrix.
By \cref{kern:lem:rbasic}(ii) it suffices to derive
$(C^{\uparrow})^{\sharp}(\Suc z)$ from $(C^{\uparrow})^{\sharp}(0)$, which
follows from the hypothesis and \cref{kern:lem:pad}.

(iii) $A\Rightarrow_* B$ is $\Q{y}\sel(y,A,B)$.  By (ii) it suffices to
derive $\sel(\Suc z,A,B)^{\uparrow}$ from $\sel(0,A,B)^{\uparrow}$.  By (i)
and $\sel(0,A,B)\proves A$, the assumption gives $A^{\uparrow}$; the
hypothesis gives $B^{\uparrow}$; and (i) with $B\proves\sel(\Suc z,A,B)$
gives the conclusion.

(iv) $\forall_* y\,C$ is $C(0)\wedge_*\Q{y}C$.  Substitution in
$\mathsf{R}$ gives $C(0)^{\uparrow}$ and $C(\Suc z)^{\uparrow}$, (ii) gives
$(\Q{y}C)^{\uparrow}$, and (i) with the introduction of $\wedge_*$ gives the
conclusion.

(v) $\exists_* x\,C$ is $\D{x}{y}\operatorname{gate}(y,C)$, and its image is
$\D{x}{y}\bigl(\operatorname{gate}(y,C)^{\uparrow}\bigr)^{\sharp}$.
\rulename{DE} in $\mathsf{R}$ with the fresh variable $b$ opens the image of
$\Q{y}\operatorname{gate}(y,C(b))$.  By (i) and
$\Q{y}\operatorname{gate}(y,C(b))\proves C(b)$, this gives
$C(b)^{\uparrow}$, and cut with the second hypothesis gives $F$.  The
formula $F$ is carried through \rulename{DE} and is not analyzed.
\end{proof}
% ed.: the package states this lemma with the proof "translate the derived
% rule constructions"; items (i)-(v) spell out which constructions are used
% and why they hold over an arbitrary context of the new calculus.

The lemma does not say that the selector and the connectives of $\Nat$ can
be applied to arbitrary formulas of $\mathsf{R}$; that would ignore the
dependence of $\Q{}$ on the root of its matrix.

\begin{lemma}[Normalization of atomic closures]
\label{kern:lem:atomicnorm}
For a literal equation $E$, over every context of $\mathsf{R}$,
\[
 \Q{y}E(y)\dashv\vdash_{\mathsf{R}}(\Q{y}E(y))^{\downarrow\uparrow},
 \qquad
 \D{x}{y}E(x,y)\dashv\vdash_{\mathsf{R}}(\D{x}{y}E(x,y))^{\downarrow\uparrow}.
\]
\end{lemma}

\begin{proof}
Equations are their own images.  From $\Q{y}E(y)$ and $E(z)$, one step
(\cref{kern:lem:rtransport}(i)) gives $E(\Suc z)$;
\cref{kern:lem:image}(iii) and (iv) give the image of
$\forall_* y\,(E(y)\Rightarrow_* E(\Suc y))$.  Conversely, from that image
and $E(z)$, \cref{kern:lem:image}(i) with the eliminations of $\forall_*$
and $\Rightarrow_*$ gives $E(\Suc z)$, and $\rulename{QI}_{\eta}$ at the
literal equation $E$ gives $\Q{y}E(y)$.

For $\D{}{}$: apply \rulename{DE} at a fresh $b$, convert the opened
$\Q{y}E(b,y)$ by the first part, and apply \cref{kern:lem:image}(i) with the
introduction of $\exists_*$.  Conversely apply \cref{kern:lem:image}(v) with
$F\defeq\D{x}{y}E(x,y)$: under the image of
$\forall_* y\,(E(b,y)\Rightarrow_* E(b,\Suc y))$ and the assumption $E(b,z)$
the eliminations give $E(b,\Suc z)$, which is the premise of
$\rulename{DI}_{\eta}$ with the witness $b$.
\end{proof}

\begin{theorem}[Round trips]
\label{kern:thm:rround}
For every native formula $A$ and every formula $F$ of $\mathsf{R}$,
\[
 A\dashv\vdash A^{\uparrow\downarrow}\ \text{in }\Nat,
 \qquad
 F\dashv\vdash_{\mathsf{R}}F^{\downarrow\uparrow},
\]
over arbitrary contexts.
\end{theorem}

\begin{proof}
\emph{In $\Nat$.}  First, $\operatorname{pad}(E)^{\downarrow}$ is
interderivable with $E$: it is
$\exists_* r\,\forall_* v\,\bigl(\operatorname{gate}(v,E)\Rightarrow_*
\operatorname{gate}(\Suc v,E)\bigr)$ with $r$ unused.  From $E$ every
$\operatorname{gate}(\Suc v,E)$ follows.  Conversely
$\operatorname{gate}(0,E)$ is derivable, and progress at $0$ gives
$\operatorname{gate}(1,E)$, hence $E$.  Now induct on $A$.  Equations are
fixed.  The matrix $(B^{\uparrow})^{\sharp}$ is compound, so
$(\Q{y}B)^{\uparrow\downarrow}$ is $\Q{y}((B^{\uparrow})^{\sharp})^{\downarrow}$.
If $B$ is an equation, the matrix is $\operatorname{pad}(B)^{\downarrow}$,
interderivable with $B$.  If $B$ is compound, the matrix is
$B^{\uparrow\downarrow}$, interderivable with $B$ by the induction
hypothesis.  The congruence of $\Q{}$ in $\Nat$
(\cref{base:sec:structural}) gives the claim, and $\D{}{}$ is the same.

\emph{In $\mathsf{R}$.}  Induct on $F$.  Equations are fixed.  If the matrix
of $F$ is a literal equation, the claim is \cref{kern:lem:atomicnorm}.  Let
the matrix $B$ be compound.  The induction hypothesis makes $B$ and
$B^{\downarrow\uparrow}$ interderivable, and both are compound.
\Cref{kern:lem:rtransport}(iv) makes $\Q{y}B$ and
$\Q{y}B^{\downarrow\uparrow}$ interderivable, and since
$B^{\downarrow\uparrow}$ is compound, $\Q{y}B^{\downarrow\uparrow}$ is
$(\Q{y}B)^{\downarrow\uparrow}$ with no padding inserted.  $\D{}{}$ is the
same.  No congruence across different kinds of root is used, and the
normalization is proved for every formula before the connectives are
extended to all of $\mathsf{R}$.
\end{proof}

\begin{theorem}[The orbit calculus and the core]
\label{kern:thm:uniform}
\begin{enumerate}[label=\textup{(\roman*)}]
\item $\Gamma\proves A$ in $\Nat$ if and only if
 $\Gamma^{\uparrow}\proves_{\mathsf{R}}A^{\uparrow}$.
\item $\Gamma\proves_{\mathsf{R}}F$ if and only if
 $\Gamma^{\downarrow}\proves F^{\downarrow}$ in $\Nat$.
\item $\mathsf{R}$ derives $0=\Suc 0$ if and only if $\Nat$ does.
\item With the progressive reading, $\mathsf{R}$ is bi-interpretable with
 $\Teq$ and with $\iPRA$, with whole domains, actual equality, identity
 numerical composites and both formula round trips.
\end{enumerate}
\end{theorem}

\begin{proof}
(i) One direction is \cref{kern:thm:rup}.  For the other, apply
\cref{kern:thm:rdown} to $\Gamma^{\uparrow}\proves_{\mathsf{R}}A^{\uparrow}$
and remove $\uparrow\downarrow$ by \cref{kern:thm:rround} and cut.  (ii) is
symmetric, with \cref{kern:thm:rdown}, \cref{kern:thm:rup} and the second
round trip.  (iii) is (i) or (ii) for the equation $0=\Suc 0$, which both
translations fix.  (iv) The translations $\uparrow$ and $\downarrow$ fix
every term, so composing them with the interpretations of
\cref{base:thm:bridge} adds no change of domain; the formula round trips
compose.
\end{proof}

\begin{corollary}[Connectives for all formulas of the orbit calculus]
\label{kern:cor:rlogic}
For formulas $F$, $G$ of $\mathsf{R}$ put
$F\Rightarrow_{\mathsf{R}}G\defeq
(F^{\downarrow}\Rightarrow_* G^{\downarrow})^{\uparrow}$, and likewise for
$\wedge$, $\vee$, $\exists$ and $\forall$; falsity is $0=\Suc 0$.  These
operations obey the rules of intuitionistic natural deduction in
$\mathsf{R}$, and $\mathsf{R}$ derives coverage and equation induction.
\end{corollary}

\begin{proof}
By \cref{kern:thm:rround}, $(F\Rightarrow_{\mathsf{R}}G)^{\downarrow}$ is
interderivable in $\Nat$ with $F^{\downarrow}\Rightarrow_* G^{\downarrow}$.
By \cref{kern:thm:uniform}(ii) each rule for $\Rightarrow_{\mathsf{R}}$ is
therefore the rule for $\Rightarrow_*$ in $\Nat$, applied to the readings.
The other connectives are the same.
\end{proof}

\begin{remark}[Sources and Lean]
\label{kern:rem:uniform-sources}
The calculus and \cref{kern:thm:uniform} are from \pkg{uniform\_progress}.
The package compares $\mathsf{R}$ with an earlier progressive presentation
that reads every closure progressively, has a recursor whose body sees the
counter, two eliminations and a rule of Boolean cases; here the iterator is
unchanged, compound closures keep the conditional reading, and one
elimination serves both.  In Lean (modules \Lean{NativeProgressive} and
following) the calculus is \Lean{RPrf}, the translation $\downarrow$ is
\Lean{NFm.rd} and $\uparrow$ is \Lean{NFm.emb}.  \Cref{kern:thm:rdown} is
\Lean{RPrf.toNPrf}, with \Lean{rprf\_sound\_pra} and
\Lean{rprf\_consistent}; \cref{kern:lem:rstruct} is \Lean{RPrf.subst};
\cref{kern:lem:rbasic}(iii) is \Lean{RPrf.eqInd};
\cref{kern:lem:rtransport}(ii) is \Lean{RPrf.leibniz}; \cref{kern:lem:pad}
is \Lean{RPrf.padI} and \Lean{RPrf.padE}; \cref{kern:thm:rup} is
\Lean{NPrf.toR}; \cref{kern:thm:uniform}(i) and (ii) are
\Lean{rprf\_emb\_iff} and \Lean{rprf\_iff\_rd}; and the two round trips are
\Lean{NEquiv.rd\_emb} and \Lean{REquivIn.emb\_rd}.  In the Lean proof of
\cref{kern:thm:rdown} the only rule that needs an argument is
$\rulename{QE}_{\omega}$ at an equation, which is an equation induction in
the kernel (\Lean{NPrf.progE}).
\end{remark}

\subsection{Transport in a formula template}
\label{kern:sec:mixed}

\begin{definition}[The calculus $\mathsf{R}_{F}$]
\label{kern:def:RF}
$\mathsf{R}_{F}$ is $\mathsf{R}$ with $\rulename{Tr}_{2}$ replaced by
$\rulename{Tr}_{F}$; its seven rules are \rulename{Hyp}, \rulename{I0},
\rulename{IS}, $\rulename{Tr}_{F}$, $\rulename{QE}_{\omega}$, \rulename{DE}
and $\rulename{DI}_{\eta}$.
\end{definition}

\begin{proposition}[Mixed closure]
\label{kern:prop:mixed}
$\mathsf{R}_{F}$ and $\mathsf{R}$ derive the same sequents.  Hence
\cref{kern:thm:rsound,kern:thm:uniform} hold for $\mathsf{R}_{F}$.
\end{proposition}

\begin{proof}
$\rulename{Tr}_{2}$ is the instance of $\rulename{Tr}_{F}$ at an equation
template, and $\rulename{Tr}_{F}$ is derivable in $\mathsf{R}$ by
\cref{kern:lem:rtransport}(ii).
\end{proof}

\begin{remark}[The account of \pkg{scope\_and\_closure}]
\label{kern:rem:mixed-sources}
$\mathsf{R}_{F}$ is the ``mixed closure'' calculus of
\pkg{scope\_and\_closure}, written there with the orbit term
$\It{n}{t}{u}{\Suc u}$ and the selector $d_B$ for $\eta_B$.  The package
proves the bi-interpretation with $\Teq$ independently of
\pkg{uniform\_progress}, and its proof differs in three places.  The gate is
extended through binders, $G_v(\Q{y}B)\defeq\Q{y}G_vB$ and
$G_v(\D{x}{y}B)\defeq\D{x}{y}G_vB$, with $\proves G_0B$ and
$G_{\Suc v}B\dashv\vdash B$ proved by induction on $B$ from homogeneous
congruence.  Every matrix is padded, $P(B)\defeq\D{r}{v}G_vB$, compound or
not.  And the second formula round trip is proved through $\Teq$: with
$\rho$ the progressive reading and $K(\varphi)$ the padded native
translation of an ordinary formula $\varphi$, the package shows
$\rho(K\varphi)\leftrightarrow\varphi$ in $\Teq$ and
$B\dashv\vdash K(B^{\rho})$ in its calculus, the latter by the cases of
\cref{kern:lem:atomicnorm,kern:thm:rround}.  For $\rulename{QE}_{\omega}$
it needs only $\omega(a,1)=\Suc a$ inside the calculus, as here.  In Lean
the calculus is \Lean{MCPrf}, with \Lean{mcprf\_iff\_rprf},
\Lean{mcprf\_iff\_rd}, \Lean{mcprf\_emb\_iff} and
\Lean{mcprf\_consistent}.  The package measures no checker for it and
states that the single introduction moves a cost into the elimination, which
builds a term and takes a second numerical parameter.
\end{remark}

\subsection{Elimination at a target}
\label{kern:sec:target}

The orbit elimination takes two terms, a source and a seed.  It can be
replaced by an elimination that takes the target alone.  For a matrix $B$
and a term $t$ put
\[
 \delta_B(t)\defeq
 \begin{cases}
  r, & B\text{ is a literal equation and }t\equiv\Suc r,\\
  0, & \text{otherwise.}
 \end{cases}
\]
The test is on the syntax of $t$; terms are not normalized before it, and
$J(\Suc r)=\It{0}{\Suc r}{u}{0}$ is not a literal successor.

\begin{definition}[The calculus $\mathsf{T}$]
\label{kern:def:T}
$\mathsf{T}$ has the seven rules \rulename{Hyp}, \rulename{I0},
\rulename{IS}, $\rulename{Tr}_{2}$, \rulename{DE}, $\rulename{DI}_{\eta}$ and
\[
 \infrule{\Gamma\proves\Q{y}B(y)\qquad\Gamma\proves B(\delta_B(t))}
         {\Gamma\proves B(t)}\ \rulename{QE}_{\delta}.
\]
\end{definition}

A target that is a literal successor asks for the instance at its
predecessor; every other target of an atomic closure, and every target of a
compound closure, asks for the instance at zero.

\begin{lemma}[Interfaces and induction in $\mathsf{T}$]
\label{kern:lem:tbasic}
$\mathsf{T}$ derives the equality rules and $\rulename{QI}_{\eta}$, and for a
literal equation $E$, a compound matrix $B$ and every term $t$:
\begin{enumerate}[label=\textup{(\roman*)}]
\item one step: from $\Q{y}E(y)$ and $E(t)$ to $E(\Suc t)$;
\item instance from the base: from $\Q{y}E(y)$ and $E(0)$ to $E(t)$, and
 from $\Q{y}B(y)$ and $B(0)$ to $B(t)$;
\item equation induction.
\end{enumerate}
\end{lemma}

\begin{proof}
The equality rules and $\rulename{QI}_{\eta}$ are derived as in
\cref{kern:lem:rbasic}.  (i) $\rulename{QE}_{\delta}$ at the target
$\Suc t$.  (ii) For $B$ compound the rule is $\rulename{QE}_{\delta}$
itself.  For $E$, take the target $J(t)$, which is not a literal successor:
$\rulename{QE}_{\delta}$ gives $E(J(t))$ from $E(0)$, and $J(t)=t$ with
$\rulename{Tr}_{2}$ gives $E(t)$.  (iii) $\rulename{QI}_{\eta}$ and (ii).
\end{proof}

This bootstrap uses no predecessor, no numerical case distinction and no
identity for the orbit term.

\begin{lemma}[Substitution in $\mathsf{T}$]
\label{kern:lem:tsubst}
If $\Gamma\proves_{\mathsf{T}}A$, then
$\Gamma[\theta]\proves_{\mathsf{T}}A[\theta]$.  Injective renaming maps rule
instances to rule instances; substitution in general does not.
\end{lemma}

\begin{proof}
By induction on the derivation, with eigenvariables renamed.  All rules but
$\rulename{QE}_{\delta}$ commute with substitution.  Consider
$\rulename{QE}_{\delta}$ with the matrix $B$ and the target $t$.  If $B$ is
compound, or $t\equiv\Suc r$, or neither $t$ nor $t[\theta]$ is a literal
successor, the substituted inference is an instance.  In the remaining case
$B$ is a literal equation, the second premise is the instance at $0$, and
$t[\theta]$ is a literal successor; then \cref{kern:lem:tbasic}(ii) gives
$B[\theta](t[\theta])$ from the substituted premises.  That lemma does not
depend on the present one.
\end{proof}

\begin{lemma}[Transport in $\mathsf{T}$]
\label{kern:lem:ttransport}
The statements (ii) and (iv) of \cref{kern:lem:rtransport} hold in
$\mathsf{T}$.
\end{lemma}

\begin{proof}
The proof of \cref{kern:lem:rtransport} uses $\rulename{Tr}_{2}$,
$\rulename{QI}_{\eta}$, \rulename{DE}, $\rulename{DI}_{\eta}$, one step at a
literal equation, the instance from the base at a compound matrix, and
substitution.  In $\mathsf{T}$ these are \cref{kern:lem:tbasic,kern:lem:tsubst},
and the instance from the base at a compound matrix is primitive, so no
redex has to be removed.
\end{proof}

\begin{theorem}[Target elimination and orbit elimination]
\label{kern:thm:target}
$\mathsf{T}$ and $\mathsf{R}$ derive the same sequents.  Hence
\cref{kern:thm:rsound,kern:thm:uniform} hold for $\mathsf{T}$.
\end{theorem}

\begin{proof}
The two calculi differ in the elimination of $\Q{}$ only.

\emph{$\rulename{QE}_{\delta}$ in $\mathsf{R}$.}  At a literal equation and
a target $\Suc r$ it is one step (\cref{kern:lem:rtransport}(i)).  In the
other cases its second premise is the instance at $0$, and it is
\cref{kern:lem:rtransport}(iii); for a compound matrix this uses transport
in a compound formula of $\mathsf{R}$, which is why that lemma precedes the
comparison.

\emph{$\rulename{QE}_{\omega}$ in $\mathsf{T}$.}  At a compound matrix,
$\rulename{QE}_{\delta}$ with the target $\omega(0,n)$ gives the conclusion
from $B(0)$ at once.  At a literal equation, let $\Q{y}E(y)$ and $E(a)$ be
given.  Equation induction on $k$ in $\mathsf{T}$ gives $E(\omega(k,a))$:
the base is $E(a)$ with $\omega(0,a)=a$, and the step is one step at
$\omega(k,a)$ with $\omega(\Suc k,a)=\Suc\,\omega(k,a)$.  At $k\defeq n$,
the equation $\omega(n,a)=\omega(a,n)$, which holds in $\mathsf{T}$ by
\cref{kern:lem:orbit}, gives $E(\omega(a,n))$.
\end{proof}

The orbit identities are needed to simulate the orbit rule in $\mathsf{T}$.
They are not needed by the rule $\rulename{QE}_{\delta}$, by its soundness
or by the induction of \cref{kern:lem:tbasic}.  Directly, $\rulename{QE}_{\delta}$
is sound for the progressive reading as follows: at a literal successor it
uses one instance of progress, at another target of an atomic closure it
uses equation induction from the instance at zero, and at a compound matrix
it is the elimination of the conditional reading.

\begin{remark}[Sources and Lean]
\label{kern:rem:target-sources}
The subsection is the first part of \pkg{target\_closure}.  In Lean the
calculus is \Lean{TQPrf}, with \Lean{tqprf\_iff\_rprf}, the orbit rule
derived as \Lean{TQPrf.orbit}, and the consequences
\Lean{tqprf\_emb\_iff}, \Lean{tqprf\_iff\_rd} and \Lean{tqprf\_consistent}.
The second part of the package concerns the typing of certificates for the
calculus with the conditional reading and is summarized in
\cref{kern:sec:sizes}.
\end{remark}

\subsection{The calculi of the chapter}
\label{kern:sec:table}

\Cref{kern:tab:calculi} lists the calculi of
\cref{kern:sec:intro,kern:sec:transport,kern:sec:aqe,kern:sec:uniq,kern:sec:prog}.  Each
has \rulename{Hyp}, \rulename{I0}, \rulename{IS} and \rulename{DE}; the
second column gives its other rules.  The calculus
$\mathsf{U}_{\mathrm{c}}$ is the exception: it is stated on normal terms and
has neither \rulename{I0} nor \rulename{IS}.  The column ``subst.'' says
whether substitution maps rule instances to rule instances
(\cref{kern:lem:struct}); for $\mathsf{U}_{\mathrm{c}}$ substitution is
followed by normalization (\cref{kern:lem:ucstruct}).

\begin{longtable}{@{}lp{0.36\linewidth}cp{0.2\linewidth}c@{}}
\caption{The calculi of the chapter.}\label{kern:tab:calculi}\\
\toprule
Calculus & Rules besides \rulename{Hyp}, \rulename{I0}, \rulename{IS},
\rulename{DE} & Rules & Relation to $\Nat$ & Subst.\\
\midrule
\endfirsthead
\toprule
Calculus & Rules besides \rulename{Hyp}, \rulename{I0}, \rulename{IS},
\rulename{DE} & Rules & Relation to $\Nat$ & Subst.\\
\midrule
\endhead
\bottomrule
\endfoot
$\Nat$ & \rulename{Tr}, \rulename{QE}, \rulename{DI}, \rulename{QI} & 8 &
 & yes\\
$\Nat_{\mathrm{D}}$ & \rulename{Tr}, \rulename{QE}, \rulename{DI},
 \rulename{DP} & 8 & equivalent, \cref{kern:thm:direct} & yes\\
$\Nat_{\det}$ & \rulename{Tr}, \rulename{QE}, \rulename{DI},
 $\rulename{QI}_{\eta}$ & 8 & equivalent, \cref{kern:cor:det} & yes\\
$\mathsf{W}$ & \rulename{Tr}, \rulename{QE}, $\rulename{DI}_{\mathrm{w}}$ &
 7 & equivalent, \cref{kern:thm:witness} & no\\
$\Nat_{2}$ & $\rulename{Tr}_{2}$, \rulename{QE}, \rulename{DI},
 \rulename{QI} & 8 & equivalent, \cref{kern:thm:onecontext} & yes\\
$\Nat_{F}$ & $\rulename{Tr}_{F}$, \rulename{QE}, \rulename{DI},
 \rulename{QI} & 8 & equivalent, \cref{kern:thm:onecontext} & yes\\
$\mathsf{O}$ & $\rulename{Tr}^{-}$, \rulename{QE}, \rulename{DI},
 \rulename{QI} & 8 & equivalent, \cref{kern:thm:onesided} & no\\
$\mathsf{M}$ & $\rulename{Tr}^{-}$, \rulename{QE},
 $\rulename{DI}_{\mathrm{w}}$ & 7 & equivalent, \cref{kern:thm:onesided} &
 no\\
$\mathsf{X}$ & \rulename{TrX}, \rulename{QX}, \rulename{DI}, \rulename{QI} &
 8 & equivalent, \cref{kern:thm:reversal} & no\\
$\mathsf{Z}$ & $\rulename{Tr}_{F}^{-}$, \rulename{QE}, \rulename{DI},
 \rulename{QI} & 8 & equivalent, \cref{kern:thm:dispatch} & no\\
$\mathsf{G}$ & \rulename{Rep}, \rulename{QE}, \rulename{DI}, \rulename{QI} &
 8 & equivalent, \cref{kern:thm:tfree} & no\\
$\mathsf{G}_7$ & \rulename{Rep}, \rulename{QE},
 $\rulename{DI}_{\mathrm{w}}$ & 7 & equivalent, \cref{kern:thm:tfree} & no\\
$\mathsf{B}$ & \rulename{BRep}, \rulename{QE}, \rulename{DI}, \rulename{QI}
 & 8 & equivalent, \cref{kern:thm:brep} & no\\
$\mathsf{B}_{7}$ & \rulename{BRep}, \rulename{QE},
 $\rulename{DI}_{\mathrm{w}}$ & 7 & equivalent, \cref{kern:thm:brep} & no\\
$\mathsf{B}^{-}$ & $\rulename{BRep}^{-}$, \rulename{QE}, \rulename{DI},
 \rulename{QI} & 8 & equivalent, \cref{kern:thm:bminus} & no\\
$\mathsf{B}^{-}_{7}$ & $\rulename{BRep}^{-}$, \rulename{QE},
 $\rulename{DI}_{\mathrm{w}}$ & 7 & equivalent, \cref{kern:thm:bminus} & no\\
$\mathsf{L}$ & \rulename{LOR}, \rulename{QE}, \rulename{DI}, \rulename{QI} &
 8 & equivalent, \cref{kern:thm:onehole} & yes\\
$\Nat_{\mathrm{u}}$ & uncrossed replacement in a term template,
 \rulename{QE}, \rulename{DI}, \rulename{QI} & 8 & equivalent,
 \cref{kern:thm:onehole} & yes\\
$\mathsf{A}_{7}$ & \rulename{BRep}, \rulename{AQE},
 $\rulename{DI}_{\mathrm{w}}$ & 7 & equivalent, \cref{kern:thm:aqe} & no\\
$\mathsf{U}$ & \rulename{Tr}, \rulename{QE}, \rulename{DI}, \rulename{UR} &
 8 & equivalent, \cref{kern:thm:urnative} & yes\\
$\mathsf{U}_{\mathrm{o}}$ & \rulename{Tr}, \rulename{QE}, \rulename{DI},
 \rulename{Orb} & 8 & equivalent, \cref{kern:prop:orbit} & yes\\
$\mathsf{U}_{\mathrm{c}}$ & \rulename{Ref}, \rulename{QE}, \rulename{DI},
 \rulename{UX}, on normal terms; no \rulename{I0}, \rulename{IS} & 6 &
 through $\nf$, \cref{kern:thm:crossed} & yes\\
$\mathsf{R}$ & $\rulename{Tr}_{2}$, $\rulename{QE}_{\omega}$,
 $\rulename{DI}_{\eta}$ & 7 & through $\uparrow$ and $\downarrow$,
 \cref{kern:thm:uniform} & yes\\
$\mathsf{R}_{F}$ & $\rulename{Tr}_{F}$, $\rulename{QE}_{\omega}$,
 $\rulename{DI}_{\eta}$ & 7 & same sequents as $\mathsf{R}$,
 \cref{kern:prop:mixed} & yes\\
$\mathsf{T}$ & $\rulename{Tr}_{2}$, $\rulename{QE}_{\delta}$,
 $\rulename{DI}_{\eta}$ & 7 & same sequents as $\mathsf{R}$,
 \cref{kern:thm:target} & no\\
\end{longtable}

% ---------------------------------------------------------------------
\section{Two results on other recursors}
\label{kern:sec:recursors}

This section draws on the appendix of \pkg{witness\_rule} on counter
congruence and on the obstruction note of \pkg{scope\_and\_closure}.  Both
concern term languages other than that of \cref{base:sec:terms}; the
formulas and the logical rules are the native ones.

\subsection{Counter congruence at a zero seed}
\label{kern:sec:zcc}

The appendix of \pkg{witness\_rule} refines an earlier package,
\pkg{counter\_congruence}, whose term language has a recursor that sees the
counter.  The terms are variables, $0$, $\Suc t$ and
$\operatorname{Rec}(n,a;\,u,v.\,h)$, which binds $u$, the previous counter,
and $v$, the previous state, in $h$.  The computation equations are
\[
 \operatorname{Rec}(0,a;\,u,v.\,h)=\Suc a,\qquad
 \operatorname{Rec}(\Suc n,a;\,u,v.\,h)
  =h\subst{n}{u}\subst{\operatorname{Rec}(n,a;\,u,v.\,h)}{v}.
\]
The formulas are equations and $\D{x}{y}A$; a closure $\Q{y}B$ occurs in a
context only, as a resource in the sense of \cref{base:sec:clauses}.

\begin{definition}[The calculus $\mathcal{C}_0$]
\label{kern:def:C0}
The rules of $\mathcal{C}_0$ are: \rulename{Hyp} for a formula of the
context; lookup of a closure of the context; the two computation equations;
\rulename{QE}, with its first premise a lookup; \rulename{DE};
\rulename{DI}; \rulename{DP}; root replacement on the left,
\[
 \infrule{\Gamma\proves s=t\qquad\Gamma\proves s=r}{\Gamma\proves t=r}\
 \rulename{RL};
\]
and counter congruence at a zero seed,
\[
 \infrule{\Gamma\proves s=t}
  {\Gamma\proves\operatorname{Rec}(s,0;\,u,v.\,h)
               =\operatorname{Rec}(t,0;\,u,v.\,h)}\ \rulename{ZCC}
 \qquad(v\text{ not free in }h).
\]
\end{definition}

The seed of \rulename{ZCC} is the literal $0$, and its body does not depend
on the state, also not below nested binders.  There is no reflexivity, no
general congruence and no transport.  The reference calculus of
\pkg{counter\_congruence} has, in place of \rulename{ZCC}, the rule
\rulename{CC} with an arbitrary seed $a$ and an arbitrary body.  Count one
node for every rule instance, lookups included.

\begin{lemma}[Root equality and cases]
\label{kern:lem:zccbasic}
In $\mathcal{C}_0$, without \rulename{ZCC}:
\begin{enumerate}[label=\textup{(\roman*)}]
\item reflexivity is derivable with three nodes, and symmetry and
 transitivity are derivable;
\item from $l=r$, $l=f$ and $r=g$, two instances of \rulename{RL} give
 $f=g$;
\item numerical cases, from $A(0)$ and $A(\Suc z)$ to $A(t)$, are derivable
 with four nodes besides the two premises, with no equality rule and no
 \rulename{DP}.
\end{enumerate}
\end{lemma}

\begin{proof}
(i) $\operatorname{Rec}(\Suc s,0;\,u,v.\,u)=s$ is a computation equation, and
\rulename{RL} applied to two copies of it gives $s=s$.  \rulename{RL}
applied to $s=t$ and $s=s$ gives $t=s$, and applied to $t=s$ and $t=r$ it
gives $s=r$.  (ii) \rulename{RL} on $l=r$ and $l=f$ gives $r=f$, and on
$r=f$ and $r=g$ it gives $f=g$.  (iii) \rulename{DI} with an unused
coordinate and the witness $0$, \rulename{DE}, a lookup of the opened
closure, and \rulename{QE}.
\end{proof}

\begin{theorem}[Congruence from \rulename{ZCC}]
\label{kern:thm:zcc}
For every term template $F(x)$, $\mathcal{C}_0$ derives $F(s)=F(t)$ from
$s=t$, with exactly $20$ nodes besides the derivation of $s=t$ and with no
instance of \rulename{DP}.  Consequently $\mathcal{C}_0$ derives transport
in every formula template, and it derives the same sequents as the calculus
with \rulename{CC} and as the calculus with unrestricted transport.
\end{theorem}

\begin{proof}
Put $U(n)\defeq\operatorname{Rec}(n,0;\,u,v.\,\Suc\Suc u)$.  Its computation
equations are $U(0)=\Suc 0$ and $U(\Suc z)=\Suc\Suc z$, so numerical cases
give $U(n)=\Suc n$ for every term $n$, with six nodes.  \rulename{ZCC}
applies to $U$ and gives $U(s)=U(t)$; with $U(s)=\Suc s$ and
$U(t)=\Suc t$, \cref{kern:lem:zccbasic}(ii) gives $\Suc s=\Suc t$.  This
stage has $1+6+6+2=15$ nodes.

Put $V_F(n)\defeq\operatorname{Rec}(\Suc n,0;\,u,v.\,F(u))$.  One
computation equation gives $V_F(n)=F(n)$.  \rulename{ZCC}, applied to
$\Suc s=\Suc t$ with the body $F(u)$, gives $V_F(s)=V_F(t)$, and
\cref{kern:lem:zccbasic}(ii) with $V_F(s)=F(s)$ and $V_F(t)=F(t)$ gives
$F(s)=F(t)$.  This stage has $1+1+1+2=5$ nodes.

Transport in a formula template $A(w)$ follows by induction on the number of
binders of $A$.  For an equation $p(w)=q(w)$, join $p(s)=p(t)$ and
$q(s)=q(t)$ with $p(s)=q(s)$ by \cref{kern:lem:zccbasic}.  For
$\D{x}{y}B(w,x,y)$, apply \rulename{DE} at a fresh $b$ and \rulename{DI} at
the witness $b$: the assumption $B(t,b,0)$ is transported back to
$B(s,b,0)$, \rulename{QE} with the opened closure gives $B(s,b,\Suc z)$, and
transport forward gives $B(t,b,\Suc z)$.

Every instance of \rulename{ZCC} is an instance of \rulename{CC}, and
\rulename{CC} is the congruence just proved for the template
$\operatorname{Rec}(x,a;\,u,v.\,h)$.  In the calculus with transport,
\rulename{RL} is transport at the left root and \rulename{ZCC} is transport
of one occurrence in an instance of reflexivity.
\end{proof}

The construction of \pkg{counter\_congruence} has $56$ nodes for the same
congruence; replacing its final assembly by \cref{kern:lem:zccbasic}(ii)
gives $44$, and the representation through $U$ and $V_F$ gives $20$.  These
are sizes of the constructions described, not lower bounds, and the terms in
the nodes grow with $F$.

\begin{proposition}[A wrapped rule]
\label{kern:prop:zccplus}
Let $\rulename{ZCC}^{+}$ infer
$\operatorname{Rec}(\Suc s,0;\,u,v.\,h)=\operatorname{Rec}(\Suc t,0;\,u,v.\,h)$
from $s=t$, with $v$ not free in $h$.  Replacing \rulename{ZCC} by
$\rulename{ZCC}^{+}$ gives the same derivable sequents, and congruence in a
term template then has $5$ nodes besides the derivation of $s=t$.
\end{proposition}

\begin{proof}
$\rulename{ZCC}^{+}$ with the body $F(u)$ gives $V_F(s)=V_F(t)$ directly;
two computation equations and two instances of \rulename{RL} finish.
Conversely \cref{kern:thm:zcc} derives $\rulename{ZCC}^{+}$, as congruence
for the template $\operatorname{Rec}(\Suc x,0;\,u,v.\,h)$.
\end{proof}

$\rulename{ZCC}^{+}$ wraps arbitrary premises in successors.  It is not
\rulename{ZCC} with a side condition on its premise, and the difference
matters:

\begin{proposition}[Successor premises do not suffice]
\label{kern:prop:zcccounter}
Restrict \rulename{ZCC} to premises of the form $\Suc p=\Suc q$, and keep
the other rules of $\mathcal{C}_0$.  The resulting calculus does not derive
$0=\Suc 0\proves\Suc 0=\Suc\Suc 0$, which $\mathcal{C}_0$ derives.
\end{proposition}

\begin{proof}
Interpret terms in the standard natural numbers, with the recursor read as
primitive recursion with the seed $a+1$.  Interpret an equation $s=t$ as
$e(s,t)$, where $e(m,n)$ holds if $m=n$ or both $m$ and $n$ are in
$\{0,1\}$; interpret closures and $\D{}{}$ by their ordinary readings over
this interpretation of equations.  The relation $e$ is an equivalence, so
\rulename{RL} preserves truth.  Computation equations are true because their
two sides have equal values.  For the restricted rule the two counters are
positive, and on positive numbers $e$ is equality, so a true premise gives
equal counters and equal values of the recursor.  Lookup, \rulename{QE} and
\rulename{DE} preserve truth by the readings, \rulename{DI} because every
number is zero or a successor, and \rulename{DP} by induction on the
natural numbers for the interpreted predicate.  The assumption $0=\Suc 0$ is
true, since $e(0,1)$ holds, and $\Suc 0=\Suc\Suc 0$ is false, since $e(1,2)$
fails.  $\mathcal{C}_0$ derives the sequent by congruence for $\Suc$.
\end{proof}

The interpretation is an interpretation of the rules; it does not read the
equation sign as equality and is not a model of the ordinary theory.  The
package assigns no checker size to $\mathcal{C}_0$ or to the wrapped rule.

\subsection{Recursors seeded at zero}
\label{kern:sec:zerostart}

The iterator has three arguments: a counter, a seed and a body.  The note of
\pkg{scope\_and\_closure} examines what is lost if the seed is fixed to
zero.  Consider two term languages over variables, $0$ and $\Suc$.  The
\emph{scalar} language has $\operatorname{It}_0(n;\,u.\,h)$ with
\[
 \operatorname{It}_0(0;\,u.\,h)=0,\qquad
 \operatorname{It}_0(\Suc n;\,u.\,h)=h\subst{\operatorname{It}_0(n;\,u.\,h)}{u}.
\]
The \emph{counter-aware} language has $\operatorname{Rec}_0(n;\,k,u.\,h)$
with
\[
 \operatorname{Rec}_0(0;\,k,u.\,h)=0,\qquad
 \operatorname{Rec}_0(\Suc n;\,k,u.\,h)
  =h\subst{n}{k}\subst{\operatorname{Rec}_0(n;\,k,u.\,h)}{u}.
\]
All statements below are about the denotations of terms in the standard
natural numbers.

\begin{theorem}[Monotonicity]
\label{kern:thm:mono}
Every term of either language denotes a function that is nondecreasing in
each of its free variables.
\end{theorem}

\begin{proof}
By induction on the term.  Variables, $0$ and $\Suc$ are clear.  For the
recursor, let the counter $n$ and the body $h$ denote nondecreasing
functions, fix an assignment $\rho$ of the parameters, and let $a^{\rho}_j$
be the state after $j$ steps: $a^{\rho}_0=0$ and
$a^{\rho}_{j+1}=h(j,a^{\rho}_j,\rho)$.  The sequence is nondecreasing in
$j$: $a^{\rho}_0=0\le a^{\rho}_1$, and $a^{\rho}_j\le a^{\rho}_{j+1}$ gives
$a^{\rho}_{j+1}=h(j,a^{\rho}_j,\rho)\le h(j+1,a^{\rho}_{j+1},\rho)
=a^{\rho}_{j+2}$.  For $\rho\le\sigma$, induction on $j$ gives
$a^{\rho}_j\le a^{\sigma}_j$.  Since $n(\rho)\le n(\sigma)$,
\[
 a^{\rho}_{n(\rho)}\le a^{\sigma}_{n(\rho)}\le a^{\sigma}_{n(\sigma)} .
\]
The scalar case is the same without the argument $j$ of $h$.  The argument
covers a variable that occurs both in the counter and in the body.
\end{proof}

\begin{theorem}[Zero sets in the scalar language]
\label{kern:thm:zeroprofile}
For a term $t(x_1,\dots,x_r)$ of the scalar language let
$\mathcal{Z}(t)\defeq\{\vec x: t(\vec x)=0\}$.  Then $\mathcal{Z}(t)$ is the
whole of $\mathbb{N}^r$ or a union $\bigcup_{j\in J}\{\vec x:x_j=0\}$ of
coordinate hyperplanes, for some $J\subseteq\{1,\dots,r\}$, possibly empty.
Every set of this form is the zero set of a term.
\end{theorem}

\begin{proof}
By induction on $t$.  The zero set of $0$ is everything, that of a variable
is its hyperplane, and that of a successor is empty.  For
$\operatorname{It}_0(n;\,u.\,h)$ fix the parameters and let $b$ be the value
of $h$ at $u=0$.  If $b=0$, every iterate from $0$ is $0$.  If $b>0$,
monotonicity of $h$ gives $h(v)\ge b>0$ for every $v$, so every iterate
after the first step is positive.  Hence
\[
 \mathcal{Z}\bigl(\operatorname{It}_0(n;\,u.\,h)\bigr)
 =\mathcal{Z}(n)\cup\mathcal{Z}\bigl(h\subst{0}{u}\bigr).
\]
The zero set of $h\subst{0}{u}$ is everything if the zero set of $h$ is
everything or contains the hyperplane $u=0$, and otherwise it is the union
of the same hyperplanes of the parameters.  A union of two sets of the
stated form is of that form.  Conversely $0$ and $\Suc 0$ realize the whole
set and the empty set, and $\operatorname{It}_0(c;\,u.\,a)$ with $u$ not in
$a$ has the zero set $\mathcal{Z}(c)\cup\mathcal{Z}(a)$, so nested terms of
this kind realize every finite union of hyperplanes.
\end{proof}

\begin{corollary}[What a zero seed loses]
\label{kern:cor:zeroseed}
\begin{enumerate}[label=\textup{(\roman*)}]
\item No term of the scalar language denotes the predecessor, whose zero set
 is $\{0,1\}$, or addition, whose zero set is $\{(0,0)\}$.
\item No term of either language denotes the function that is $1$ at $0$
 and $0$ elsewhere.  Hence no term denotes the selector $\sel(c,a,b)$, and
 no translation of terms recovers the iterator of \cref{base:sec:terms}
 from either language on the same domain with the standard $0$ and $\Suc$.
\item In the counter-aware language
 $\operatorname{Rec}_0(n;\,k,u.\,k)$ denotes the predecessor, so
 \cref{kern:thm:zeroprofile} does not extend to it.
\end{enumerate}
\end{corollary}

\begin{proof}
(i) \Cref{kern:thm:zeroprofile}.  (ii) The function decreases from $0$ to
$1$, against \cref{kern:thm:mono}; it is $\sel(c,1,0)$.  (iii) The
computation equations give $0$ at $0$ and $n$ at $\Suc n$.
\end{proof}

These are statements about definability by terms on the same domain in the
standard model.  They do not exclude graphs of functions defined by
formulas, definable subdomains, domains of tuples or quotient equalities;
an interpretation of that kind would need its own proofs of totality and
functionality, of the transfer of induction and of both composites.  Whether
addition is denoted by a term of the counter-aware language is not decided
in the package.  In Lean, \cref{kern:thm:mono} and part (ii) of the
corollary are \Lean{ZTm.mono} and \Lean{ZTm.no\_zero\_test} in the module
\Lean{ZeroSeedObstruction}, for a zero-seeded recursor whose body sees the
counter, with \Lean{Tm.val\_zf} for the selector of the kernel's language.

% ---------------------------------------------------------------------
\section{Binder depth and finite axiomatizations}
\label{kern:sec:depth}

This section draws on the binder-depth section and the appendix on finite
axiomatizations of \pkg{reductions}.  Neither result changes a rule.  The
first determines what native formulas of bounded binder depth express under
the standard reading; the second recalls why a presentation of $\iPRA$ with
finitely many principles has to use schemata, as the eight rules of $\Nat$
do.

\subsection{The relations defined at bounded depth}
\label{kern:sec:depthclass}

For a native formula $A$ let $\delta(A)$ be its \emph{binder depth}:
$\delta(s=t)\defeq0$ and $\delta(\Q{y}B)\defeq\delta(\D{x}{y}B)\defeq
1+\delta(B)$.  Each binder has one immediate subformula, so $\delta(A)$ is
also the number of occurrences of $\Q{}$ and $\D{}{}$ in $A$; it is the
binder rank $r(A)$ of \cref{base:rem:single-binder-size} and the depth used
in \cref{seed:thm:prefix}.  Substitution of terms preserves it.

Let $\mathbb N$ be the standard model, in which the iterator denotes
iteration.  A native formula $A$ whose free variables are among
$p_1,\dots,p_k$ \emph{defines} the relation
$\{\vec a\in\mathbb N^k:\mathbb N\models A^\circ[\vec a]\}$, where $A^\circ$
is the ordinary reading of \cref{base:def:reading}; this is the
\emph{standard reading} of $A$, and it is classical.  Let $\mathcal R_d(k)$
be the class of relations on $\mathbb N^k$ defined by native formulas of
depth at most $d$.  Numerical parameters may be among the $p_i$ or be fixed
as numerals; there are no set or function parameters.  Let $\mathrm{PR}(k)$
be the class of primitive recursive relations on $\mathbb N^k$, and for
$m\ge1$ let $\Sigma^0_m(k)$ and $\Pi^0_m(k)$ be the classes of relations
\[
 \exists x_1\,\forall x_2\,\exists x_3\cdots x_m\;
 R(x_1,\dots,x_m,\vec p),\qquad
 \forall x_1\,\exists x_2\,\forall x_3\cdots x_m\;
 R(x_1,\dots,x_m,\vec p)
\]
with $m$ alternating single quantifiers and $R$ primitive recursive.  The
following facts are the standard theory of the arithmetical hierarchy
(S.~C.~Kleene, Recursive predicates and quantifiers, Trans.\ Amer.\ Math.\
Soc.\ 53 (1943), 41--73); they rest on primitive recursive pairing, which
contracts adjacent quantifiers of the same kind.
\begin{enumerate}[label=\textup{(H\arabic*)}]
\item $\mathrm{PR}(k)\subseteq\Sigma^0_1(k)\cap\Pi^0_1(k)$ and
 $\Sigma^0_m(k)\cup\Pi^0_m(k)\subseteq\Sigma^0_{m+1}(k)\cap\Pi^0_{m+1}(k)$.
\item The complement of a $\Sigma^0_m$ relation is $\Pi^0_m$, and
 conversely.
\item $\Sigma^0_m$ is closed under existential quantification of an
 argument, $\Pi^0_m$ under universal quantification; both are closed under
 finite conjunction and disjunction and under substitution of primitive
 recursive functions, in particular of constants, for arguments.
\item For every $m\ge1$ there is a relation $U_m\in\Sigma^0_m(2)$ such that
 every relation in $\Sigma^0_m(1)$ is $\{n:U_m(e,n)\}$ for some $e$.
\end{enumerate}

\begin{lemma}[Strictness of the hierarchy]
\label{kern:lem:hier}
Let $\langle e,n\rangle$ be a primitive recursive pairing with primitive
recursive projections $\pi_0\langle e,n\rangle=e$ and
$\pi_1\langle e,n\rangle=n$, and put $H_m(q)\defeq U_m(\pi_0q,\pi_1q)$ for
$m\ge1$.  Then $H_m\in\Sigma^0_m(1)\setminus\Pi^0_m(1)$.
\end{lemma}

\begin{proof}
$H_m\in\Sigma^0_m(1)$ by (H3).  Suppose $H_m\in\Pi^0_m(1)$.  Then
$U_m(e,n)\leftrightarrow H_m(\langle e,n\rangle)$ is $\Pi^0_m$ by (H3), so
the relation $\neg U_m(n,n)$ is $\Sigma^0_m$ by (H2) and (H3).  By (H4) it
is $\{n:U_m(e,n)\}$ for some $e$, and at $n=e$ this gives
$U_m(e,e)\leftrightarrow\neg U_m(e,e)$.
\end{proof}

This is the diagonal argument for the hierarchy theorem.  It uses the
universal relations of (H4) and their primitive recursive matrices; it does
not use an evaluator for all primitive recursive functions.

\begin{theorem}[The depth classes]
\label{kern:thm:depth}
For every arity $k$,
\[
 \mathcal R_0(k)=\mathrm{PR}(k),\qquad
 \mathcal R_d(k)=\Sigma^0_{2d}(k)\quad(d\ge1).
\]
Every relation in $\Sigma^0_{2d}(k)$ is defined by a formula of depth $d$
without $\Q{}$.
\end{theorem}

The equalities are equalities of classes of relations on the standard
natural numbers.  They do not say that every native formula of depth $d$ is
intuitionistically equivalent, in $\Teq$ or in $\iPRA$, to the reading of a
prenex formula; the reading of a nested native formula contains
implications with quantified antecedents.

\begin{proof}
\emph{Depth zero.}  By induction on terms, every term denotes a primitive
recursive function of its free variables: $\It{n}{a}{u}{h}$ denotes the
function $F$ with $F(0,a,\vec p)=a$ and
$F(m+1,a,\vec p)=h(F(m,a,\vec p),\vec p)$, where $h(\cdot,\vec p)$ is the
function denoted by the body.  Equality of two primitive recursive functions
is a primitive recursive relation.  Conversely let $R\in\mathrm{PR}(k)$ and
let $\chi$ be a function symbol of $\iPRA$ (\cref{base:def:ipra}) whose
function is $0$ exactly on $R$.  By \cref{base:thm:primrec}(ii), $\Teq$
proves the translations $\sigma$ of the defining equations of every
function symbol, and $\Teq$ is valid in $\mathbb N$; by induction on the
generation of $\chi$, the term $\sigma(\chi(\vec p))$ denotes the function
of $\chi$.  The equation $\sigma(\chi(\vec p))=0$ defines $R$.

\emph{Lower inclusion at positive depth.}  For a variable $z$ let $G_zA$ be
$A$ with every equation $s=t$ replaced by $\sel(z,0,s)=\sel(z,0,t)$, and let
$P$ be the predecessor term of \cref{base:thm:pred}; $\delta(G_zA)=\delta(A)$.
In $\mathbb N$, $\sel(0,0,s)$ and $\sel(0,0,t)$ both denote $0$, and
$\sel(w+1,0,s)$ denotes the value of $s$.  Hence $G_0A$ is true under every
valuation, by induction on $A$, since a formula all of whose equations are
true under every valuation is true; and $G_{w+1}A$ has the same reading as
$A$.  Since $\Teq$ is valid in $\mathbb N$, $P(y+1)=y$.  For a formula
$A(x,y,\vec q)$ and a fresh $z$ put
$\mathcal B_{x,y}A\defeq\D{x}{z}\,G_z\bigl(A(x,P(z),\vec q)\bigr)$.  Its
standard reading is
\[
 \exists x\,\bigl(G_0A(x,P(0),\vec q)\to\forall z\,G_zA(x,P(z),\vec q)\bigr),
\]
whose antecedent is true; the instance $z=0$ of the consequent is true, and
the instance $z=y+1$ has the reading of $A(x,y,\vec q)$.  So
$\mathcal B_{x,y}A$ defines $\exists x\,\forall y\,A(x,y,\vec q)$, and
$\delta(\mathcal B_{x,y}A)=1+\delta(A)$.  Now let
$X(\vec p)\leftrightarrow\exists x_1\forall y_1\cdots\exists x_d\forall y_d\,
R(\vec x,\vec y,\vec p)$ with $R$ primitive recursive, and let
$r(\vec x,\vec y,\vec p)=0$ define $R$, by the case of depth zero.  Apply
$\mathcal B$ to this equation $d$ times, from the innermost pair
$x_d,y_d$ outward.  The result defines $X$, has depth $d$ and contains no
$\Q{}$.  This is the construction of \cref{seed:thm:prefix}, which proves
the equivalence in $\iPRA$; only its standard reading is used here.

\emph{Upper inclusion.}  The argument is classical and concerns
$\mathbb N$.  It is proved by induction on $d\ge1$ that every native formula
$A$ with $\delta(A)\le d$ defines a $\Sigma^0_{2d}$ relation, and a
$\Pi^0_{2d-1}$ relation if its root is $\Q{}$.  An equation defines a
primitive recursive relation, which is $\Sigma^0_{2d}$ by (H1).  Let $A$ be
$\D{x}{y}B$ or $\Q{y}B$ with $\delta(B)\le d-1$.  If $d=1$, $B$ is an
equation and defines a primitive recursive relation; if $d\ge2$, $B$
defines a $\Sigma^0_{2d-2}$ relation by the induction hypothesis.  In both
cases, by (H1)--(H3), $\neg B(x,0)$ is $\Pi^0_{2d-1}$ and $\forall y\,B(x,y)$
is $\Pi^0_{2d-1}$, and so is their disjunction.  Classically
\[
 (\Q{y}B)^\circ\leftrightarrow\neg B^\circ(0)\vee\forall y\,B^\circ(y),
 \qquad
 (\D{x}{y}B)^\circ\leftrightarrow
 \exists x\,\bigl(\neg B^\circ(x,0)\vee\forall y\,B^\circ(x,y)\bigr),
\]
so $\Q{y}B$ defines a $\Pi^0_{2d-1}$ relation, which is $\Sigma^0_{2d}$ by
(H1), and $\D{x}{y}B$ defines a $\Sigma^0_{2d}$ relation by (H3).
\end{proof}

\begin{corollary}[Every level is needed]
\label{kern:cor:depthstrict}
\begin{enumerate}[label=\textup{(\roman*)}]
\item For $m\ge1$, a relation in $\Sigma^0_m\setminus\Pi^0_m$ is not
 defined by a native formula of depth less than $\lceil m/2\rceil$, and a
 relation in $\Pi^0_m\setminus\Sigma^0_m$ is not defined by a native formula
 of depth less than $\lfloor m/2\rfloor+1$.
\item For every $d\ge1$ the unary relation $H_{2d}$ of \cref{kern:lem:hier}
 is defined by a native formula of depth $d$ and by none of smaller depth.
\item $\mathcal R_0(1)\subsetneq\mathcal R_1(1)\subsetneq\mathcal R_2(1)
 \subsetneq\cdots$.
\end{enumerate}
\end{corollary}
% ed.: part (i) is the general form of the argument that the package gives
% for H_{2d}; the package states parts (ii) and (iii).

\begin{proof}
(i) Let $X\in\Sigma^0_m\setminus\Pi^0_m$ be defined at depth $e$.  If
$e=0$, $X$ is primitive recursive and hence $\Pi^0_m$ by (H1).  If $e\ge1$
and $2e\le m-1$, then $X\in\Sigma^0_{2e}\subseteq\Sigma^0_{m-1}
\subseteq\Pi^0_m$ by \cref{kern:thm:depth} and (H1).  Both are excluded, so
$2e\ge m$.  Let $X\in\Pi^0_m\setminus\Sigma^0_m$ be defined at depth $e$.  If
$e=0$, $X\in\Sigma^0_m$; if $e\ge1$ and $2e\le m$, then
$X\in\Sigma^0_{2e}\subseteq\Sigma^0_m$.  So $2e\ge m+1$.
(ii) $H_{2d}\in\Sigma^0_{2d}=\mathcal R_d(1)$ by \cref{kern:thm:depth}, and
(i) with \cref{kern:lem:hier} excludes smaller depth.
(iii) The classes increase with $d$, and $H_{2d+2}$ belongs to
$\mathcal R_{d+1}(1)$ and not to $\mathcal R_d(1)$, by (ii).
\end{proof}

\begin{corollary}[No translation of bounded depth]
\label{kern:cor:nofixed}
For every $d$ there is an ordinary formula $\varphi(q)$ with one free
variable such that no native formula $A$ of depth at most $d$, with free
variables among $q$, satisfies $\mathbb N\models\forall q\,
(\varphi\leftrightarrow A^\circ)$.  Hence no map, computable or not, from
ordinary formulas to native formulas of depth at most $d$ has
$\Teq\proves\varphi\leftrightarrow(\varphi')^\circ$ for every $\varphi$ with
one free variable.  In particular the translations $\varphi^{*}$ of
\cref{base:thm:bridge} and $\varphi^{\dagger}$ of
\cref{base:cor:single-binder-bridge} have unbounded depth on such formulas.
\end{corollary}

\begin{proof}
$H_{2d+2}$ is $\Sigma^0_{2d+2}$; replacing its primitive recursive matrix by
a defining equation (the case of depth zero above) gives a prenex ordinary
formula $\varphi(q)$ that defines it.  By
\cref{kern:cor:depthstrict}(ii) no native formula of depth at most $d$
defines it.  $\Teq$ is valid in $\mathbb N$, so a provable equivalence would
be true there; and by \cref{base:thm:bridge}(i),
$\Teq\proves\varphi\leftrightarrow(\varphi^{*})^\circ$, and likewise for
$\varphi^{\dagger}$ by \cref{base:cor:single-binder-bridge}(i).
\end{proof}

\begin{remark}[Depth and the size of derived quantifiers]
\label{kern:rem:depth-rank}
\Cref{base:rem:single-binder-size} computes the binder rank of the derived
quantifiers: $r(\forall_*x\,A)=2r(A)+2$, while
$r(\mathcal U_xA)=r(A)+1$, and $r(\exists_*x\,A)=r(A)+1$ because the gate
preserves rank.  For a prenex ordinary formula $\varphi$ with $m$ single
quantifiers over one equation, $\varphi^{\dagger}$ therefore has depth $m$,
against $2^{m+1}-2$ for $\varphi^{*}$ when all quantifiers are universal.
\Cref{kern:cor:depthstrict}(i) bounds every native formula with the same
standard reading from below: if $\varphi$ defines a relation in
$\Sigma^0_m\setminus\Pi^0_m$, its depth is at least $\lceil m/2\rceil$, and
\cref{kern:thm:depth} attains $m/2$ for even $m$ by pairing each
existential quantifier with the following universal one in one binder.  On
such formulas the single-binder translation is thus within a factor two of
the least depth.  The depth bound concerns binders only; it says nothing
about the size of terms, about the selectors that the gates add, about
formula sharing or about the length of proofs, and the node counts of
\cref{base:q:single-binder-size} are a different measure.
\end{remark}
% ed.: the comparison with the single-binder translation is editorial; the
% package states the depth classification only.

\begin{remark}[Sources and scope]
\label{kern:rem:depth-sources}
\pkg{reductions} states \cref{kern:thm:depth} for formulas built from
equations and $\D{}{}$, and adds that $\Q{}$, counted as one level, does not
change the classes, because a closure of outer depth $d$ defines a
$\Pi^0_{2d-1}$ relation; the proof above treats both binders at once.  The
lower inclusion is \cref{seed:thm:prefix}; the upper inclusion, the
strictness and \cref{kern:cor:nofixed} are the contribution of the
package.  The classification is a statement about the standard reading and
the numerical variables of a formula.  It is not an intuitionistic
classification by provable prenex equivalence, and it does not exclude a
different primitive connective, a different interpreted domain, or a
shorter proof checker.
\end{remark}

\subsection{Finite ordinary axiomatizations}
\label{kern:sec:finax}

Here a \emph{finite ordinary axiomatization} is a theory in first-order
logic with finitely many nonlogical axioms in a finite signature; finitely
many rule schemata with binding syntax, such as the eight rules of $\Nat$,
are not of this kind.

\begin{proposition}[No finite ordinary axiomatization]
\label{kern:prop:finax}
No finitely axiomatized theory $F$ in intuitionistic first-order logic over
a finite signature is mutually interpretable with $\iPRA$ by ordinary
finite-dimensional relative interpretations.  In particular no such theory
is bi-interpretable with $\iPRA$.
\end{proposition}

\begin{proof}
The proof applies two published theorems.  Suppose $F$ is such a theory.
Replace intuitionistic logic by classical logic in $F$ and in $\iPRA$,
keeping the nonlogical axioms; the second theory becomes classical $\mathrm{PRA}$.
A relative interpretation commutes with the connectives and relativizes the
quantifiers, so it maps every instance of excluded middle to a relativized
instance of excluded middle, which is a classical theorem; the conditions on
the domain and on the interpreted functions are proved in the intuitionistic
theories already.  Hence the two given interpretations remain
interpretations, and the finite classical theory $F^{\mathrm{cl}}$ is
mutually interpretable with $\mathrm{PRA}$.  It is consistent, because it
is interpretable in $\mathrm{PRA}$, which has the standard model.

The function symbols of $\mathrm{PRA}$ form an infinite signature.
M.~B\'ilkov\'a, D.~de~Jongh and J.~J.~Joosten (Interpretability in PRA,
preprint, 2008, Section~2.1), following Beklemishev's analysis of reflection,
state that $\mathrm{PRA}$ is mutually interpretable with the theory
$(\mathrm{EA})^{\omega}_{2}$ in the finite signature of arithmetic, obtained
from elementary arithmetic by $\omega$ iterations of uniform
$\Pi^0_2$ reflection.  A.~Visser (The small-is-very-small principle, Math.\
Logic Quart.\ 65 (2019), 453--478, Section~7.3) identifies $\mathrm{PRA}$ as
strongly locally reflexive, and his Theorem~7.5 states that a strongly
locally reflexive theory is not locally mutually interpretable with a
finitely axiomatized theory.  Global mutual interpretability implies local
mutual interpretability, which gives the contradiction.
\end{proof}

The reduction behind the cited theorem is the following.  An
interpretation of the finitely many axioms of $F^{\mathrm{cl}}$ is verified
in one finite fragment $U_0$ of $(\mathrm{EA})^{\omega}_{2}$, so composition
makes $U_0$, enlarged by the finitely many facts that witness
sequentiality, interpret the whole theory; strong local reflexivity then
interprets in that fragment arithmetic with its own consistency statement,
which the generalized second incompleteness theorem excludes.  The
interpretations in \cref{kern:prop:finax} may use definable domains, finite
tuples and definable equality relations, so the obstruction is not limited
to the whole domain or to actual equality.  It does not exclude a finite
grammar with schematic formation of recursors, finitely many inference
schemata, or a compact checker for an infinite family of axioms.
\pkg{reductions} notes that $\mathrm{PRA}$ is reflexive and is not
essentially reflexive, so that the argument has to use local reflexivity.

% ---------------------------------------------------------------------
\section{Checker sizes}
\label{kern:sec:sizes}

This section collects the measurements reported by \pkg{direct\_closure},
\pkg{direct\_closure\_qprogress}, \pkg{witness\_rule},
\pkg{source\_dispatch}, \pkg{reversal\_core}, \pkg{one\_context},
\pkg{one\_sided\_transport}, \pkg{template\_free},
\pkg{backward\_replacement},
\pkg{scalar\_uniqueness}, \pkg{orbit\_closure}, \pkg{crossed\_uniqueness},
\pkg{uniform\_progress} and \pkg{target\_closure}.  All numbers are for complete standalone Python programs under the
protocol of \cref{base:sec:measure}, as the packages state them.
\Cref{kern:tab:sizes} gives, for each program, its significant tokens, AST
nodes and normalized bytes, and the same three numbers for the control it
is matched against.  Programs in different blocks have different input
contracts and are not comparable across blocks.  The compact programs of
\pkg{kernel\_refinement} and of \pkg{scope\_and\_closure} implement calculi
of other chapters and are reported there; neither package measures a checker
for the calculi $\mathsf{L}$ and $\mathsf{R}_{F}$ of this chapter.
The readable complete checker of \pkg{abstraction\_and\_refinements} for
$\mathsf{A}_7$ has 1672 tokens, 1818 AST nodes and 7132 normalized bytes,
against 1590 / 1741 / 6830 for its matched control, the same program for
$\mathsf{B}_7$ with \rulename{QE} and a supplied matrix, so the removal of
the matrix from the certificate does not make the checker smaller; the other
programs of that package check calculi of other chapters and are reported
there.

{\small
\begin{longtable}{@{}p{0.19\linewidth}p{0.31\linewidth}rrrp{0.2\linewidth}@{}}
\caption{Checker sizes reported by the packages of the chapter: tokens, AST
nodes, normalized bytes; the control is given as tokens / AST nodes /
bytes.}\label{kern:tab:sizes}\\
\toprule
Package & Program & Tok. & AST & Bytes & Matched control\\
\midrule
\endfirsthead
\toprule
Package & Program & Tok. & AST & Bytes & Matched control\\
\midrule
\endhead
\bottomrule
\endfoot
\pkg{direct\_closure} & strict tagged checker with \rulename{QI} &
 1169 & 1337 & 5116 & with \rulename{DP}: 1187 / 1358 / 5199\\
 & compact checker with counter vectors (experiment) &
 271 & 314 & 533 & byte profile of \pkg{neutral\_scope}: 266 / 310 / 529\\
\midrule
\pkg{direct\_closure\_\allowbreak qprogress} & compact checker,
 \rulename{QI} at literal equations only &
 285 & 331 & 568 & \pkg{neutral\_scope}: 266 / 310 / 529\\
 & compact checker for $\Nat_{\det}$ &
 278 & 322 & 552 & \pkg{neutral\_scope}: 266 / 310 / 529\\
\midrule
\pkg{witness\_rule} & compact checker, two introductions, Boolean dispatch &
 264 & 310 & 550 & byte profile of \pkg{certified\_zero}: 264 / 310 / 558\\
 & compact checker for $\mathsf{W}$ &
 266 & 312 & 555 & the same; joint profile of \pkg{certified\_zero}:
 262 / 296 / 598\\
 & the same with an explicit type test &
 266 & 312 & 559 & the same\\
 & the same with an explicit list test &
 266 & 310 & 563 & the same\\
\midrule
\pkg{source\_dispatch} & byte profile & 266 & 305 & 526 &
 266 / 310 / 529\\
 & AST profile & 263 & 283 & 607 & 263 / 289 / 613\\
 & token profile & 259 & 289 & 572 & 261 / 292 / 576\\
\midrule
\pkg{reversal\_core} & typed checker for $\mathsf{X}$ &
 456 & 520 & 1055 & old typed control: 532 / 618 / 1365; same refactoring
 without reversal: 450 / 514 / 1049\\
 & $\mathsf{X}$ with a validator for goals of equations and $\D{}{}$ &
 379 & 432 & 830 & old: 450 / 521 / 1109; without reversal:
 373 / 426 / 824\\
 & $\mathsf{X}$ with a validator for all goals &
 405 & 462 & 879 & old: 481 / 560 / 1189; without reversal:
 399 / 456 / 873\\
 & compact byte profile with a scalar neutral premise &
 265 & 304 & 525 & 266 / 305 / 526\\
 & compact AST profile with a scalar neutral premise &
 262 & 282 & 606 & 263 / 283 / 607\\
 & compact checker for $\mathsf{X}$ &
 271 & 310 & 531 & none stated\\
\midrule
\pkg{one\_context} & tagged checker for $\Nat$ &
 1089 & 1150 & 3358 & for $\Nat_{2}$: 1101 / 1160 / 3382\\
 & strict validator, equations and $\D{}{}$ &
 359 & 414 & 803 & 373 / 426 / 824\\
 & strict validator, all formulas &
 391 & 444 & 863 & 399 / 456 / 873\\
 & full typing of the certificate &
 442 & 502 & 1039 & 450 / 514 / 1049\\
 & strict validator with reversal, equations and $\D{}{}$ &
 365 & 420 & 809 & 379 / 432 / 830\\
 & strict validator with reversal, all formulas &
 397 & 450 & 869 & 405 / 462 / 879\\
 & full typing with reversal &
 448 & 508 & 1045 & 456 / 520 / 1055\\
\midrule
\pkg{one\_sided\_\allowbreak transport} & strict checker for $\mathsf{M}$ &
 1171 & 1339 & 3246 & $\rulename{Tr}_{2}$ at the same sources:
 1196 / 1364 / 3302\\
\midrule
\pkg{template\_free} & strict tagged checker for $\mathsf{G}_7$ &
 1027 & 1072 & 3147 & $\mathsf{W}$ with three traversals:
 1029 / 1089 / 3181; 1059 / 1124 / 3215; 1044 / 1096 / 3206\\
 & compact forward checker with \rulename{Rep} &
 397 & 464 & 865 & crossed rule, another packet layout: 384 / 448 / 830\\
\midrule
\pkg{backward\_\allowbreak replacement} & readable strict checker for
 $\mathsf{B}$ &
 1647 & 1807 & 7121 & the same with \rulename{Tr} and a supplied template:
 1581 / 1738 / 6971\\
 & compact checker of the guarded calculus with regrouped packets, not of
 $\mathsf{B}$ &
 250 & 281 & 604 & the same with flat packets: 251 / 283 / 605\\
 & strict companion of that checker &
 402 & 451 & 971 & the same with flat packets: 403 / 453 / 972\\
\midrule
\pkg{scalar\_\allowbreak uniqueness} & tagged checker for $\mathsf{U}$ &
 1124 & 1175 & 3393 & for $\Nat$: 1084 / 1141 / 3322\\
 & validated codec with scope reuse, calculus unchanged &
 381 & 439 & 866 & 383 / 442 / 870\\
\midrule
\pkg{orbit\_closure} & tagged checker for $\mathsf{U}_{\mathrm{o}}$,
 templates supplied &
 1132 & 1178 & 3400 & for $\mathsf{U}$: 1124 / 1175 / 3393\\
 & the same, templates read off the closure &
 1145 & 1194 & 3455 & the same\\
\midrule
\pkg{crossed\_\allowbreak uniqueness} & strict checker for
 $\mathsf{U}_{\mathrm{c}}$, normalizer included &
 1909 & 2087 & 5618 & none stated\\
 & compact checker of the guarded calculus, not of
 $\mathsf{U}_{\mathrm{c}}$ &
 269 & 301 & 502 & byte profile of
 \pkg{deterministic\_\allowbreak introduction}: 267 / 303 / 506\\
\midrule
\pkg{uniform\_\allowbreak progress} & typed checker for $\mathsf{R}$ &
 1206 & 1366 & 4553 & for $\Nat$: 1235 / 1407 / 4692; for $\mathsf{W}$:
 1173 / 1334 / 4456\\
 & compact token profile, guarded-source calculus &
 256 & 283 & 593 & 257 / 285 / 568\\
 & compact AST profile, the same calculus &
 260 & 278 & 602 & 261 / 279 / 603\\
 & compact byte profile, the same calculus &
 263 & 300 & 521 & 264 / 301 / 522\\
\midrule
\pkg{target\_closure} & typed checker for $\mathsf{T}$ &
 1195 & 1358 & 4532 & for $\mathsf{R}$: 1206 / 1366 / 4553\\
 & typed checker of the conditional calculus with a lexical gate &
 399 & 453 & 895 & recursive typing: 442 / 502 / 1039\\
 & the same with a descriptive error &
 404 & 455 & 963 & the same\\
\end{longtable}
}

The following paragraphs summarize what the programs do beyond the rules.

\emph{\pkg{direct\_closure}.}  The two strict checkers are built from one
body and two fragments for the progress rule; they differ in that branch
only.  The saving is the witness, its scope check and one substitution; the
normalized source compressed by gzip is 19 bytes larger.  The experiment
with counter vectors passes the counter of a computation as an empty or
nonempty list, so that $\It{1}{a}{u}{h}$ is certified without the repair of
\cref{kern:lem:comprepair}; it costs four bytes.  A variant of 528 bytes
writes the empty list where the protocol uses the scalar zero, a canonical
certificate fails on it, and it is not counted as a complete checker.

\emph{\pkg{direct\_closure\_qprogress}.}  The compact programs keep the
arity encoding of the control.  A failed test on the matrix must reject and
not return a false value, which in this encoding would enter the context as
the scalar zero; the first program uses a sentinel that is not JSON for the
failure, and the second gives a compound matrix the mode of cases.

\emph{\pkg{witness\_rule}.}  Every selected first premise is truthy and
every generated computation equation is a nonempty list, so one Boolean
expression selects the branch; null markers become empty lists.  This saves
eight bytes with the calculus unchanged, and the test of the witness adds
five.  The package proves that every accepting run is simulated, after a
fixed change of markers, by an accepting run of the control.

\emph{\pkg{source\_dispatch}.}  In canonical certificates the source is
falsy exactly when it is the literal $0$, so the comparison of the first
premise against two patterns becomes one comparison selected by the source.
The byte and AST profiles also move the pure computation of the replaced
goal before the check of the first premise, and the byte profile uses the
checker function itself as the invalid image of a substitution.

\emph{\pkg{reversal\_core}.}  The typed checker shares one validator between
terms and formulas and writes the neutral premise as a scalar; these two
changes give the reduction of 310 bytes against the old control, and
reversal itself adds six bytes.  The compact programs of 525 bytes and 282
AST nodes check certificates without reversal, for the equivalent calculus.

\emph{\pkg{one\_context}.}  In the tagged checker only the branch of
transport changes.  Independently, the strict validator of the goal omits
the upper bound on variable indices: for a tree of the right shape that
bound holds exactly when the first traversal of the unchanged kernel
succeeds, which happens before any instruction is read, so the sets of
accepted certificates are equal.  The sizes after gzip increase in all six
pairs.

\emph{\pkg{one\_sided\_transport}.}  The fixed side $r$ of
$\rulename{Tr}^{-}$ is read off the conclusion, so the instruction carries
three terms where the control carries four.

\emph{\pkg{template\_free}.}  One traversal serves lifting, substitution and
replacement; an instruction for \rulename{Rep} consists of its two premise
derivations.  Against the control with the same traversal the saving is 59
bytes.  The compact forward program synthesizes its conclusion and is 35
bytes larger than its control, which has another packet layout.

\emph{\pkg{backward\_replacement}.}  The readable checker validates the
syntax and scope of every term, formula and proof node, rejects extra fields
and Boolean values in place of integers, and accepts closed theorems at its
driver; a separate entry point checks open sequents.  Its control differs
only in the branch of the equality rule, which takes a supplied template,
and in the absence of the replacement traversal; that traversal costs 66
tokens, 69 AST nodes and 150 bytes, so removing the template from the
certificate does not make the checker smaller.  The two compact programs
check the guarded calculus of another chapter, unchanged, and regroup the
last three fields of a packet, $[m,s,u,t]$ becoming $[m,[s,u,t]]$; the
package proves by a simulation of runs that the two formats accept
certificates for the same canonical public theorems.  The package compares
with the least counts it collects from the corpus, 251 tokens, 275 AST nodes
and 502 bytes, attained by three different programs: the 250 tokens lower
the first by one, the other two are not improved, and the 250-token program
does not implement \rulename{BRep}.

\emph{\pkg{scalar\_uniqueness}.}  The instruction for \rulename{UR} carries
the two templates, the target and the premise, checked with no added
assumption; the direct replacement of \rulename{QI} by \rulename{UR} costs
71 bytes.  The second pair keeps the calculus of $\Nat$: a fully typed codec
whose goal validator relaxes the test \texttt{2 < a < d} on a variable index
to \texttt{a > 2} and leaves the upper bound to the first scope traversal of
the kernel, the factorization described for \pkg{one\_context}; the change
is four bytes.

\emph{\pkg{orbit\_closure}.}  The instruction for \rulename{Orb} carries a
proposed closure and one premise.  The selected checker reads $f$ and $h$
off the closure, compares the seed with the instance of $f$ at zero, and
rejects a body that mentions the counter (\cref{kern:rem:orbcounter}); it is
62 bytes larger than its control, the checker for $\mathsf{U}$ of
\pkg{scalar\_uniqueness}.  The translator of the package is not trusted; its
output is checked by the target checker.  The package also audits every
Python source of the corpus, 5,825 paths with 1,380 distinct contents, by
parsing and measuring without executing them, and finds the least values of
the three coordinates, 256 tokens, 276 AST nodes and 506 bytes, in three
different programs of other packages; it records which smaller files are
not complete standalone checkers.

\emph{\pkg{crossed\_uniqueness}.}  The checker for
$\mathsf{U}_{\mathrm{c}}$ counts the normalizer, substitution and all scope
checks, introduces the eigenvariable of \rulename{UX} itself, and checks
scopes before normalizing, so that a subterm erased by a computation is
still checked.  Separately, the package composes three known refinements
(delayed elimination, eager selection from a pair, a shorter selector of
the mode) into a compact checker of the guarded calculus of another chapter
with 502 bytes against 506, and a validating companion with 1172 bytes
against 1183; neither implements $\mathsf{U}_{\mathrm{c}}$, and the least
token and AST counts of the corpus are not improved.

\emph{\pkg{uniform\_progress}.}  The typed checker for $\mathsf{R}$ is 139
bytes smaller than the matched checker for $\Nat$ and 97 bytes larger than
the matched checker for $\mathsf{W}$.  Its three compact programs are not
checkers for $\mathsf{R}$; they turn the continuation into a loop over an
immutable context and write the neutral premise as a scalar.

\emph{\pkg{target\_closure}.}  The checker for $\mathsf{T}$ saves 21 bytes
against that for $\mathsf{R}$ by dropping the source and the seed of the
orbit, and remains 76 bytes larger than that for $\mathsf{W}$.  The lexical
gate replaces the recursive test that every value of a certificate is an
integer or a list by a test on the input text: a text that the JSON parser
accepts decodes to such a value exactly when all its characters are among
the brackets, the comma, the minus sign, the digits and the four whitespace
characters of JSON.

% ---------------------------------------------------------------------
\section{Further questions}
\label{kern:sec:questions}

The packages of the chapter state the following questions, limits and
unproved claims.  None of them claims a smallest calculus or a lower bound
on the size of a checker.

\begin{question}[Checkers for the direct closure]
\label{kern:q:directsize}
The compact checkers of \pkg{direct\_closure\_qprogress} with \rulename{QI}
are larger than their control, because the two openings at the root must be
distinguished.  Is there an encoding in which the direct closure rule gives
a smaller complete checker, with malformed, aliased and scope-violating
inputs counted?
\end{question}

\begin{question}[The variant with counter vectors]
\label{kern:q:countervec}
The 528-byte variant in \pkg{direct\_closure} fails on a canonical
certificate because the base counter is written as an empty list.  Is there
a conversion between the two spellings of zero that makes it complete?  No
false theorem is known to be accepted by it.
\end{question}

\begin{question}[Deleting an inference case]
\label{kern:q:deletecase}
The direct closure rule removes a witness and keeps eight rules.  Is there a
basis on the same formulas that deletes a case of the checker and preserves
arbitrary contexts and both formula round trips?
(\pkg{direct\_closure\_qprogress}.)
\end{question}

\begin{question}[Focused derivations]
\label{kern:q:focused}
Can a normalization theorem that permutes introductions and eliminations
justify a simpler protocol for first premises?  Equation induction makes the
theorem nontrivial, and it has to precede any removal of checking.
(\pkg{direct\_closure\_qprogress}.)
\end{question}

\begin{question}[One-template transport in a compact protocol]
\label{kern:q:compactone}
Can a certificate format expose \rulename{Tr} or $\rulename{Tr}^{-}$ and
keep the opening routine shared with \rulename{QE}, so that the complete
checker is smaller than the compact programs under the same contract for
the goal?  (\pkg{one\_context}, \pkg{one\_sided\_transport}.)
\end{question}

\begin{question}[Sharing in certificates]
\label{kern:q:dag}
What is the least overhead in the size of certificates when
$\rulename{Tr}_{2}$ is replaced by \rulename{Tr}?  A certificate with
sharing avoids the copy of the first premise (\cref{kern:rem:dag}); can a
checker for such certificates be smaller?  (\pkg{one\_context},
\pkg{one\_sided\_transport}.)
\end{question}

\begin{question}[Constant templates]
\label{kern:q:constant}
Symmetry is derived from \rulename{Tr} with a constant template.  Is there
an equality rule simpler than $\rulename{Tr}^{-}$ that keeps this step?  A
restriction to templates in which the variable occurs must replace it.
(\pkg{one\_sided\_transport}.)
\end{question}

\begin{question}[Removing both syntactic tests]
\label{kern:q:bothtests}
The test $s\not\equiv 0$ of $\rulename{Tr}^{-}$ and the test of the witness
in $\rulename{DI}_{\mathrm{w}}$ serve the dispatch of compact checkers; the
unrestricted rules are admissible.  Can both tests be removed while a
compact complete checker is kept?  (\pkg{one\_sided\_transport}.)
\end{question}

\begin{question}[The forward computation equation]
\label{kern:q:forward}
If the successor computation is replaced by
$\It{\Suc n}{a}{u}{h}=\It{n}{h(a)}{u}{h}$, is the equation \rulename{IS}
derivable?  The induction that suggests itself needs its hypothesis at the
changed seed $h(a)$, while equation induction keeps parameters fixed.
\pkg{one\_sided\_transport} leaves this open and claims no independence.
\end{question}

\begin{question}[A fixed input contract]
\label{kern:q:contract}
Which contract should later comparisons fix: goals of equations and
$\D{}{}$, goals of all native formulas, or exact typing of the whole
certificate?  The minima reported by the packages are attained by different
programs under different contracts.  (\pkg{one\_context},
\pkg{scalar\_uniqueness}.)
\end{question}

\begin{question}[Replacement without a template in a compact checker]
\label{kern:q:tfreesize}
Does the rule \rulename{Rep} give a smaller complete checker under a fixed
compact contract?  The compact forward program of \pkg{template\_free} is
larger than its control.
\end{question}

\begin{question}[A checker for backward replacement]
\label{kern:q:brepsize}
Can matching under binders, with its scope shifts, share one traversal with
the opening routine of the compact checkers, so that \rulename{BRep} removes
more trusted code than it adds?  The readable checker for $\mathsf{B}$ is
larger than its matched control in all three coordinates.  Does
$\mathsf{B}^{-}_{7}$ admit a smaller complete checker than either of its two
restrictions alone?  An answer has to report one complete checker under a
fixed input contract with all three counts.  (\pkg{backward\_replacement}.)
\end{question}

\begin{question}[Fusing backward replacement]
\label{kern:q:brepfuse}
Can \rulename{BRep} be fused with the computation equations or with the
discharge of a witness without a second traversal of syntax?  Such a rule
has to handle the zero and the successor computation, constant templates and
replacement below binders.  What overhead in the size of certificates is
unavoidable when a chosen replacement is compiled through the marker of
\cref{kern:lem:marker}, whose size is not optimized?
(\pkg{backward\_replacement}.)
\end{question}

\begin{question}[A checker for elimination by abstraction]
\label{kern:q:aqesize}
The readable checker for $\mathsf{A}_7$ is larger than its matched control
for $\mathsf{B}_7$ in all three coordinates: the abstraction traverses
compound formulas and shifts its patterns below binders.  Can the
abstraction of \rulename{AQE} share one traversal with \rulename{BRep} and
with the opening routine, or be combined with another sound protocol of
certificates, so that one complete checker under a fixed input contract is
below the compact records of 244 tokens, 265 AST nodes and 477 bytes
together, with an explicit interpretation?  The package makes no assumption
that such a combination exists.  (\pkg{abstraction\_and\_refinements}.)
\end{question}

\begin{question}[Replacement of one occurrence in a checker]
\label{kern:q:lor}
Does \rulename{LOR} help a readable or a formally verified checker, even
though a traversal by position costs more than it saves in the compact
programs?  (\pkg{kernel\_refinement}.)
\end{question}

\begin{question}[Uniqueness with fewer bodies]
\label{kern:q:urweak}
Can \rulename{UR} be restricted to a family of bodies and still give the
predecessor of \cref{kern:thm:urpred}?  The reflection of equality uses
mostly the case of a constant recurrence (\cref{kern:lem:cu}(ii)), while the
predecessor uses other bodies.  For the retractions the question is
answered: \cref{kern:lem:rigid} needs the body $\Suc u$ only, and by
\cref{kern:prop:rinv} it needs no coverage.
(\pkg{scalar\_uniqueness}, \pkg{orbit\_closure}, \pkg{reductions}.)
\end{question}

\begin{question}[The strength of successor uniqueness]
\label{kern:q:su}
Does $T_{\mathrm{SU}}$ derive induction for all literal equations?  A
positive answer has to supply a predecessor and a term that reflects
equality from the body $\Suc u$ alone; a negative answer has to give a model
of the computation equations for every body and of every instance of
\rulename{SU} over arbitrary contexts in which equation induction fails.
The model that separates induction for zero equations from induction for
all equations is not such a model (\cref{kern:rem:sugap}).  The concrete
missing step is a term $\Delta$ with $\Delta(a,b)=0\leftrightarrow a=b$,
derived from \rulename{SU} and the weak scalar axioms; the right inverse of
\cref{kern:prop:rinv} does not give it.
(\pkg{orbit\_closure}, \pkg{reductions}.)
\end{question}

\begin{question}[Coverage on a changed domain]
\label{kern:q:coverdomain}
The product model of \cref{kern:rem:rinv-product} shows that coverage cannot
be deleted from the scalar presentation on the whole domain with the fixed
language.  Can a presentation without coverage be bi-interpretable with
$\iPRA$ through a definable domain?  Such an interpretation needs definable
comparison isomorphisms in both theories; an encoding that is correct in
the standard model alone does not give bi-interpretability.
(\pkg{reductions}.)
\end{question}

\begin{question}[From depth to size]
\label{kern:q:depthsize}
\Cref{kern:thm:depth,kern:cor:depthstrict} show that some relations need
every binder of a given depth.  Is there a corresponding lower bound on the
size of numerical terms, on the cost of the gates and of pairing, or on the
length of proofs?  Such a bound would measure a different coordinate of
simplicity; compare \cref{base:q:single-binder-size}.  (\pkg{reductions}.)
\end{question}

\begin{question}[Successor and identity bodies]
\label{kern:q:suid}
Does \rulename{SU} together with uniqueness for the body $u$
(\cref{kern:lem:cu}(ii)) give induction for all equations?  Does
$T_{\mathrm{SU}}$ prove that addition is cancellative, which by
\cref{kern:prop:suadd} would give uniqueness for the body $u$?
(\pkg{orbit\_closure}.)
\end{question}

\begin{question}[Compact encodings for the other forms of uniqueness]
\label{kern:q:urcompact}
Can \rulename{Orb} be combined with the compact encodings of goals without
adding more recognition of syntax than it removes?  Can \rulename{UX}
receive a compact encoding that keeps the freshness of its eigenvariable
and the independence of its body from the counter, and can the normalizer
share the traversal by arity while respecting sorted positions, protected
slots and reindexing?  The checkers for $\mathsf{U}_{\mathrm{o}}$ and
$\mathsf{U}_{\mathrm{c}}$ are larger than their comparators.
(\pkg{orbit\_closure}, \pkg{crossed\_uniqueness}.)
\end{question}

\begin{question}[Crossed uniqueness with clauses]
\label{kern:q:uxclauses}
Can a calculus with clause contexts fuse \rulename{UX} with the discharge
of a witness without complicating the derivation of equality in
\cref{kern:lem:uxeq}?  (\pkg{crossed\_uniqueness}.)
\end{question}

\begin{question}[Uses of the accumulator]
\label{kern:q:accumulator}
The term $\varphi$ of \cref{kern:lem:acc} is a composition of subtraction
and addition with a recurrence that depends on the counter.  Does it give a
more direct compiler for primitive recursion, or a shorter predecessor,
without coding tuples?  (\pkg{scalar\_uniqueness}.)
\end{question}

\begin{question}[A checker for uniqueness]
\label{kern:q:ursize}
The checker for $\mathsf{U}$ is 71 bytes larger than its control.  Is there
an architecture that shares the check of a computation equation with the
check of \rulename{UR} and reverses the comparison?
(\pkg{scalar\_uniqueness}.)
\end{question}

\begin{question}[Cost of the root-sensitive closure]
\label{kern:q:mixedcost}
Can $\mathsf{R}$, $\mathsf{R}_{F}$ or $\mathsf{T}$ be checked with less
total work than the test of the witness in $\mathsf{W}$?  The matched
checkers for $\mathsf{R}$ and $\mathsf{T}$ are larger than that for
$\mathsf{W}$, and the comparison has to include the instances at the source
and at the target, the construction of the orbit term and the test of the
root.  (\pkg{uniform\_progress}, \pkg{target\_closure},
\pkg{scope\_and\_closure}.)
\end{question}

\begin{question}[A zero seed through graphs]
\label{kern:q:graph}
\Cref{kern:cor:zeroseed} excludes a translation of terms.  Can an
interpretation by graphs of functions recover an arbitrary seed from a
language with the seed fixed to zero, with proofs of totality and
functionality, the transfer of induction and both composites?  Is addition
denoted by a term of the counter-aware language of
\cref{kern:sec:zerostart}?  (\pkg{scope\_and\_closure}.)
\end{question}

\begin{question}[Lower bounds]
\label{kern:q:lower}
Is there a lower bound on the size of presentations in a fixed class, with
a fixed encoding and cost model?  A census of existing programs is not such
a bound.  (\pkg{one\_context}, \pkg{scope\_and\_closure}.)
\end{question}

The following statements of the packages are recorded without proof or as
results that were not achieved.
\begin{itemize}
\item \pkg{direct\_closure\_qprogress} states that replacing the two
 computation equations by a rule of conversion in a context would still
 need both rewriting patterns and an account of reflexivity, and that a
 conversion engine left out of the count would not be a smaller checker.
 No calculus of this kind is given.
\item \pkg{witness\_rule}: no program of the package is below all three
 earlier coordinate minima; the restriction of the introduction is not
 credited with a saving of bytes.
\item \pkg{one\_context} states that removing a primitive equality test
 from the self-seeded proposals of other packages is a missing theorem and
 does not follow from coverage of a protected orbit.
\item \pkg{scalar\_uniqueness}: the converse of \cref{kern:prop:uniq} for
 the ordinary theory is formalized in Lean (\Lean{tur\_iff\_teq},
 \cref{kern:rem:ur-sources}); the native form \cref{kern:thm:urnative} is a
 written proof only.
\item \pkg{orbit\_closure}: the equivalence of $T_{\mathrm{SU}}$ with $\Teq$
 is not established, and the package does not claim it.  Its statement
 about the model of \pkg{canonical\_recursors} (\cref{kern:rem:sugap}) rests
 on the construction of that package.  The rule \rulename{Orb} removes a
 field of the certificate and makes the matched checker larger, by 7 bytes
 with supplied templates and by 62 bytes with templates read off the
 closure.
\item \pkg{crossed\_uniqueness}: six rule names do not by themselves make a
 smaller presentation than the three-family clause calculi of other
 packages; the checker for $\mathsf{U}_{\mathrm{c}}$ is the largest of
 \cref{kern:tab:sizes}, and its normalizer is counted.  The 502-byte
 program lowers one coordinate for a calculus of another chapter.  The
 package asks whether a relational bi-interpretation on a changed domain
 could reduce a complete checker; the obstructions to definability by
 terms do not exclude it, and both composites would need proofs.
\item \pkg{orbit\_closure}, \pkg{crossed\_uniqueness}: neither is
 formalized in Lean, and both rest on the written proof of
 \cref{kern:thm:ur}.
\item \pkg{template\_free}: the calculus $\mathsf{G}_7$, with \rulename{Rep}
 and the single introduction together, is formalized in Lean
 (\Lean{g7prf\_iff\_nprf}, generically over \Lean{TFKernel}).
\item \pkg{backward\_replacement}: $\mathsf{B}$ and $\mathsf{B}_{7}$ are
 formalized in Lean (\Lean{NativeReplacementCalculi}), and so are
 $\mathsf{B}^{-}$ and $\mathsf{B}^{-}_{7}$ (\Lean{NativeNonzeroSources}); the package
 suggests extending the Lean calculus \Lean{TPrf} of $\mathsf{G}$ and reusing its marker
 lemma \Lean{Tm.rep\_marker}, with the bridge unchanged.  Its checker is not
 formally verified, and the restricted calculi are not implemented.  No
 program of the package is below the three least counts of the corpus
 together, and the readable checker is not claimed to bound from below the
 size of an optimized checker for \rulename{BRep}.  The package records
 three routes that gave no result: removing the seed of the iterator meets
 the obstructions for a single body, reusing the compact opening routine as
 a normalizer needs an account of dummy slots and of shifts of bound
 variables, and replacing a stored scope bound by the length of the context
 has to keep apart the check of a packet at the old scope and the opening of
 a branch at the new one.
\item \pkg{abstraction\_and\_refinements}: \cref{kern:thm:aqe} is written in
 the package and formalized in Lean in this repository
 (\Lean{NativeAbstraction}, \Lean{aprf\_iff\_nprf}), not in Rocq; the
 bridge to $\iPRA$ is the cited \cref{base:thm:bridge}.  The package does
 not claim that seven rules, the signature, the protocol of certificates or
 any measured source representation is minimal, and states that the
 coordinatewise envelope of the corpus is not a lower bound on
 bi-interpretable presentations.  \rulename{AQE} is a reduction of the data
 of an inference and not of the size of the checker.
\item \pkg{uniform\_progress}, \pkg{target\_closure}: neither calculus gives
 a new minimum in any coordinate of the matched comparison; the compact
 minima of other packages are not improved by a checker for $\mathsf{R}$ or
 $\mathsf{T}$.
\item \pkg{target\_closure} reports two routes that gave no result: making
 the guard of an atomic matrix implicit simplifies the datatype of goals
 and adds a cost of elaboration that loses against a validator factored in
 the same way; and a stopped-clock construction in a self-seeded arithmetic
 gives a diagonal equation and does not prove that the successor is
 injective; two paired stopping orbits do (\cref{seed:thm:tf-inj}).  No
 independence result is inferred from either.
\item \pkg{scope\_and\_closure}: the obstruction of
 \cref{kern:sec:zerostart} is about terms in the standard model and is not
 an impossibility theorem for simpler axiomatizations.
\item \pkg{reductions} states that it does not establish a new
 axiomatization below the complete-checker records, deletes no arithmetic
 or logical axiom, claims no absolute minimality and gives no new Lean or
 Rocq verification.  \Cref{kern:prop:rinv} does not bootstrap equality
 reflection or induction for all equations from \rulename{SU}.
 \Cref{kern:thm:depth} is a theorem about the classical standard reading;
 it is not an intuitionistic prenex normal form theorem and not a lower
 bound on the size of terms, of proofs or of a checker.
 \Cref{kern:prop:finax} rests on the cited theorems of Visser and of
 B\'ilkov\'a, de~Jongh and Joosten, which are not proved here, and it does
 not apply to finite rule schemata.
\item The proofs of this chapter are written proofs; the Lean names in the
 remarks mark the statements that are formalized.  The Python checkers are
 not formally verified, and the tests reported by the packages concern
 particular certificates.
\end{itemize}
